% Lutte contre les perturbations du système immunitaire — SVT Tle D — Cours signé OnBuch+
% Programme officiel MINESEC. Compiler avec : tectonic N11-systeme-immunitaire-perturbations.tex
\newcommand{\DOCMATIERE}{SVT}
\newcommand{\DOCNIVEAU}{Tle D}
\newcommand{\DOCMODULE}{Module II : L'Éducation à la Santé — Mécanismes et dysfonctionnements du système immunitaire}
\newcommand{\DOCLECON}{N11}
\newcommand{\DOCTITRE}{Lutte contre les perturbations du système immunitaire}
% OnBuch+ — préambule « cours signés » ENRICHI (compilé avec tectonic / XeLaTeX)
% Système visuel « Pop » d'OnBuch+ : fond crème (jamais blanc), bordures encre
% épaisses, ombres « dures » décalées, titres Archivo Black en capitales.
% Version MATHÉMATIQUES : graphes (pgfplots/tikz), boîtes pédagogiques
% (définition, propriété, méthode, attention, le savais-tu, exercice résolu,
% exercice, TP, bilan), page de garde et signature OnBuch+ ancrée en bas.
%
% Macros attendues dans main.tex AVANT \input{preamble} :
%   \DOCMATIERE  \DOCNIVEAU  \DOCMODULE  \DOCLECON (numéro)  \DOCTITRE
\documentclass[11pt]{article}

\usepackage{fontspec}
\usepackage[french]{babel}
\usepackage[a4paper,margin=2cm,top=3.4cm,bottom=2.8cm,headheight=1.4cm,footskip=1.3cm]{geometry}
\usepackage{amsmath,amssymb,amsfonts}
\usepackage{mathtools} % \xmapsto, \coloneqq... souvent utilisés par les modèles
\usepackage{cancel} % simplifications barrées (bilans de forces)
% \abs{z} (module, valeur absolue) et \norm{u} (norme), utilisés par les modèles
\providecommand{\abs}{}\let\abs\relax\DeclarePairedDelimiter{\abs}{\lvert}{\rvert}
\providecommand{\norm}{}\let\norm\relax\DeclarePairedDelimiter{\norm}{\lVert}{\rVert}
\usepackage[table]{xcolor} % [table] : fournit aussi \rowcolors (alternance de lignes)
\usepackage{pagecolor}
\usepackage{tikz}
\usetikzlibrary{arrows.meta,positioning,calc,shapes.geometric,shapes.misc,shapes.arrows,decorations.pathmorphing,decorations.pathreplacing,decorations.markings,patterns,shadows,mindmap,trees,fit,backgrounds,matrix,intersections,angles,quotes}
\usepackage{pgfplots}
\usepackage[version=4]{mhchem} % équations nucléaires \ce{^{235}_{92}U}
\usepackage[european,siunitx]{circuitikz} % schémas électriques
\pgfplotsset{compat=1.18}
\usepgfplotslibrary{fillbetween,groupplots} % aires entre courbes (calcul intégral)
\usepackage{enumitem}
\usepackage{fancyhdr}
% [most] charge toutes les bibliothèques tcolorbox (listings, minted, raster,
% documentation...) et gonfle inutilement la mémoire TeX sur un document long
% (~300 boîtes) : on ne charge que ce qui est réellement utilisé.
\usepackage[skins,breakable]{tcolorbox}
\usepackage{graphicx}
\usepackage{colortbl}
\usepackage{booktabs}
\usepackage{tabularx}
\usepackage{diagbox} % case d'en-tête diagonale des tableaux à double entrée
% Colonnes extensibles centrée / gauche / droite (C, L, R), souvent utilisées par les modèles
\newcolumntype{C}{>{\centering\arraybackslash}X}
\newcolumntype{L}{>{\raggedright\arraybackslash}X}
\newcolumntype{R}{>{\raggedleft\arraybackslash}X}
\usepackage{array}
\usepackage{multicol}
\usepackage{siunitx}
% parse-numbers=false : accepte aussi \SI{2,5 \times 10^{-3}}{...}
\sisetup{locale=FR,output-decimal-marker={,},group-minimum-digits=4,parse-numbers=false}
\DeclareSIUnit\ohm{\ensuremath{\symup{\Omega}}} % U+2126 absent de la police
\DeclareSIUnit\division{div} % réglages d'oscilloscope (ms/div, V/div)
\DeclareSIUnit\tr{tr} % tours (vitesse de rotation, ex. \SI{1500}{\tr\per\minute})
\DeclareSIUnit\molar{M} % concentration molaire (ex. \SI{3}{\molar}, \SI{1}{\micro\molar})
\DeclareSIUnit\an{an} % année (le modèle confond parfois avec \year, un compteur TeX) — ex. \SI{5}{\kilo\meter\per\an}
\DeclareSIUnit\FCFA{FCFA} % franc CFA (ex. \SI{500}{\FCFA\per\kilo\gram})
\DeclareSIUnit\degreeBrix{\textdegree Brix} % teneur en sucre d'un sirop/jus (réfractométrie)
\DeclareSIUnit\month{mois} % le modèle utilise parfois \month, un compteur TeX, comme unité de durée
\usepackage{qrcode}
% \degree : commande fréquemment utilisée par les modèles pour un angle ou une
% température en dehors de \SI{}{} (non fournie par les paquets chargés).
\providecommand{\degree}{^{\circ}}
% \qed (fin de démonstration), utilisé par les modèles sans amsthm
\providecommand{\qed}{\hfill$\square$}
% \barre : double barre des valeurs interdites dans un tableau de variations (tkz-tab)
\providecommand{\barre}{\ensuremath{\|}}
% environnement proof (amsthm non chargé) utilisé par les modèles
\ifdefined\proof\else\newenvironment{proof}[1][Démonstration]{\par\noindent\textit{#1.}\ }{\qed\par}\fi
\usepackage{lastpage}
\usepackage{hyperref}
\hypersetup{hidelinks}

% ============================================================
% Palette « Pop » OnBuch+ — mêmes hex que la classe P (plus_ui.dart)
% ============================================================
\definecolor{poporange}{HTML}{F59321}
\definecolor{popdark}{HTML}{E07A0C}
\definecolor{popcream}{HTML}{FFF6E8}
\definecolor{popink}{HTML}{1C1714}
\definecolor{poppurple}{HTML}{7A5AE0}
\definecolor{popgreen}{HTML}{2E9E6A}
\definecolor{poppink}{HTML}{E0457B}
\definecolor{popblue}{HTML}{2F7DE1}
\definecolor{popgold}{HTML}{C98A12}
\definecolor{popmuted}{HTML}{8A7F76}
\definecolor{popline}{HTML}{EADFD2}
\definecolor{popgray}{HTML}{8A7F76}
\colorlet{popgrey}{popgray}
% Alias tolérants (le modèle nomme parfois les couleurs différemment)
\colorlet{popwhite}{white}
\colorlet{popblack}{popink}
\colorlet{popred}{poppink}
\colorlet{popyellow}{popgold}
% Teintes claires (fonds de tableaux/encadrés) — le modèle nomme parfois
% les couleurs avec un suffixe L (ex. popblueL) sans les avoir définies.
\colorlet{poporangeL}{poporange!20}
\colorlet{popdarkL}{popdark!20}
\colorlet{poppurpleL}{poppurple!20}
\colorlet{popgreenL}{popgreen!20}
\colorlet{poppinkL}{poppink!20}
\colorlet{popblueL}{popblue!20}
\colorlet{popgoldL}{popgold!20}
\colorlet{popredL}{popred!20}
\colorlet{popgrayL}{popgray!20}
% Teintes claires pour fonds de tableaux / zones
\colorlet{popblueL}{popblue!10!white}
\colorlet{poporangeL}{poporange!14!white}
\colorlet{popgreenL}{popgreen!12!white}
\colorlet{poppurpleL}{poppurple!12!white}
\colorlet{poppinkL}{poppink!12!white}
\colorlet{popgoldL}{popgold!14!white}

\pagecolor{popcream}

% ============================================================
% Typographie
% ============================================================
\defaultfontfeatures{Ligatures=TeX}
\setmainfont{PlusJakartaSans-Medium.ttf}[Path=../../../fonts/,BoldFont=PlusJakartaSans-Bold.ttf]
% Maths en sans-serif (Fira Math) avec chiffres et lettres droites de Plus
% Jakarta, pour que les formules se fondent dans le texte.
\usepackage[math-style=ISO,bold-style=ISO]{unicode-math}
\setmathfont{FiraMath-Regular.otf}[Path=../../../fonts/]
\setmathfont{PlusJakartaSans-Medium.ttf}[Path=../../../fonts/,range={up/num,up/{Latin,latin},\mathup}]
\usepackage{icomma} % virgule décimale sans espace en mode math
% Caractères mathématiques Unicode absents des polices de texte (Plus Jakarta,
% Archivo Black) : rendus par leur équivalent mathématique, en texte comme en
% formule — sinon ils s'affichent en carré vide (ex. « Arithmétique dans ℤ »).
\usepackage{newunicodechar}
\newunicodechar{ℝ}{\ensuremath{\mathbb{R}}}
\newunicodechar{ℤ}{\ensuremath{\mathbb{Z}}}
\newunicodechar{ℂ}{\ensuremath{\mathbb{C}}}
\newunicodechar{ℕ}{\ensuremath{\mathbb{N}}}
\newunicodechar{ℚ}{\ensuremath{\mathbb{Q}}}
\newunicodechar{𝔻}{\ensuremath{\mathbb{D}}}
\newunicodechar{α}{\ensuremath{\alpha}}
\newunicodechar{β}{\ensuremath{\beta}}
\newunicodechar{γ}{\ensuremath{\gamma}}
\newunicodechar{δ}{\ensuremath{\delta}}
\newunicodechar{ε}{\ensuremath{\varepsilon}}
\newunicodechar{θ}{\ensuremath{\theta}}
\newunicodechar{λ}{\ensuremath{\lambda}}
\newunicodechar{μ}{\ensuremath{\mu}}
\newunicodechar{σ}{\ensuremath{\sigma}}
\newunicodechar{τ}{\ensuremath{\tau}}
\newunicodechar{φ}{\ensuremath{\varphi}}
\newunicodechar{ψ}{\ensuremath{\psi}}
\newunicodechar{ω}{\ensuremath{\omega}}
\newunicodechar{Σ}{\ensuremath{\Sigma}}
% majuscules produites par \MakeUppercase dans les titres (α-aminés -> Α)
\newunicodechar{Α}{\ensuremath{\alpha}}
\newunicodechar{Β}{\ensuremath{\beta}}
\newunicodechar{Γ}{\ensuremath{\Gamma}}
\newunicodechar{Δ}{\ensuremath{\Delta}}
\newunicodechar{Ε}{\ensuremath{\varepsilon}}
\newunicodechar{Θ}{\ensuremath{\Theta}}
\newunicodechar{Λ}{\ensuremath{\Lambda}}
\newunicodechar{Μ}{\ensuremath{\mu}}
\newunicodechar{Τ}{\ensuremath{\tau}}
\newunicodechar{Φ}{\ensuremath{\Phi}}
\newunicodechar{Ψ}{\ensuremath{\Psi}}
\newunicodechar{Ω}{\ensuremath{\Omega}}
\newunicodechar{↦}{\ensuremath{\mapsto}}
\newunicodechar{∈}{\ensuremath{\in}}
\newunicodechar{∉}{\ensuremath{\notin}}
\newunicodechar{∧}{\ensuremath{\wedge}}
\newunicodechar{⇔}{\ensuremath{\Leftrightarrow}}
\newunicodechar{⟺}{\ensuremath{\Longleftrightarrow}}
\newunicodechar{⇒}{\ensuremath{\Rightarrow}}
\newunicodechar{⟹}{\ensuremath{\Longrightarrow}}
\newunicodechar{∘}{\ensuremath{\circ}}
\newunicodechar{∪}{\ensuremath{\cup}}
\newunicodechar{∩}{\ensuremath{\cap}}
\newunicodechar{≡}{\ensuremath{\equiv}}
\newunicodechar{⊂}{\ensuremath{\subset}}
\newunicodechar{⊥}{\ensuremath{\perp}}
\newunicodechar{∃}{\ensuremath{\exists}}
\newunicodechar{∀}{\ensuremath{\forall}}
\newunicodechar{∛}{\ensuremath{\sqrt[3]{\,}}}
\newunicodechar{∜}{\ensuremath{\sqrt[4]{\,}}}
\newunicodechar{⃗}{\ensuremath{{}^{\to}}}
\newunicodechar{ⁿ}{\ifmmode^{n}\else\textsuperscript{\ensuremath{n}}\fi}
\newunicodechar{ˣ}{\ifmmode^{x}\else\textsuperscript{\ensuremath{x}}\fi}
\newunicodechar{ᵇ}{\ifmmode^{b}\else\textsuperscript{\ensuremath{b}}\fi}
\newunicodechar{ʳ}{\ifmmode^{r}\else\textsuperscript{\ensuremath{r}}\fi}
\newunicodechar{ᵘ}{\ifmmode^{u}\else\textsuperscript{\ensuremath{u}}\fi}
\newunicodechar{ʸ}{\ifmmode^{y}\else\textsuperscript{\ensuremath{y}}\fi}
\newunicodechar{ᵃ}{\ifmmode^{a}\else\textsuperscript{\ensuremath{a}}\fi}
\newunicodechar{ᵉ}{\ifmmode^{e}\else\textsuperscript{\ensuremath{e}}\fi}
\newunicodechar{ᵏ}{\ifmmode^{k}\else\textsuperscript{\ensuremath{k}}\fi}
\newunicodechar{ᵖ}{\ifmmode^{p}\else\textsuperscript{\ensuremath{p}}\fi}
\newunicodechar{ᴺ}{\ifmmode^{N}\else\textsuperscript{\ensuremath{N}}\fi}
\newunicodechar{ᶜ}{\ifmmode^{c}\else\textsuperscript{\ensuremath{c}}\fi}
\newunicodechar{ʰ}{\ifmmode^{h}\else\textsuperscript{\ensuremath{h}}\fi}
\newunicodechar{ᵠ}{\ifmmode^{q}\else\textsuperscript{\ensuremath{q}}\fi}
\newunicodechar{ₙ}{\ifmmode_{n}\else\textsubscript{\ensuremath{n}}\fi}
\newunicodechar{ₐ}{\ifmmode_{a}\else\textsubscript{\ensuremath{a}}\fi}
\newunicodechar{ᵢ}{\ifmmode_{i}\else\textsubscript{\ensuremath{i}}\fi}
\newunicodechar{ᵣ}{\ifmmode_{r}\else\textsubscript{\ensuremath{r}}\fi}
\newunicodechar{ₖ}{\ifmmode_{k}\else\textsubscript{\ensuremath{k}}\fi}
\newunicodechar{ᵦ}{\ifmmode_{\beta}\else\textsubscript{\ensuremath{\beta}}\fi}
\newunicodechar{ₓ}{\ifmmode_{x}\else\textsubscript{\ensuremath{x}}\fi}
\newfontfamily\popdisplay{ArchivoBlack-Regular.ttf}[Path=../../../fonts/]
\newfontfamily\poplogo{SpaceGrotesk-Bold.ttf}[Path=../../../fonts/]
\newfontfamily\popsemi{PlusJakartaSans-SemiBold.ttf}[Path=../../../fonts/]
\newfontfamily\popxbold{PlusJakartaSans-ExtraBold.ttf}[Path=../../../fonts/]
\newfontfamily\popxboldls{PlusJakartaSans-ExtraBold.ttf}[Path=../../../fonts/,LetterSpace=3.0]

\setlist[enumerate]{leftmargin=1.7em,itemsep=5pt,topsep=4pt}
\setlist[itemize]{leftmargin=1.7em,itemsep=4pt,topsep=4pt,label={\color{poporange}\textbullet}}
\setlength{\parindent}{0pt}
\setlength{\parskip}{5pt}
\renewcommand{\arraystretch}{1.25}

% pgfplots : style Pop par défaut
\pgfplotsset{
  every axis/.append style={axis line style={popink,line width=1pt},tick style={popink},
    grid style={popline,line width=0.5pt},label style={font=\small},tick label style={font=\footnotesize}},
  popcurve/.style={poporange,line width=1.8pt,no markers,smooth},
}
% Même style disponible dans un \draw TikZ simple (hors pgfplots)
\tikzset{popcurve/.style={poporange,line width=1.8pt}}
\tikzset{popgrid/.style={popline,very thin}}
\colorlet{popcurve}{poporange} % le nom du style est parfois utilisé comme couleur

% ============================================================
% Pill / squiggle
% ============================================================
\newcommand{\poppill}[3]{%
  \tikz[baseline=-0.55ex]\node[draw=popink,line width=1.2pt,rounded corners=2.6pt,%
    fill=#2,inner xsep=6.5pt,inner ysep=3pt]{%
    \popxboldls\fontsize{7}{7}\selectfont\color{#3}\MakeUppercase{#1}};%
}
\newcommand{\popsquiggle}[1]{%
  \tikz[baseline=-0.4ex]{
    \draw[#1,line width=2.6pt,line cap=round]
      plot [smooth] coordinates {(0,0.09) (0.34,0.01) (0.68,0.21) (1.02,0.01) (1.36,0.10)};
  }%
}
% Mot-clé surligné (marqueur orange sous le texte)
\newcommand{\cle}[1]{{\popxbold\color{popdark}#1}}

% ============================================================
% En-tête / pied
% ============================================================
\pagestyle{fancy}
\fancyhf{}
\renewcommand{\headrulewidth}{0pt}
\renewcommand{\footrulewidth}{0pt}
\lhead{\raisebox{-0.15ex}{\poplogo\fontsize{13}{13}\selectfont\color{popink}OnBuch\color{poporange}+}}
\rhead{\poppill{\DOCMATIERE\ \textbullet\ \DOCNIVEAU\ \textbullet\ Leçon \DOCLECON}{white}{popink}}
\lfoot{\popxbold\fontsize{7.6}{7.6}\selectfont\color{popmuted}\MakeUppercase{Cours signé OnBuch+}\ \textbullet\ {\popsemi\DOCTITRE}}
\rfoot{\tikz[baseline=-0.6ex]\node[rounded corners=8.5pt,fill=popink,minimum height=17pt,minimum width=17pt,inner xsep=5pt,inner sep=0]{\popxbold\fontsize{8}{8}\selectfont\color{popcream}\thepage\,/\,\pageref*{LastPage}};}

% ============================================================
% Titres
% ============================================================
\newcommand{\chaptitre}[2]{%
  \begin{tcolorbox}[enhanced,colback=poporange,colframe=popink,boxrule=2.6pt,arc=11pt,
      drop shadow=popink,left=15pt,right=15pt,top=12pt,bottom=11pt,before skip=3pt,after skip=18pt]
    \poppill{#2}{popink}{popcream}\par\vspace{9pt}
    {\raggedright\hyphenpenalty=10000\exhyphenpenalty=10000\popdisplay\fontsize{22}{24}\selectfont\color{popcream}\MakeUppercase{#1}\par}
  \end{tcolorbox}
}
% Section numérotée : pastille orange avec le numéro + titre Archivo
\newcounter{popsec}
\newcommand{\coursec}[1]{%
  \stepcounter{popsec}\setcounter{popsub}{0}%
  \par\vspace{16pt}\noindent
  \begin{minipage}{\linewidth}
  \tikz[baseline=-0.7ex]\node[draw=popink,line width=1.6pt,fill=poporange,rounded corners=4pt,minimum size=22pt,inner sep=0]{\popdisplay\fontsize{12}{12}\selectfont\color{popcream}\thepopsec};%
  \hspace{8pt}{\popdisplay\fontsize{15}{17}\selectfont\color{popink}\MakeUppercase{#1}}\par
  \vspace{4pt}\noindent{\color{poporange}\rule{36pt}{3.6pt}}
  \end{minipage}\par\nopagebreak\vspace{8pt}
}
\newcounter{popsub}
\newcommand{\courssub}[1]{%
  \stepcounter{popsub}%
  \par\vspace{9pt}\noindent{\popxbold\fontsize{11.5}{13}\selectfont\color{popink}{\color{poporange}\thepopsec.\thepopsub}\hspace{6pt}#1}\par\nopagebreak\vspace{4pt}
}

% ============================================================
% Boîtes pédagogiques — même recette : fond blanc, bordure épaisse colorée,
% ombre dure encre, badge plein. Titre optionnel [#1].
% ============================================================
\tcbset{popbase/.style={breakable,enhanced,colback=white,boxrule=2.3pt,arc=9pt,
  shadow={3pt}{-3pt}{0pt}{popink},left=14pt,right=14pt,top=11pt,bottom=11pt,before skip=11pt,after skip=15pt,
  fontupper=\popsemi\fontsize{10.6}{14.5}\selectfont\color{popink},
  fontlower=\popsemi\fontsize{10.6}{14.5}\selectfont\color{popink}}}
% Coche dessinée (le glyphe ✓ n'existe pas dans les polices du document)
\newcommand{\popcheck}[1]{\tikz[baseline=-0.5ex]\draw[#1,line width=1.6pt,line cap=round,line join=round](-0.13,0.0)--(-0.04,-0.09)--(0.13,0.1);}
\providecommand{\coche}{\popcheck{popgreen}}
\providecommand{\resultat}[1]{\par\smallskip\noindent\fcolorbox{popgreen}{popgreenL}{\parbox{\dimexpr\linewidth-2\fboxsep-2\fboxrule\relax}{\textbf{Résultat.} #1}}\par\smallskip}
\newcommand{\popboxhead}[3]{\poppill{#1}{#2}{white}\ifx\relax#3\relax\else\hspace{6pt}{\popxbold\fontsize{10.6}{12}\selectfont\color{popink}#3}\fi\par\vspace{7pt}}

\newtcolorbox{aretenir}[1][]{popbase,colframe=popgreen,before upper={\popboxhead{À retenir}{popgreen}{#1}}}
\newtcolorbox{exemplebox}[1][]{popbase,colframe=poporange,before upper={\popboxhead{Exemple}{poporange}{#1}}}
\newtcolorbox{definition}[1][]{popbase,colframe=popblue,before upper={\popboxhead{Définition}{popblue}{#1}}}
\newtcolorbox{propriete}[1][]{popbase,colframe=poppurple,before upper={\popboxhead{Propriété}{poppurple}{#1}}}
\newtcolorbox{methode}[1][]{popbase,colframe=popgold,before upper={\popboxhead{Méthode}{popgold}{#1}}}
\newtcolorbox{attention}[1][]{popbase,colframe=poppink,before upper={\popboxhead{Attention}{poppink}{#1}}}
\newtcolorbox{experience}[1][]{popbase,colframe=popblue,colback=popblueL,before upper={\popboxhead{Expérience}{popblue}{#1}}}
% « Le savais-tu ? » : carte violette pleine (contexte, culture, Cameroun)
\newtcolorbox{savaistu}[1][]{popbase,colback=poppurple,colframe=popink,
  fontupper=\popsemi\fontsize{10.4}{14.2}\selectfont\color{white},
  before upper={\poppill{Le savais-tu\,?}{white}{poppurple}\ifx\relax#1\relax\else\hspace{6pt}{\popxbold\color{white}#1}\fi\par\vspace{7pt}}}
% Exercice résolu : énoncé en haut, \tcblower puis la solution sur fond crème
\newcounter{popexo}\newcounter{popexr}
\newtcolorbox{exoresolu}[1][]{popbase,colframe=poporange,colbacklower=popcream,
  segmentation style={popink,line width=1.2pt,dashed},
  before upper={\stepcounter{popexr}\popboxhead{Exercice résolu \thepopexr}{poporange}{#1}},
  before lower={\poppill{Solution}{popink}{popcream}\par\vspace{6pt}}}
% Exercice (non résolu en place) — la correction est dans la partie Corrigés
\newtcolorbox{exercice}[1][]{popbase,colframe=popink,
  before upper={\stepcounter{popexo}\popboxhead{Exercice \thepopexo}{popink}{#1}}}
% Corrigé
\newtcolorbox{corrige}[1][]{popbase,colframe=popgreen,colback=popgreenL,
  before upper={\popboxhead{Corrigé}{popgreen}{#1}}}
% Objectifs / prérequis / bilan
\newtcolorbox{objectifs}{popbase,colframe=popink,colback=poporangeL,
  before upper={\popboxhead{Objectifs de la leçon}{popink}{}}}
\newtcolorbox{prerequis}{popbase,colframe=popink,colback=popblueL,
  before upper={\popboxhead{Prérequis}{popblue}{}}}
\newtcolorbox{bilan}{popbase,colframe=popink,colback=white,boxrule=2.8pt,
  before upper={\popboxhead{Fiche bilan}{poporange}{L'essentiel à connaître pour l'examen}}}
% Figure encadrée avec légende
\newtcolorbox{popfigure}[1][]{enhanced,breakable=false,colback=white,colframe=popline,boxrule=1.4pt,arc=8pt,
  left=10pt,right=10pt,top=10pt,bottom=8pt,before skip=10pt,after skip=12pt,halign=center,
  lower separated=false,halign lower=center,
  fontlower=\popsemi\fontsize{9}{11.5}\selectfont\color{popmuted},
  #1}
\newcommand{\legende}[1]{\par\vspace{6pt}{\popsemi\fontsize{9}{11.5}\selectfont\color{popmuted}#1}}

% ============================================================
% Signature OnBuch+
% ============================================================
\newcommand{\poplogoinline}[1]{{\poplogo\fontsize{#1}{#1}\selectfont\color{popink}OnBuch\color{poporange}+}}
\newcommand{\signatureonbuch}{%
  \begin{tcolorbox}[enhanced,colback=popink,colframe=popink,boxrule=2.2pt,arc=10pt,
      drop shadow=poporange,left=14pt,right=14pt,top=10pt,bottom=10pt,before skip=0pt,after skip=0pt]
    \tikz[baseline=-0.7ex]\node[circle,fill=popgreen,draw=popcream,line width=1.3pt,minimum size=22pt,inner sep=0]{\popcheck{white}};%
    \hspace{9pt}\begin{minipage}[c]{0.62\linewidth}
      {\popxbold\fontsize{12}{13}\selectfont\color{popcream}Signé\ \textbullet\ {\poplogo OnBuch\color{poporange}+}}\par\vspace{2pt}
      {\raggedright\popsemi\fontsize{8}{10.5}\selectfont\color{popline}\MakeUppercase{Équipe pédagogique OnBuch+}\\\MakeUppercase{Conforme au programme officiel MINESEC}\par}
    \end{minipage}\hfill
    {\poplogo\fontsize{22}{22}\selectfont\color{popcream}OnBuch\color{poporange}+}
  \end{tcolorbox}%
}

% ============================================================
% Page de garde
% #1 titre · #2 matière · #3 niveau · #4 module · #5 n° leçon · #6 durée ·
% #7 description / objectifs (texte ou itemize)
% ============================================================
\newcommand{\pagegarde}[7]{%
  \thispagestyle{empty}
  \newgeometry{a4paper,margin=1.7cm,top=1.6cm,bottom=1.4cm}%
  \begin{tikzpicture}[remember picture,overlay]
    \fill[poporange!18] ([xshift=-3.2cm,yshift=-3.6cm]current page.north east) circle (4.6cm);
    \fill[poppurple!14] ([xshift=2.4cm,yshift=7.6cm]current page.south west) circle (3.8cm);
  \end{tikzpicture}%
  {\poplogo\fontsize{30}{30}\selectfont\color{popink}OnBuch\color{poporange}+}\hfill
  \raisebox{0.6ex}{\poppill{Cours signé \textbullet\ Premium}{popink}{popcream}}\par
  \vspace{1pt}\popsquiggle{poppurple}\par
  \vspace{8pt}
  \begin{tcolorbox}[enhanced,colback=poporange,colframe=popink,boxrule=2.8pt,arc=13pt,
      drop shadow=popink,left=18pt,right=18pt,top=12pt,bottom=13pt,after skip=10pt]
    \poppill{#2 \textbullet\ #3}{popink}{popcream}\hspace{5pt}\poppill{#4}{white}{popink}\par\vspace{8pt}
    {\popdisplay\fontsize{14}{14}\selectfont\color{popink}LEÇON #5}\par\vspace{4pt}
    {\raggedright\hyphenpenalty=10000\exhyphenpenalty=10000\popdisplay\fontsize{25}{27}\selectfont\color{popcream}\MakeUppercase{#1}\par}
    \vspace{8pt}
    {\popxbold\fontsize{9.5}{9.5}\selectfont\color{popink}\MakeUppercase{Durée conseillée : #6}}
  \end{tcolorbox}
  \begin{tcolorbox}[enhanced,colback=white,colframe=popink,boxrule=2.2pt,arc=10pt,
      drop shadow=popink,left=16pt,right=16pt,top=10pt,bottom=10pt,after skip=8pt]
    \poppill{Dans cette leçon}{poppurple}{white}\par\vspace{5pt}
    \popsemi\fontsize{9.3}{12.6}\selectfont\color{popink}#7
  \end{tcolorbox}
  \noindent\begin{minipage}[t]{0.487\linewidth}
    \begin{tcolorbox}[enhanced,colback=popgreen,colframe=popink,boxrule=2.2pt,arc=10pt,
        drop shadow=popink,left=11pt,right=11pt,top=10pt,bottom=10pt,height=3.1cm,valign=center]
      \tcbox[colback=white,colframe=popink,boxrule=1.3pt,arc=5pt,boxsep=0pt,nobeforeafter,
        left=3pt,right=3pt,top=3pt,bottom=3pt,valign=center]{\qrcode[height=1.9cm]{https://play.google.com/store/apps/details?id=com.luvvix.onbuch&hl=fr}}%
      \hspace{8pt}\begin{minipage}[c]{0.5\linewidth}
        {\popdisplay\fontsize{11}{12.5}\selectfont\color{white}\MakeUppercase{Télécharge OnBuch}}\par\vspace{4pt}
        {\raggedright\popsemi\fontsize{8.4}{11}\selectfont\color{white}Gratuit \textbullet\ Google Play\\Tuteur IA, annales, concours, résultats\par}
      \end{minipage}
    \end{tcolorbox}
  \end{minipage}\hfill
  \begin{minipage}[t]{0.487\linewidth}
    \begin{tcolorbox}[enhanced,colback=poppurple,colframe=popink,boxrule=2.2pt,arc=10pt,
        drop shadow=popink,left=11pt,right=11pt,top=10pt,bottom=10pt,height=3.1cm,valign=center]
      \tikz[baseline=-0.7ex]\node[circle,fill=white,draw=popink,line width=1.3pt,minimum size=1.6cm,inner sep=0]{\popdisplay\fontsize{24}{24}\selectfont\color{poppurple}+};%
      \hspace{8pt}\begin{minipage}[c]{0.6\linewidth}
        {\popdisplay\fontsize{11}{12.5}\selectfont\color{white}\MakeUppercase{Passe à OnBuch+}}\par\vspace{4pt}
        {\raggedright\popsemi\fontsize{8.4}{11}\selectfont\color{white}Tous les cours signés, Tuteur IA illimité, fiches PDF — onglet OnBuch+ de l'app\par}
      \end{minipage}
    \end{tcolorbox}
  \end{minipage}
  \vspace{8pt}
  \popcontenu
  \vfill
  \signatureonbuch
  \restoregeometry
  \newpage
}

% Rangée « Ce cours contient » (page de garde)
\newcommand{\poptile}[3]{% #1 couleur #2 gros texte #3 légende
  \begin{tcolorbox}[enhanced,colback=#1,colframe=popink,boxrule=2pt,arc=9pt,shadow={2.5pt}{-2.5pt}{0pt}{popink},
      left=6pt,right=6pt,top=9pt,bottom=9pt,halign=center,valign=center,height=1.75cm,width=\linewidth,nobeforeafter]
    {\popdisplay\fontsize{12.5}{13}\selectfont\color{white}\MakeUppercase{#2}}\par\vspace{3pt}
    {\centering\popsemi\fontsize{7.8}{9.5}\selectfont\color{white}#3\par}
  \end{tcolorbox}}
\newcommand{\popcontenu}{%
  \par\noindent{\popxboldls\fontsize{7.5}{7.5}\selectfont\color{popmuted}CE COURS SIGNÉ CONTIENT}\par\vspace{7pt}
  \noindent\begin{minipage}[t]{0.235\linewidth}\poptile{popblue}{Cours}{complet, illustré, pas à pas}\end{minipage}\hfill
  \begin{minipage}[t]{0.235\linewidth}\poptile{poporange}{Méthodes}{et exercices résolus}\end{minipage}\hfill
  \begin{minipage}[t]{0.235\linewidth}\poptile{poppink}{Exercices}{type examen + corrigés}\end{minipage}\hfill
  \begin{minipage}[t]{0.235\linewidth}\poptile{popgreen}{Fiche bilan}{l'essentiel à retenir}\end{minipage}\par}

% Sommaire
\newtcolorbox{sommaire}{enhanced,colback=white,colframe=popink,boxrule=2.2pt,arc=10pt,
  drop shadow=popink,left=18pt,right=18pt,top=13pt,bottom=13pt,before skip=6pt,after skip=20pt,
  before upper={\poppill{Sommaire}{poppurple}{white}\par\vspace{9pt}\popsemi\fontsize{10.6}{14.5}\selectfont\color{popink}}}

% Signature de fin de cours — poussée tout en bas de la dernière page
\newcommand{\findecours}{%
  \par\vfill\nopagebreak
  \begin{center}\popsquiggle{poporange}\end{center}\vspace{4pt}
  \signatureonbuch
}
\AtBeginDocument{\def\llbracket{\mathopen{[\mkern-3.5mu[}}\def\rrbracket{\mathclose{]\mkern-3.5mu]}}} % intervalles d'entiers (glyphe ⟦ absent de la police)
% Glyphes absents de Fira Math : redéfinis à partir de glyphes disponibles
\AtBeginDocument{%
  \def\setminus{\mathbin{\backslash}}\let\smallsetminus\setminus
  \def\vdots{\mathord{\vbox{\baselineskip3.2pt\lineskiplimit0pt\kern4pt\hbox{.}\hbox{.}\hbox{.}}}}%
  \def\ddots{\mathinner{\mkern1mu\raise6.5pt\hbox{.}\mkern2mu\raise3.6pt\hbox{.}\mkern2mu\raise0.7pt\hbox{.}\mkern1mu}}%
  \def\triangle{\mathord{\scalebox{0.9}{$\symup{\Delta}$}}}%
  \def\nabla{\mathord{\raisebox{\depth}{\scalebox{1}[-1]{$\symup{\Delta}$}}}}%
  \def\checkmark{\popcheck{popdark}}%
  \def\star{\mathbin{\ast}}\def\bigstar{\mathord{\ast}}%
  \def\diamond{\mathbin{\scalebox{0.6}{$\Diamond$}}}%
  \def\bigcup{\mathop{\mathchoice{\raisebox{-0.3em}{\scalebox{1.7}{$\cup$}}}{\scalebox{1.25}{$\cup$}}{\cup}{\cup}}\displaylimits}%
  \def\bigcap{\mathop{\mathchoice{\raisebox{-0.3em}{\scalebox{1.7}{$\cap$}}}{\scalebox{1.25}{$\cap$}}{\cap}{\cap}}\displaylimits}%
  \def\lesseqqgtr{\mathrel{\lessgtr}}\def\bigoplus{\mathop{\oplus}\displaylimits}%
}
\providecommand{\ssi}{si et seulement si} % abréviation fréquente
\AtBeginDocument{\providecommand{\wideparen}{\overparen}} % arc de cercle

%% FIGFIX-BEGIN v5 — correctifs globaux des figures (docs/onbuch-generation/tools/figfix)
\usepackage{adjustbox}\usepackage{etoolbox}\tcbuselibrary{raster}
% toute figure TikZ/pgfplots plus large que la ligne est réduite à la largeur disponible
\BeforeBeginEnvironment{tikzpicture}{\begin{adjustbox}{max width=\linewidth}}
\AfterEndEnvironment{tikzpicture}{\end{adjustbox}}
% légendes pgfplots lisibles par-dessus les courbes
\pgfplotsset{every axis/.append style={legend style={fill=white,fill opacity=0.92,text opacity=1,draw=black!15,inner sep=3pt}}}
% étiquettes d'arêtes (syntaxe edge["..."]) avec fond blanc
\tikzset{every edge quotes/.append style={fill=white,inner sep=1.2pt}}
%% FIGFIX-END
\begin{document}

\pagegarde{Lutte contre les perturbations du système immunitaire}{SVT}{Tle D}{Module II : L'Éducation à la Santé — Mécanismes et dysfonctionnements du système immunitaire}{N11}{12 h}{Cette leçon étudie l'origine des cellules immunitaires, les mécanismes des réponses immunitaires non spécifique et spécifique (cellulaire et humorale), puis les dysfonctionnements du système immunitaire : maladies auto-immunes et VIH/Sida. Elle développe des compétences d'interprétation de résultats expérimentaux et des compétences psychosociales pour la lutte contre le VIH/Sida.
\vspace{4pt}
\begin{multicols}{2}\small
\begin{itemize}
\item À la fin de cette leçon, je sais identifier les lieux de formation et de maturation des cellules immunitaires et les structures de reconnaissance du non-soi.
\item À la fin de cette leçon, je sais rappeler les mécanismes de la réponse immunitaire non spécifique.
\item À la fin de cette leçon, je sais interpréter les résultats d'expériences sur les mécanismes de l'immunité contre la tuberculose et la diphtérie ou le tétanos.
\item À la fin de cette leçon, je sais expliquer les mécanismes aboutissant à la neutralisation et à l'élimination des antigènes (réponses à médiation cellulaire et humorale).
\item À la fin de cette leçon, je sais réaliser et utiliser des maquettes de la structure d'un anticorps et d'un complexe immun.
\item À la fin de cette leçon, je sais décrire un cas de dysfonctionnement du système immunitaire (maladie auto-immune, Sida).
\end{itemize}\end{multicols}}

\begin{sommaire}
\begin{multicols}{2}
\begin{enumerate}
\item Origine des cellules immunitaires et reconnaissance du non-soi
\item Rappels : la réponse immunitaire non spécifique
\item La réponse immunitaire spécifique à médiation humorale
\item La réponse immunitaire spécifique à médiation cellulaire
\item Les maladies auto-immunes
\item Le VIH/Sida : immunodéficience acquise et enjeux socioculturels
\item Activité d'intégration
\item Méthodes et exercices résolus
\item Exercices
\item Corrigés des exercices
\item Fiche bilan
\end{enumerate}
\end{multicols}
\end{sommaire}

\begin{objectifs}
\begin{itemize}
\item À la fin de cette leçon, je sais identifier les lieux de formation et de maturation des cellules immunitaires et les structures de reconnaissance du non-soi.
\item À la fin de cette leçon, je sais rappeler les mécanismes de la réponse immunitaire non spécifique.
\item À la fin de cette leçon, je sais interpréter les résultats d'expériences sur les mécanismes de l'immunité contre la tuberculose et la diphtérie ou le tétanos.
\item À la fin de cette leçon, je sais expliquer les mécanismes aboutissant à la neutralisation et à l'élimination des antigènes (réponses à médiation cellulaire et humorale).
\item À la fin de cette leçon, je sais réaliser et utiliser des maquettes de la structure d'un anticorps et d'un complexe immun.
\item À la fin de cette leçon, je sais décrire un cas de dysfonctionnement du système immunitaire (maladie auto-immune, Sida).
\item À la fin de cette leçon, je sais expliquer le mécanisme d'action du VIH et ses conséquences socioculturelles.
\item À la fin de cette leçon, je sais développer des compétences psychosociales pour vivre avec le VIH ou avec les personnes affectées.
\end{itemize}
\end{objectifs}

\begin{prerequis}
\begin{itemize}
\item Notion d'antigène, d'anticorps et de réaction antigène-anticorps (Première)
\item Barrières naturelles de l'organisme et réaction inflammatoire (Première)
\item Structure et rôle des cellules sanguines (hématies, leucocytes)
\item Notion de protéine et de spécificité de structure (liaison enzyme-substrat, récepteur-ligand)
\item Notion de virus et de cycle de multiplication virale (début d'année de Terminale)
\end{itemize}
\end{prerequis}

\newpage

\coursec{Origine des cellules immunitaires et reconnaissance du non-soi}

À l'hôpital de district de Bafoussam, deux patients illustrent des drames immunitaires opposés. Le premier, séropositif pour le VIH, enchaîne les infections banales — angines, diarrhées, mycoses cutanées — qui s'aggravent inexorablement, là où un organisme sain les résorbe en quelques jours. Le second, un adolescent de 15 ans, découvre un diabète sucré de type 1 (diabète juvénile) survenu brutalement, sans antécédent familial : ses propres défenses ont détruit les cellules \cle{$\beta$} pancréatiques productrices d'insuline. Pourquoi le corps du premier ne se défend-il plus ? Pourquoi celui du second s'attaque-t-il à lui-même ? Et pourquoi, après une vaccination antitétanique, une personne blessée par un clou rouillé reste-t-elle protégée des années plus tard ?

Pour répondre, il faut comprendre d'où viennent les soldats de notre armée intérieure — les cellules immunitaires — et comment elles apprennent à distinguer l'ami de l'ennemi. La problématique de cette section est : \emph{où naissent et où se forment les cellules immunitaires, et grâce à quelles structures reconnaissent-elles le non-soi ?}

\courssub{La moelle osseuse rouge : l'usine à cellules sanguines}

\textbf{Observation.} Chez un patient atteint de leucémie aiguë, les cellules sanguines anormales envahissent la moelle osseuse et le sang. Le traitement de référence consiste à détruire la moelle malade par irradiation totale du corps (ou chimiothérapie à haute dose), puis à greffer de la moelle osseuse saine d'un donneur compatible (allogreffe). Quelques semaines plus tard, le patient reconstitue \emph{tous} les lignages cellulaires sanguins : érythrocytes, thrombocytes, et l'ensemble des cellules immunitaires. Cette observation clinique prouve que toutes les cellules du sang dérivent d'un même réservoir : la \cle{moelle osseuse rouge}.

\begin{definition}[Cellules souches hématopoïétiques (CSH)]
Toutes les cellules immunitaires, comme l'ensemble des cellules sanguines, dérivent de \cle{cellules souches hématopoïétiques}, cellules indifférenciées, rares et quiescentes, nichées dans la \cle{moelle osseuse rouge} (os plats : sternum, côtes, bassin, vertèbres ; épiphyses des os longs). Ces cellules souches possèdent deux propriétés fondamentales :
\begin{itemize}
\item l'\emph{autorenouvellement} : division asymétrique produisant une cellule souche identique (maintien du stock) et une cellule engagée dans la différenciation ;
\item la \emph{potentialité multipotente} : capacité à engendrer l'ensemble des lignées cellulaires sanguines (hématopoïèse).
\end{itemize}
\end{definition}

\begin{savaistu}[Une production colossale et continue]
Chez l'adulte sain, la moelle osseuse rouge produit chaque jour environ \num{2e11} (200 milliards) globules rouges et \num{1e10} (10 milliards) globules blancs, soit plusieurs millions de cellules par seconde. Ce turn-over extrême explique la toxicité hématologique majeure des radiations ionisantes et des chimiothérapies anticancéreuses, qui ciblent les cellules en division rapide.
\end{savaistu}

À partir de la CSH, deux grandes voies de différenciation, appelées \cle{lignées}, se séparent précocement sous l'effet de facteurs de croissance (cytokines) spécifiques (GM-CSF, G-CSF, IL-3, IL-7, EPO, TPO...).

\textbf{La lignée myéloïde.} Elle se déroule entièrement dans la moelle osseuse et donne naissance aux effecteurs de l'immunité innée et aux cellules présentatrices d'antigène :
\begin{itemize}
\item les \cle{polynucléaires} (ou granulocytes) : \emph{neutrophiles} (phagocytose bactérienne, premiers recrues, durée de vie courte : quelques heures à quelques jours dans les tissus), \emph{éosinophiles} (défense anti-parasitaire, modulation allergique) et \emph{basophiles} (libération d'histamine, héparine, rôle dans l'allergie immédiate) ;
\item les \emph{monocytes}, cellules circulantes (demi-vie ~1-3 jours) qui migrent dans les tissus pour se différencier en \cle{macrophages} (phagocytose massive, présentation d'antigène, sécrétion de cytokines inflammatoires, réparation tissulaire) ou en \emph{cellules dendritiques} myéloïdes (sentinelles spécialisées dans la capture, le traitement et la présentation de l'antigène aux lymphocytes T naïfs) ;
\item les \emph{mastocytes} (issus d'un progéniteur myéloïde circulant, maturation dans les tissus conjonctifs et muqueux, acteurs clés de l'allergie).
\end{itemize}

\textbf{La lignée lymphoïde.} Elle produit les \cle{lymphocytes}, cellules effectrices de la réponse immunitaire spécifique (adaptative) :
\begin{itemize}
\item les \cle{lymphocytes B} (LB), qui naissent \emph{et} achèvent leur maturation (réarrangement des gènes des récepteurs, sélection négative) dans la \textbf{m}oelle \textbf{o}sseuse (mnémotechnique : « B » pour \emph{Bone marrow} ou \emph{Bursa} équivalent chez l'oiseau) ; ils expriment le BCR et, une activés, se différencient en plasmocytes sécréteurs d'anticorps ;
\item les \cle{lymphocytes T} (LT), qui naissent dans la moelle osseuse sous forme de progéniteurs immatures (thymocytes), puis migrent via le sang vers le \textbf{t}hymus (mnémotechnique : « T » pour \emph{Thymus}) où ils achèvent leur maturation (réarrangement TCR, sélection positive et négative) ; ils expriment le TCR et le co-récepteur CD4 ou CD8 ;
\item les lymphocytes \cle{NK} (\emph{Natural Killer}), effecteurs de l'immunité innée à activité cytotoxique spontanée contre cellules tumorales ou infectées (reconnaissance du « self manquant » : absence de CMH de classe I), sans récepteur réarrangé spécifique.
\end{itemize}

\begin{popfigure}
\begin{center}
\begin{tikzpicture}[
  font=\small,
  souche/.style={rectangle, rounded corners, draw=popdark, fill=popgoldL, thick, minimum width=4.5cm, minimum height=1.0cm, align=center},
  lignee/.style={rectangle, rounded corners, draw=popdark, fill=popblueL, thick, minimum width=4.0cm, minimum height=0.9cm, align=center},
  cell/.style={rectangle, rounded corners, draw=popmuted, fill=popcream, minimum width=3.6cm, minimum height=0.75cm, align=center},
  lieu/.style={font=\footnotesize\itshape, text=popmuted},
  fleche/.style={-{Stealth[length=2.5mm]}, thick, popdark}
]
\node[souche, align=center] (cs) at (0,0){Cellule souche\\ hématopoïétique (CSH)};
\node[lignee] (mye) at (-5,-2.2) {Lignée myéloïde};
\node[lignee] (lym) at (5,-2.2) {Lignée lymphoïde};
\node[cell, align=center] (pn) at (-7.2,-4.0){Polynucléaires\\ (Neutrophiles, Éosinophiles, Basophiles)};
\node[cell, align=center] (mono) at (-2.8,-4.0){Monocytes $\rightarrow$ Macrophages\\ Cellules dendritiques myéloïdes\\ Mastocytes (tissus)};
\node[cell, align=center] (lb) at (2.2,-4.0){Lymphocytes B\\ Lymphocytes NK};
\node[cell, align=center] (lt) at (7.8,-4.0){Lymphocytes T immatures\\ (Prothymocytes)};
\node[cell, fill=popgreenL, draw=popgreen, align=center] (ltm) at (7.8,-5.6){Lymphocytes T matures\\ (CD4$^+$ ou CD8$^+$)};
\draw[fleche] (cs.south west) -- (mye.north);
\draw[fleche] (cs.south east) -- (lym.north);
\draw[fleche] (mye) -- (pn);
\draw[fleche] (mye) -- (mono);
\draw[fleche] (lym) -- (lb);
\draw[fleche] (lym) -- (lt);
\draw[fleche, popgreen] (lt.south) -- node[right, font=\footnotesize, text=popgreen]{Migration sanguine} (ltm.north);
\node[lieu] at (0,1.1) {Moelle osseuse rouge (adulte : sternum, iliaques, côtes, vertèbres)};
\node[lieu] at (7.8,-6.5) {Maturation dans le thymus (involution après puberté)};
\draw[dashed, popmuted, rounded corners] (-9.5,-4.8) rectangle (9.5,1.4);
\node[lieu, anchor=north west] at (-9.5,1.4) {Tout se passe dans la moelle osseuse, sauf maturation des LT};
\end{tikzpicture}
\end{center}
\legende{Schéma simplifié de l'hématopoïèse. La CSH de la moelle osseuse rouge engendre deux lignées. La lignée myéloïde (polynucléaires, monocytes/macrophages, cellules dendritiques, mastocytes) achève sa différenciation sur place. La lignée lymphoïde donne les LB et NK (maturation médullaire) et les LT immatures qui migrent vers le thymus pour y mûrir.}
\end{popfigure}

\begin{attention}[Erreur fréquente au Bac : lieu de production vs lieu de maturation]
Les lymphocytes T \textbf{ne sont pas produits dans le thymus} ! Ils \emph{naissent} (différenciation à partir de la CSH) dans la moelle osseuse, puis \emph{migrent} vers le thymus pour y \emph{mûrir} (réarrangement TCR, sélection).
\begin{itemize}
\item \textbf{Formulation correcte :} « Les lymphocytes T dérivent de cellules souches de la moelle osseuse et achèvent leur maturation dans le thymus. »
\item \textbf{Formulation fautive (0 point) :} « Les lymphocytes T sont produits/fabriqués par le thymus. »
\end{itemize}
De même, les lymphocytes B mûrissent dans la moelle osseuse (chez l'adulte), pas dans la rate ni les ganglions.
\end{attention}

\courssub{Organes lymphoïdes primaires et secondaires}

L'architecture du système immunitaire repose sur une séparation spatiale entre la \emph{formation} des récepteurs (éducation) et leur \emph{utilisation} (réponse).

\begin{definition}[Organes lymphoïdes primaires (centraux) et secondaires (périphériques)]
\begin{itemize}
\item \textbf{Organes lymphoïdes primaires} : lieux de \textbf{génération et maturation} des lymphocytes à partir de progéniteurs. Chez l'Homme : \textbf{moelle osseuse rouge} (maturation des LB, site de production des LT immatures) et \textbf{thymus} (maturation des LT, sélection positive/négative). Le thymus, volumineux chez l'enfant, subit une involution graisseuse progressive après la puberté, mais conserve une activité résiduelle toute la vie.
\item \textbf{Organes lymphoïdes secondaires} : lieux de \textbf{concentration des antigènes} et de \textbf{rencontre lymphocyte--antigène} (activation de la réponse adaptative). Ils sont stratégiquement placés sur les voies de drainage : \textbf{rate} (filtre sanguin, antigènes sanguins), \textbf{ganglions lymphatiques} (filtres lymphatiques, antigènes tissulaires), \textbf{amygdales} et \textbf{plaques de Peyer} (muqueuses ORL et digestives, antigènes inhalés/ingérés).
\end{itemize}
\end{definition}

\begin{methode}[Identifier et classifier un organe lymphoïde sur un document (schéma, coupe, IRM)]
\begin{enumerate}
\item \textbf{Localiser} l'organe sur la figure (position anatomique : thymus médiastinal antérieur, rate loge splénique, ganglions aux carrefours vasculaires, moelle dans les cavités médullaires).
\item \textbf{Interroger sa fonction} : Est-ce un site de \emph{maturation} des lymphocytes naïfs (sortie de moelle/thymus) $\rightarrow$ \textbf{Primaire} ? Ou un site de \emph{réponse} (activation, prolifération clonale, gonflement = adénopathie) $\rightarrow$ \textbf{Secondaire} ?
\item \textbf{Justifier par la vascularisation} : Organes primaires = vascularisation artérielle nutritive, pas de vaisseaux lymphatiques afférents (thymus) ou sinusoïdes médullaires (moelle). Organes secondaires = vaisseaux lymphatiques afférents (apport antigène depuis tissus) et veine efférente (sortie lymphocytes activés) ; rate = artère splénique (sang) $\rightarrow$ sinusoïdes $\rightarrow$ veine splénique.
\item \textbf{Vérifier la cohérence clinique} : Adénopathie satellite d'une infection $\rightarrow$ ganglion = organe secondaire. Aplasie médullaire (aplasie de la moelle) $\rightarrow$ pancytopénie $\rightarrow$ moelle = organe primaire (usine).
\end{enumerate}
\end{methode}

\courssub{La preuve expérimentale du rôle du thymus : l'expérience de Miller (1961)}

Comment démontrer formellement qu'un organe est indispensable à une fonction ? Par l'ablation (expérience de perte de fonction) chez un modèle animal.

\begin{experience}[Thymectomie néonatale chez la souris (J.F.A.P. Miller, 1961) — Preuve du rôle du thymus dans l'immunité cellulaire]
\textbf{Protocole.}
\begin{enumerate}
\item Deux lots de souris nouveau-nées (système immunitaire immature, thymus actif) : Lot \textbf{Tx} (thymectomisées : ablation chirurgicale du thymus) ; Lot \textbf{Témoin} (sham-opérées : ouverture thoracique sans ablation).
\item À l'âge adulte (8-10 semaines), greffe cutanée allogénique (peau d'une souche génétiquement différente, ex. CBA $\rightarrow$ C57BL/6) sur le dos de chaque souris.
\item Suivi de la greffe et de l'état de santé.
\end{enumerate}

\textbf{Observations.}
\begin{itemize}
\item \textbf{Lot Témoin :} Rejet de la greffe en 10-15 jours (nécrose, chute du greffon). Souris en bonne santé par ailleurs.
\item \textbf{Lot Tx :} \textbf{Non-rejet de la greffe} (le greffon cicatrise, se vascularise, persiste indéfiniment = tolérance acquise). Mais les souris \textbf{maigrissent, présentent un thymus absent à l'autopsie, une rate et des ganglions atrophiés (dépourvus de zones paracorticales riches en LT), et meurent précocement} d'infections opportunistes (bactéries, virus, champignons) banales pour le témoin.
\end{itemize}

\textbf{Interprétation.}
\begin{itemize}
\item Le rejet de greffe est médié par les lymphocytes T cytotoxiques (LT CD8$^+$) reconnaissant le CMH de classe I du donneur. L'absence de rejet prouve l'absence de LT fonctionnels.
\item La sensibilité aux infections et l'atrophie des organes secondaires (zones T vides) confirment un défaut global de l'immunité spécifique à médiation cellulaire.
\item Comme la moelle osseuse est intacte, les progéniteurs LT sont produits mais ne peuvent pas mûrir sans thymus.
\end{itemize}

\textbf{Conclusion.} Le thymus est l'organe \textbf{indispensable et non redondant} de la maturation des lymphocytes T. Sans thymus, pas de LT matures $\rightarrow$ pas d'immunité spécifique cellulaire $\rightarrow$ pas de rejet de greffe, pas de défense antivirale/anticancéreuse efficace.
\end{experience}

\courssub{La reconnaissance du non-soi : CMH, TCR et BCR}

\textbf{Le problème de l'identification.} Comment une cellule immunitaire distingue-t-elle une cellule saine du soi d'une cellule infectée ou cancéreuse, ou une bactérie extracellulaire ? La réponse repose sur un système de « carte d'identité » moléculaire universel : le CMH.

\begin{definition}[Complexe Majeur d'Histocompatibilité (CMH) — Système HLA chez l'Homme]
Le \cle{CMH} (Human Leukocyte Antigen : \cle{HLA}) est un ensemble de gènes hautement polymorphes (codominance, hétérozygotie avantageuse) situés sur le chromosome 6 (région 6p21.3). Ils codent pour des \textbf{glycoprotéines membranaires} exprimées à la surface des cellules nucléées, qui présentent des fragments peptidiques (antigènes) aux lymphocytes T. Ce sont les \textbf{marqueurs du « soi »} : chaque individu (sauf jumeaux monozygotes) possède un profil CMH unique. L'incompatibilité CMH est la cause majeure du rejet de greffe d'organe solide.
\end{definition}

\begin{propriete}[Les deux classes de molécules du CMH — Distinction cruciale pour la réponse T]
\begin{center}
\begin{tabularx}{\linewidth}{|l|X|X|}\hline
\rowcolor{popblueL}
\textbf{Caractéristique} & \textbf{CMH de Classe I (HLA-A, -B, -C)} & \textbf{CMH de Classe II (HLA-DP, -DQ, -DR)} \\ \hline
\textbf{Expression cellulaire} & \textbf{Quasi-ubiquitaire} : toutes les cellules nucléées de l'organisme. & \textbf{Restreinte} : Cellules Présentatrices d'Antigène Professionnelles (CPA : cellules dendritiques, macrophages, lymphocytes B, thymocytes corticaux). \\ \hline
\textbf{Structure} & Chaîne lourde ($\alpha$, 44 kDa) + $\beta_2$-microglobuline (12 kDa, non codée par CMH). & Deux chaînes transmembranaires ($\alpha$ et $\beta$, ~30 kDa chacune), toutes deux codées par le CMH. \\ \hline
\textbf{Origine du peptide présenté} & \textbf{Endogène} : protéines synthétisées \emph{dans} la cellule (protéines virales, oncoprotéines, protéines self normales). Traitement dans le cytosol $\rightarrow$ protéasome $\rightarrow$ TAP $\rightarrow$ RE. & \textbf{Exogène} : protéines internalisées par \emph{endocytose/phagocytose} (bactéries, toxines, débris cellulaires). Traitement dans l'endosome/lysosome. \\ \hline
\textbf{Lymphocyte T cible} & \textbf{LT CD8$^+$} (cytotoxiques / effecteurs tueurs). & \textbf{LT CD4$^+$} (auxiliaires / helpers : Th1, Th2, Th17, Tfh, Treg). \\ \hline
\textbf{Fonction physiologique} & Surveillance de l'intégrité intracellulaire (virus, cancer). & Orchestration de la réponse immune (activation macrophages, aide aux B, inflammation). \\ \hline
\end{tabularx}
\end{center}
\end{propriete}

\begin{definition}[Antigène et Épitope]
Un \cle{antigène} (Ag) est toute substance (protéine, polysaccharide, glycolipide, nucléique) \emph{étrangère} (non-soi) ou perçue comme telle, capable de déclencher une réponse immunitaire spécifique. La portion précise reconnue par un récepteur lymphocytaire (BCR ou TCR) ou un anticorps est l'\cle{épitope} (ou déterminant antigénique). Un antigène porte généralement plusieurs épitopes distincts.
\end{definition}

\textbf{Les récepteurs clonaux des lymphocytes.} Chaque lymphocyte naïf exprime à sa surface \textbf{un unique type de récepteur} (spécificité clonale), généré par réarrangement somatique aléatoire des gènes V(D)J lors de la maturation.

\begin{itemize}
\item \textbf{Lymphocyte B :} \cle{BCR} (\emph{B Cell Receptor}). C'est un \emph{anticorps membranaire} (structure IgM ou IgD monomérique ancrée). Il reconnaît l'antigène \textbf{natif, conformationnel, sous forme soluble ou de surface} (protéine repliée, polysaccharide), \textbf{sans présentation par CMH}. La fixation de l'antigène au BCR déclenche l'internalisation, le traitement et la présentation sur CMH II (fonction de CPA).
\item \textbf{Lymphocyte T :} \cle{TCR} (\emph{T Cell Receptor}). Hétérodimère $\alpha\beta$ (majoritaire) ou $\gamma\delta$ (minoritaire, tissulaire). Il ne reconnaît \textbf{que le complexe peptide--CMH} (pCMH) à la surface d'une cellule cible ou présentatrice. Le co-récepteur \textbf{CD4} se lie au CMH II (stabilisation, recrutement kinase Lck) ; le co-récepteur \textbf{CD8} se lie au CMH I.
\end{itemize}

\begin{popfigure}
\begin{center}
\begin{tikzpicture}[font=\small, scale=0.9, transform shape]
% Styles
\tikzset{
  lympho/.style={circle, draw, thick, minimum size=2.2cm},
  noyau/.style={circle, draw, fill=popmuted!30, minimum size=0.9cm},
  recB/.style={popblue, very thick, -{Stealth[length=2mm]}},
  recT/.style={popgreen, very thick, -{Stealth[length=2mm]}},
  cmhI/.style={poporange, thick, rectangle, minimum width=0.15cm, minimum height=0.4cm},
  cmhII/.style={poppurple, thick, rectangle, minimum width=0.15cm, minimum height=0.4cm, rounded corners=1pt},
  pep/.style={fill=poppink, circle, minimum size=0.15cm},
  label/.style={font=\footnotesize, text=popdark},
  title/.style={font=\bfseries, below=0.2cm}
}
% Lymphocyte B (Gauche)
\node[lympho, fill=popblueL, draw=popblue] (LB) at (-4,0) {};
\node[noyau] at (-4,0) {};
\node[title] at (-4,-1.3) {Lymphocyte B};
\foreach \angle in {30, 90, 150} {
  \draw[recB] (-4,0) -- ++(\angle:1.4) node[label, pos=1.2, anchor=west, rotate=\angle-90] {BCR (IgM/IgD mem.)};
}
% Antigène libre
\draw[fill=poppink, draw=poppink] (-2.5,1.5) circle (0.15);
\node[label, text=poppink] at (-2.5,1.9) {Ag natif};
\draw[dashed, popmuted, ->] (-2.5,1.5) -- (-3.2,1.1);

% Lymphocyte T (Centre)
\node[lympho, fill=popgreenL, draw=popgreen] (LT) at (2,0) {};
\node[noyau] at (2,0) {};
\node[title] at (2,-1.3) {Lymphocyte T CD4$^+$};
\foreach \angle in {40, 90, 140} {
  \draw[recT] (2,0) -- ++(\angle:1.3) node[label, pos=1.15, anchor=west, rotate=\angle-90] {TCR $\alpha\beta$ + CD4};
}

% Cellule Cible / CPA (Droite)
\node[ellipse, draw=poporange, thick, fill=poporangeL, minimum width=3cm, minimum height=2cm] (Cible) at (7.5,0) {};
\node[noyau, fill=poporange!30] at (7.5,0) {};
\node[title] at (7.5,-1.3) {Cellule Cible / CPA};
% CMH I
\foreach \y in {0.6, 0.2, -0.2, -0.6} {
  \draw[cmhI] (6.3,\y) -- (6.6,\y);
  \node[pep] at (6.45,\y) {};
}
\node[label, text=poporange] at (5.2,0.9) {CMH I + peptide endogène};
\node[label, text=poporange] at (5.2,0.5) {(Virus, Cancer, Self)};
\node[label, font=\tiny] at (6.45,1.0) {CD8$^+$};
% CMH II
\foreach \y in {-1.0, -1.4} {
  \draw[cmhII] (6.3,\y) -- (6.6,\y);
  \node[pep] at (6.45,\y) {};
}
\node[label, text=poppurple] at (5.2,-1.0) {CMH II + peptide exogène};
\node[label, text=poppurple] at (5.2,-1.4) {(Bactéries, Toxines)};
\node[label, font=\tiny] at (6.45,-1.6) {CD4$^+$};
\end{tikzpicture}
\end{center}
\legende{Structures de reconnaissance du non-soi. \textbf{Gauche :} Lymphocyte B avec BCR (anticorps membranaires) fixant l'antigène natif (conformationnel). \textbf{Centre :} Lymphocyte T CD4$^+$ avec TCR + co-récepteur CD4. \textbf{Droite :} Cellule présentatrice (CPA) ou cellule cible nucléée. Elle expose des peptides sur CMH de classe I (peptides endogènes, reconnus par LT CD8$^+$) et, si CPA, sur CMH de classe II (peptides exogènes, reconnus par LT CD4$^+$). Le TCR ne voit \textbf{jamais} l'antigène libre, seulement le complexe pCMH.}
\end{popfigure}

\textbf{La tolérance du soi (Éducation lymphocytaire).} Au cours de la maturation dans les organes primaires, les lymphocytes subissent une sélection drastique :
\begin{enumerate}
\item \textbf{Sélection positive (thymus cortex) :} Seuls les thymocytes dont le TCR reconnaît \emph{faiblement} le CMH propre (self-CMH) survivent (utilité = restriction CMH). Les autres meurent par négligence.
\item \textbf{Sélection négative (thymus médulla / moelle osseuse pour LB) :} Les lymphocytes dont le récepteur reconnaît \emph{fortement} un peptide self présenté par le CMH self sont éliminés par apoptose (mort cellulaire programmée). Cela élimine les clones autoréactifs à haute affinité.
\end{enumerate}
Seuls survivent les lymphocytes \textbf{tolérants au soi} (affinité faible/nulle pour le self) mais \textbf{compétents pour le non-soi} (capables de reconnaître un peptide étranger dans le groove du CMH self). C'est la \cle{tolérance centrale}. Des mécanismes périphériques (Treg, anergie, ignorance) complètent ce contrôle. \textbf{L'échec de cette tolérance est à l'origine des maladies auto-immunes} (section 5).

\begin{aretenir}[L'essentiel de la section : Origine et Reconnaissance]
\begin{itemize}
\item \textbf{Origine unique :} Toutes les cellules immunitaires dérivent de la \textbf{cellule souche hématopoïétique (CSH)} de la \textbf{moelle osseuse rouge}.
\item \textbf{Deux lignées :} Myéloïde (innée : polynucléaires, monocytes/macrophages, cellules dendritiques, mastocytes) et Lymphoïde (adaptative : LB, LT, NK).
\item \textbf{Maturation distincte :} LB et NK $\rightarrow$ Moelle osseuse. LT $\rightarrow$ Thymus (migreation de progéniteurs médullaires). \textbf{Le thymus est indispensable à la maturation des LT (expérience de Miller).}
\item \textbf{Organes :} Primaires (Moelle, Thymus) = Maturation. Secondaires (Rate, Ganglions, Amygdales, Plaques de Peyer) = Rencontre Ag / Activation.
\item \textbf{Marqueur du Soi :} \textbf{CMH (HLA)}. Classe I (ubiquitaire, peptides endogènes $\rightarrow$ LT CD8$^+$). Classe II (CPA, peptides exogènes $\rightarrow$ LT CD4$^+$).
\item \textbf{Récepteurs :} \textbf{BCR} (LB, Ag natif, conformationnel). \textbf{TCR} (LT, peptide linéaire présenté par CMH self + co-récepteur CD4/CD8).
\item \textbf{Tolérance :} Acquise par sélection négative centrale (élimination clones autoréactifs). Échec = Auto-immunité.
\end{itemize}
\end{aretenir}

\begin{exoresolu}[Greffe de moelle allogénique et reconstitution immunitaire]
Un patient leucémique reçoit une allogreffe de moelle osseuse (donneur HLA-compatible) après aplasie médullaire totale (conditionnement). Trois mois post-greffe, l'analyse du sang périphérique montre : 100\% des polynucléaires neutrophiles, monocytes, lymphocytes B, lymphocytes T CD4$^+$ et CD8$^+$ portent les allèles HLA du donneur.
\begin{enumerate}
\item Que démontre cette observation sur l'origine des cellules immunitaires ?
\item Pourquoi les lymphocytes T du donneur, maturés dans le thymus du \emph{receveur}, ne déclenchent-ils pas de maladie du greffon contre l'hôte (GVH) massive immédiate ? (Indice : tolérance).
\end{enumerate}
\tcblower
\textbf{Solution.}
\begin{enumerate}
\item Puisque \textbf{tous} les lignages myéloïdes (polynucléaires, monocytes) et lymphoïdes (LB, LT) du receveur portent désormais le génotype HLA du donneur, ils proviennent exclusivement des cellules souches hématopoïétiques contenues dans la greffe. Cela démontre formellement que \textbf{la CSH de la moelle osseuse est le précurseur commun et exclusif de toutes les cellules immunitaires (et sanguines)}.
\item Les progéniteurs LT du donneur migrent de la moelle greffée vers le \textbf{thymus du receveur} (qui a survécu au conditionnement ou est résiduel). Là, ils subissent la \textbf{sélection positive sur le CMH du receveur} (restriction au self du receveur) et la \textbf{sélection négative sur les peptides self du receveur}. Ils deviennent ainsi \textbf{tolérants aux antigènes du receveur}. C'est pourquoi ils ne reconnaissent pas le receveur comme « non-soi » (pas de GVH immédiate). \textit{Note : La GVH survient si des LT matures du donneur (contenus dans la greffe) reconnaissent le receveur comme étranger, ou si la sélection thymique est imparfaite.}
\end{enumerate}
\end{exoresolu}

\begin{exercice}[Application : Conséquences d'ablations chirurgicales chez la souris]
On réalise chez trois lots de souris nouveau-nées (système immunitaire en développement) les interventions suivantes :
\begin{itemize}
\item Lot 1 : Thymectomie (ablation du thymus).
\item Lot 2 : Splénectomie (ablation de la rate).
\item Lot 3 : Témoin (sham-opérées).
\end{itemize}
À l'âge adulte (8 semaines), on injecte par voie intraveineuse une dose létale de \emph{Listeria monocytogenes} (bactérie intracellulaire facultative, contrôle par immunité cellulaire Th1/LT CD8$^+$).
\begin{enumerate}
\item Prédire l'issue de l'infection pour chaque lot (Survie / Décès) en justifiant par les mécanismes immunitaires mis en jeu.
\item Quel lot sert de contrôle pour l'effet de l'anesthésie/chirurgie ? Quel lot teste spécifiquement le rôle de la maturation des LT ?
\end{enumerate}
\end{exercice}

\begin{corrige}[Exercice : Thymectomie vs Splénectomie]
\begin{enumerate}
\item \textbf{Lot 1 (Thymectomisé) : Décès rapide.} Absence de thymus $\rightarrow$ pas de maturation des LT (ni CD4$^+$, ni CD8$^+$) $\rightarrow$ pas d'immunité cellulaire spécifique $\rightarrow$ incapacité à activer les macrophages (pas d'IFN-$\gamma$ des Th1) ni à lyser les cellules infectées (pas de LT CD8$^+$ cytotoxiques). La réponse innée (polynucléaires, macrophages non activés) est insuffisante contre \emph{Listeria}. Phénotype type « nude » / SCID-like.
\item \textbf{Lot 2 (Splénectomisé) : Survie probable (ou décès retardé/partiel).} La rate est un organe \textbf{secondaire} (filtre sanguin). Son ablation retire un site majeur de rencontre Ag/LT pour les antigènes sanguins, mais les \textbf{ganglions lymphatiques et la moelle/thymus sont intacts}. Les LT mûrissent normalement, sont activés dans les ganglions (drainage tissulaire) ou la moelle, et contrôlent l'infection. La rate est importante pour les antigènes sanguins (encapsulées), mais \emph{Listeria} étant intracellulaire, la réponse se fait surtout en tissus/ganglions.
\item \textbf{Lot 3 (Témoin) : Survie.} Immunité complète.
\item \textbf{Lot 3} contrôle la chirurgie/anesthésie. \textbf{Lot 1} teste spécifiquement le rôle du thymus (maturation LT).
\end{enumerate}
\end{corrige}

\coursec{Rappels : la réponse immunitaire non spécifique}

Dans la section précédente, nous avons vu que les cellules immunitaires (lymphocytes, polynucléaires, macrophages, cellules dendritiques, cellules NK) prennent toutes leur origine dans la \cle{moelle osseuse rouge} (hématopoïèse), et que leur capacité à reconnaître le \cle{non-soi} repose sur des récepteurs spécifiques (TCR, BCR) pour les lymphocytes, ou sur des récepteurs de reconnaissance de motifs (\cle{PRR}, \emph{Pattern Recognition Receptors}) pour les cellules de l'immunité innée. Avant que ces mécanismes sophistiqués et lents de l'immunité adaptative ne se mettent en place, l'organisme dispose d'une première ligne de défense immédiate, héritée génétiquement et identique chez tous les individus : l'\cle{immunité non spécifique} (ou \cle{immunité innée}). Elle ne discrimine pas finement entre les agents pathogènes, mais les reconnaît par des motifs moléculaires communs (\cle{PAMP}, \emph{Pathogen-Associated Molecular Patterns}) comme le LPS des bactéries Gram-négatif, le peptidoglycane, l'ARN double brin viral, l'ADN non méthylé CpG, le flagelline.

\courssub{Les barrières naturelles : le premier rempart}

La protection commence bien avant toute réaction cellulaire ou moléculaire interne. C'est une défense \emph{passive}, physique, chimique et microbiologique, qui empêche l'entrée des agents pathogènes.

\begin{exemplebox}[L'acidité gastrique, un piège chimique redoutable]
Au Cameroun comme ailleurs, la contamination alimentaire et hydrique est une cause majeure de diarrhées infectieuses (choléra, salmonellose, amibiase, fièvre typhoïde). Pourtant, la grande majorité des bactéries ingérées avec l'eau ou les aliments mal lavés ne parviennent jamais à l'intestin. Le \cle{suc gastrique}, sécrété par les cellules pariétales de l'estomac, maintient un pH très acide, \textbf{compris entre 1,5 et 2,5}. Cette acidité extrême dénature les protéines bactériennes (enzymes, toxines, flagelles) et lyse les membranes cytoplasmiques. Seuls quelques pathogènes adaptés (\emph{Helicobacter pylori}, qui neutralise l'acide par son uréase ; kystes d'amibes ; spores de \emph{Clostridium} ou \emph{Bacillus} ; virus entériques non enveloppés) survivent à ce barrage.
\end{exemplebox}

\begin{center}
\begin{tabularx}{\linewidth}{|l|X|X|}
\hline
\rowcolor{popblueL}
\textbf{Barrière} & \textbf{Localisation / Nature} & \textbf{Mécanisme de défense} \\
\hline
\cle{Peau} & Épiderme kératinisé, sec, pH acide (≈ 5,5) & Barrière physique infranchissable ; desquamation élimine les microbes ; sébum et sueur (acides gras, lysozyme) inhibent la croissance bactérienne. \\
\hline
\cle{Muqueuses} & Voies respiratoires, digestives, génito-urinaires & Épithélium cilié (piégeage + clairance mucociliaire) ; jonctions serrées ; production de mucus visqueux. \\
\hline
\cle{Sécrétions} & Larmes, salive, sueur, mucus, sperme & \cle{Lysozyme} (hydrolyse le peptidoglycane des parois bactériennes) ; \cle{lactoferrine} (séquestre le fer indispensable aux bactéries) ; \cle{défensines} (peptides antimicrobiens perforant les membranes). \\
\hline
\cle{Suc gastrique} & Estomac & pH ≈ 2 (HCl) : dénaturation protéique, lyse bactérienne. \\
\hline
\cle{Flore commensale (Microbiote)} & Peau, intestin, vagin, bouche & \textbf{Concurrence nutritionnelle} et \textbf{exclusion de niche} ; production de bactériocines et d'acides gras à chaîne courte (butyrate) abaissant le pH local ; stimulation de l'immunité locale (IgA sécrétoires, maturation des tissus lymphoïdes). \\
\hline
\end{tabularx}
\end{center}

\begin{attention}[Erreur fréquente : la flore commensale n'est pas « sale »]
Ne confonds pas \cle{flore commensale} (microbiote normal, symbiote, protecteur) et \cle{flore pathogène}. L'usage d'antibiotiques à large spectre (fréquent en automédication à Yaoundé ou Douala) détruit ce microbiote protecteur, favorisant les \emph{infections opportunistes} (candidose buccale ou vulvo-vaginale, diarrhée à \emph{Clostridioides difficile}). \textbf{La flore commensale est un acteur à part entière de l'immunité non spécifique.}
\end{attention}

\courssub{Les effecteurs solubles de l'immunité innée : le système du complément, les interférons et les cellules NK}

Outre les barrières et la phagocytose, l'immunité innée mobilise des armes solubles et des lymphocytes innés.

\begin{definition}[Système du complément]
Ensemble d'une trentaine de protéines plasmatiques (facteurs C1 à C9, facteurs B, D, P, properdine, régulateurs) circulantes sous forme inactive (zymogènes). Leur activation en cascade aboutit à trois effets effecteurs majeurs : \textbf{opsonisation} (C3b), \textbf{chimiotactisme et inflammation} (C3a, C5a : anaphylatoxines), \textbf{lyse cellulaire} (Complexe d'Attaque Membranaire MAC = C5b-9).
\end{definition}

\begin{propriete}[Trois voies d'activation du complément (Admis au Bac)]
\begin{itemize}
\item \textbf{Voie classique :} Déclenchée par la fixation de C1q sur des \textbf{complexes immuns Ag-Ac (IgG, IgM)}. Pont entre immunité innée et adaptative.
\item \textbf{Voie alternative :} Activation spontanée sur surfaces étrangères (paroi bactérienne, LPS) \textbf{sans anticorps}. Amplification majeure. Facteur B, D, Properdine.
\item \textbf{Voie des lectines :} Déclenchée par la fixation de la \textbf{MBL (Mannose Binding Lectin)} ou des ficolines sur des motifs glucidiques (mannose) des pathogènes. Type « PRR soluble ».
\end{itemize}
\textbf{Point de convergence :} Clivage de C3 en C3a + C3b par les C3 convertases (C4b2a ou C3bBb). C3b $\rightarrow$ opsonisation / formation C5 convertase $\rightarrow$ C5a (chimiotactisme) + C5b $\rightarrow$ MAC (lyse).
\end{propriete}

\begin{definition}[Interférons de type I (IFN-$\alpha$, IFN-$\beta$)]
Cytokines produites précocement par toute cellule infectée par un virus (détection ARNdb par RIG-I/MDA5 ou ADNc par cGAS). Elles induisent un \textbf{état antiviral} dans les cellules voisines (synthèse de protéines bloquant la réplication virale : PKR, OAS/RNase L, Mx) et activent les \cle{cellules NK} (Natural Killers) et les macrophages.
\end{definition}

\begin{definition}[Cellules NK (Natural Killers)]
Lymphocytes innés (CD3$^-$, CD56$^+$, CD16$^+$) issus de la moelle osseuse. Reconnaissent les cellules cibles \textbf{sans CMH de classe I} (hypothèse du « missing self » : évasion virale/tumorale) ou porteuses d'anticorps (ADCC via CD16/Fc$\gamma$RIII). Tuent par \textbf{perforines/granzymes} (apoptose) et sécrètent de l'IFN-$\gamma$ (activation macrophages).
\end{definition}

\begin{savaistu}[Le complément et le diagnostic : le test de Coombs]
La recherche de C3b déposé sur les globules rouges (test de Coombs direct) permet de diagnostiquer les anémies hémolytiques auto-immunes (maladie auto-immune) ou la maladie hémolytique du nouveau-né (incompatibilité Rhésus). Au Cameroun, ce test est réalisé dans les laboratoires des hôpitaux de référence (HGOPY, HGOD, CHU) pour la prise en charge des anémies chroniques.

\begin{experience}[Mise en évidence de l'activité hémolytique du complément (Technique de Mayer / Hémolyse en gelose)]
\textbf{Principe :} Le complément lysé les globules rouges (GR) sensibilisés par des anticorps (voie classique).
\textbf{Matériel :} Sérum frais (source de complément), GR de mouton + anti-GR mouton (hémolysine), gélose, puits.
\textbf{Protocole :} Mélanger sérum dilué + hémolysine + GR. Incuber 37°C. Mesurer l'hémolyse (surnageant coloré) ou déposer en puits sur gélose + GR : zone de lyse (halo transparent) proportionnelle à l'activité complémentaire.
\textbf{Interprétation :} Sérum chauffé 56°C 30 min (inactivation complément) $\rightarrow$ pas de lyse. Le complément est thermolabile.
\end{experience}
\end{savaistu}

\courssub{La phagocytose : le « manger » cellulaire}

Si les barrières sont franchies (plaie, inhalation profonde, inoculation par moustique), la deuxième ligne de défense cellulaire intervient : la \cle{phagocytose}. Découverte par \textbf{Élie Metchnikoff} (Prix Nobel 1908) sur des larves d'étoile de mer, elle est réalisée par les \cle{cellules phagocytaires professionnelles} : les \cle{polynucléaires neutrophiles} (PNN, ou granulocytes neutrophiles), premiers arrivés sur le site d'infection (6 à 12 h), et les \cle{macrophages} tissulaires (issus des monocytes sanguins), plus lents mais plus endurants et capables de présenter l'antigène.

\begin{definition}[Phagocytose]
Mécanisme actif d'ingestion et de digestion intracellulaire de particules étrangères (bactéries, champignons, débris cellulaires, poussières) par une cellule phagocytaire, aboutissant à la formation d'une vacuole digestive (\cle{phagolysosome}) où des enzymes hydrolytiques détruisent la proie.
\end{definition}

\begin{experience}[L'observation fondatrice de Metchnikoff (1882)]
\textbf{Observation :} Metchnikoff introduit une épine de rosier (corps étranger stérile) dans le corps transparent de larves d'étoile de mer. Au microscope, il voit des cellules amiboïdes (mésenchymateuses) s'agglutiner autour de l'épine et l'englober.
\textbf{Hypothèse :} Ces cellules « mangent » les corps étrangers pour protéger l'organisme.
\textbf{Expérience décisive :} Il injecte des bactéries vivantes (\emph{Bacillus anthracis}) dans des puces d'eau (\emph{Daphnia}). Les cellules phagocytaires les ingèrent et les digèrent. Si les phagocytes sont saturés ou absents, l'animal meurt.
\textbf{Conclusion :} La phagocytose est une réaction de défense cellulaire universelle, conservée de l'invertébré au vertébré. Chez l'homme, ce sont les PNN et les macrophages qui assurent ce rôle.
\end{experience}

\begin{popfigure}
\begin{tikzpicture}[scale=0.9, every node/.style={font=\small}]
% Styles
\tikzset{
    cell/.style={draw=popdark, thick, fill=popblue!10, rounded corners=5pt, minimum width=3cm, minimum height=2.5cm},
    bact/.style={draw=poporange, thick, fill=poporange!30, circle, minimum size=0.5cm},
    lysosome/.style={draw=poppurple, thick, fill=poppurple!20, ellipse, minimum width=1cm, minimum height=0.6cm},
    arrow/.style={->, thick, popdark, >=Stealth},
    label/.style={font=\small\bfseries, popdark},
    stepbox/.style={font=\small\bfseries, popblue, anchor=north}
}

% --- Vignette 1 : Adhérence ---
\begin{scope}[xshift=0cm, yshift=0cm]
\node[cell] (c1) at (1.5,1.25) {};
\node[bact] (b1) at (1.5,2.8) {};
\draw[arrow] (b1.south) -- (c1.north) node[midway, right] {PRR / Opsonines};
\node[stepbox] at (1.5, 3.3) {1. Adhérence};
\node[font=\small\bfseries] at (1.5, -0.3) {PNN / Macrophage};
\node[font=\small\bfseries] at (1.5, 3.1) {Bactérie (PAMP)};
\draw[popline, thick] (0,0) rectangle (3,2.5);
\end{scope}

% --- Vignette 2 : Ingestion (Pseudopodes) ---
\begin{scope}[xshift=4cm, yshift=0cm]
\node[cell] (c2) at (1.5,1.25) {};
% Pseudopodes
\draw[popblue, thick, fill=popblue!20] (1.5,2.5) .. controls (1.0,2.0) and (0.8,1.5) .. (1.5,1.0) .. controls (2.2,1.5) and (2.0,2.0) .. (1.5,2.5);
\node[bact] (b2) at (1.5, 1.8) {};
\node[stepbox] at (1.5, 3.3) {2. Ingestion};
\node[font=\small\bfseries] at (1.5, -0.3) {Formation des pseudopodes};
\draw[popline, thick] (4,0) rectangle (7,2.5);
\end{scope}

% --- Vignette 3 : Phagosome ---
\begin{scope}[xshift=8cm, yshift=0cm]
\node[cell] (c3) at (1.5,1.25) {};
% Phagosome
\draw[popblue, thick, fill=popblue!30] (1.5,1.5) circle (0.5cm);
\node[bact] (b3) at (1.5, 1.5) {};
\node[lysosome] (ly3) at (2.8, 0.5) {};
\draw[arrow, dashed] (ly3) -- (1.5,1.0) node[midway, below, sloped] {Fusion};
\node[stepbox] at (1.5, 3.3) {3. Phagosome};
\node[font=\small\bfseries] at (1.5, -0.3) {Vacuole intracellulaire};
\node[font=\small\bfseries] at (3.5, 0.5) {Lysosome};
\draw[popline, thick] (8,0) rectangle (11,2.5);
\end{scope}

% --- Vignette 4 : Digestion ---
\begin{scope}[xshift=12cm, yshift=0cm]
\node[cell] (c4) at (1.5,1.25) {};
% Phagolysosome
\draw[poppurple, thick, fill=poppurple!30] (1.5,1.5) circle (0.55cm);
% Debris
\foreach \ang/\rad in {30/0.3, 150/0.25, 270/0.35} {
    \fill[popmuted] ($(1.5,1.5)+(\ang:\rad)$) circle (0.08cm);
}
\node[stepbox] at (1.5, 3.3) {4. Digestion};
\node[font=\small\bfseries] at (1.5, -0.3) {Phagolysosome : enzymes hydrolytiques};
\node[font=\small\bfseries] at (1.5, 1.5) {Résidus};
\draw[popline, thick] (12,0) rectangle (15,2.5);
\end{scope}
\end{tikzpicture}
\legende{Les quatre étapes majeures de la phagocytose : (1) Adhérence via récepteurs PRR (ex: récepteurs au mannose, TLR) ou opsonines (IgG, C3b, CRP) ; (2) Émission de pseudopodes et internalisation formant le phagosome ; (3) Fusion du phagosome avec les lysosomes (contenant lysozyme, cathepsines, NADPH oxydase) ; (4) Digestion dans le phagolysosome (pH acide, burst oxydatif) et élimination des résidus par exocytose.}
\end{popfigure}

\begin{methode}[Étapes biochimiques de la destruction intracellulaire (Burst oxydatif)]
\begin{enumerate}
\item \textbf{Acidification :} Pompe à protons (V-ATPase) abaisse le pH du phagolysosome à \textbf{≈ 4,5--5,0}, activant les hydrolases acides (protéases, nucléases, lipases).
\item \textbf{Burst oxydatif (Respiratory burst) :} Activation de la \cle{NADPH oxydase} (complexe NOX2) à la membrane du phagosome $\rightarrow$ production d'\textbf{anion superoxyde} ($O_2^{\bullet -}$).
\item \textbf{Espèces réactives de l'oxygène (ERO) :} Dismutation spontanée ou par SOD $\rightarrow$ $H_2O_2$ (eau oxygénée). Réaction de Fenton ($Fe^{2+}$) $\rightarrow$ \textbf{radical hydroxyle} ($HO^{\bullet}$, très toxique). Myélopéroxydase (MPO) + $H_2O_2$ + $Cl^-$ $\rightarrow$ \textbf{acide hypochloreux} ($HOCl$, eau de Javel endogène).
\item \textbf{Enzymes non oxydatives :} Lysozyme, élastase, cathépsine G, défensines, protéase-3.
\end{enumerate}
\textbf{Résultat :} Destruction de la plupart des bactéries en 10--30 min. \emph{Mycobacterium tuberculosis} inhibe la fusion phagosome-lysosome (blocage maturation) et résiste aux ERO (catalase, superoxyde dismutase, mycolates) $\rightarrow$ survie intracellulaire (granulome tuberculeux, voir Section 5/6).
\end{methode}

\begin{attention}[Le pus : signe de victoire, pas d'échec !]
\emph{« Docteur, j'ai du pus, c'est grave ? »} \textbf{Non, c'est souvent bon signe.} Le \cle{pus} (matière purulente) est un amas blanchâtre-jaunâtre composé de \textbf{polynucléaires neutrophiles morts} (après phagocytose), de bactéries digérées, de débris tissulaires et de lymphe. Sa présence indique une \textbf{phagocytose massive et efficace}. Une absence de pus chez un patient immunodéprimé (neutropénie, Sida avancé) face à une infection bactérienne est un signe de \textbf{gravité extrême} (septicémie foudroyante sans réaction locale).
\end{attention}

\courssub{La réaction inflammatoire : l'alarme tissulaire}

L'inflammation est la réponse vasculaire et cellulaire locale à une agression (infection, traumatisme, brûlure). Elle isole l'agent, recrute les effecteurs et répare le tissu. Ses quatre signes cardinaux (Celse, Ier siècle ; Galien) sont : \cle{Rubor} (rougeur), \cle{Calor} (chaleur), \cle{Tumor} (gonflement), \cle{Dolor} (douleur). Un cinquième signe, \cle{Functio laesa} (perte de fonction), est souvent ajouté.

\begin{popfigure}
\begin{tikzpicture}[scale=0.85, every node/.style={font=\small}]
% Vaisseau sanguin
\draw[popred, line width=3pt] (0,1.5) -- (14,1.5) node[right, pos=1.05] {Lumière du capillaire};
\draw[popred, line width=3pt] (0,0.5) -- (14,0.5);
% Endothélium
\foreach \x in {0.5, 1.5, 2.5, 3.5, 4.5, 5.5, 6.5, 7.5, 8.5, 9.5, 10.5, 11.5, 12.5, 13.5} {
    \draw[popdark, thick] (\x, 0.5) -- (\x, 1.5);
    \fill[popblue!30] (\x-0.05, 0.5) rectangle (\x+0.05, 1.5);
}
% Histamine release (Mastocyte)
\node[draw=poppurple, thick, fill=poppurple!20, circle, minimum size=0.8cm] (mast) at (3, -1.2) {};
\node[font=\small\bfseries] at (3, -2.2) {Mastocyte};
\draw[poppurple, ->, thick, decorate, decoration={snake, amplitude=1pt, segment length=4pt}] (mast.north) -- (3, 0.5) node[midway, right] {Histamine, Héparine};
% Vasodilatation
\draw[<->, poporange, thick] (2, 1.7) -- (4, 1.7) node[midway, above] {Vasodilatation};
% Diapedèse
\foreach \y in {0.7, 0.9, 1.1, 1.3} {
    \draw[popblue, ->, thick] (5.5, \y) .. controls (5.5, 1.6) and (6.5, 1.8) .. (6.5, 2.2);
}
\node[draw=popblue, fill=popblue!20, circle, minimum size=0.4cm] (pn1) at (6.5, 2.5) {};
\node[draw=popblue, fill=popblue!20, circle, minimum size=0.4cm] (pn2) at (7.2, 2.8) {};
\node[draw=popblue, fill=popblue!20, circle, minimum size=0.4cm] (pn3) at (5.8, 3.0) {};
\node[font=\small\bfseries] at (6.5, 3.3) {PNN (Diapedèse)};
% Chimiotactisme
\draw[popgreen, ->, thick, dashed] (7, 2.5) -- (9, 3.5) node[above right] {Chimiotactisme (C5a, LTB4, IL-8, fMLP)};
% Bactéries dans le tissu
\foreach \pos in {(9,3.5), (9.5,3.2), (10,3.6), (9.2,3.8)} {
    \node[draw=popred, fill=popred!40, ellipse, minimum width=0.35cm, minimum height=0.2cm] at \pos {};
}
\node[font=\small\bfseries] at (9.5, 4.3) {Foyer infectieux};
% Oedème
\draw[popmuted, decorate, decoration={zigzag, amplitude=2pt, segment length=4pt}, thick] (0, -0.2) -- (14, -0.2) node[midway, below] {Espace interstitiel : Œdème (exsudat)};
% Légendes des signes
\node[font=\small\bfseries, popblue, anchor=west] at (0.2, 2.5) {Calor / Rubor : Vasodilatation + afflux sanguin};
\node[font=\small\bfseries, popblue, anchor=west] at (0.2, 2.0) {Tumor : Exsudat (protéines, eau) + cellules};
\node[font=\small\bfseries, popblue, anchor=west] at (0.2, 1.5) {Dolor : Pression sur nerfs + médiateurs (bradykinine, PGE2)};
\end{tikzpicture}
\legende{Schéma de la réaction inflammatoire aiguë locale. 1. Dégranulation des mastocytes (histamine) $\rightarrow$ vasodilatation (Calor, Rubor) et augmentation de la perméabilité vasculaire (jonctions endothéliales écartées). 2. Sortie du plasma (exsudat riche en protéines : fibrinogène, anticorps, complément) $\rightarrow$ Œdème (Tumor). 3. Adhérence des PNN sur l'endothélium activé (molécules d'adhésion : selectines $\rightarrow$ roulement, puis intégrines $\rightarrow$ arrêt ferme) et diapedèse (migration transendothéliale). 4. Chimiotactisme vers le foyer infectieux guidé par le gradient de chimiotactines (C5a, LTB4, IL-8, fMLP).}
\end{popfigure}

\begin{propriete}[Caractères fondamentaux de l'immunité non spécifique (Admis au Bac)]
L'immunité innée (non spécifique) se caractérise par :
\begin{itemize}
\item \textbf{Imédiateté :} Réponse en \textbf{minutes à heures} (pas de phase de prolifération clonale nécessaire).
\item \textbf{Non-spécificité :} Même réponse quel que soit l'agresseur (reconnaissance de motifs conservés \emph{PAMP} par \emph{PRR} : TLR, NLR, RLR, CLR). Pas de réarrangement génique des récepteurs.
\item \textbf{Absence de mémoire immunologique :} Une seconde infection par le même agent déclenche une réponse \textbf{identique} en intensité et en cinétique (pas de réponse secondaire amplifiée). \textit{Nuance récente : existence d'un « entraînement immunitaire inné » (réprogrammation épigénétique des monocytes par BCG ou $\beta$-glucanes), mais non exigible au programme Terminale D.}
\item \textbf{Héréditaire :} Les gènes codant les PRR, le complément, les cytokines inflammatoires sont transmis selon les lois de Mendel (polymorphismes influençant la susceptibilité aux infections, ex: déficit en MBL, mutation TLR5).
\end{itemize}
\end{propriete}

\courssub{Limites de l'immunité non spécifique et transition vers l'immunité spécifique}

Malgré sa rapidité et sa puissance, l'immunité innée a des failles que les pathogènes exploitent :
\begin{itemize}
\item \textbf{Évasion de la phagocytose :} Capsules polysaccharidiques (\emph{Streptococcus pneumoniae}, \emph{Klebsiella pneumoniae}, \emph{Cryptococcus neoformans}) empêchent l'adhérence ; toxines leucocidines (\emph{Staphylococcus aureus} : PVL, Hlg) tuent les phagocytes ; inhibition de la fusion phagosome-lysosome (\emph{Mycobacterium tuberculosis}, \emph{Legionella pneumophila}, \emph{Salmonella} Typhi).
\item \textbf{Variabilité antigénique :} Changement de surface (Trypanosomes \emph{VSG}, \emph{Neisseria gonorrhoeae} \emph{pili}, Virus de la grippe \emph{drift/shift}, VIH) rend la reconnaissance par PRR moins efficace.
\item \textbf{Localisation intracellulaire stricte :} Virus, \emph{Chlamydia}, \emph{Rickettsia}, \emph{Plasmodium} (stade hépatique/érythrocytaire) se cachent hors d'atteinte des PNN et du complément.
\item \textbf{Absence de mémoire :} Pas de protection durable après guérison (sauf entraînement inné, voir supra).
\end{itemize}

\emph{Conséquence :} Lorsque l'immunité innée est submergée ou contournée, l'organisme \textbf{doit} activer l'immunité adaptative (spécifique). Le \textbf{pont} entre les deux mondes est assuré par les \cle{Cellules Présentatrices d'Antigène (CPA)} : principalement les \cle{cellules dendritiques} (les plus efficaces, seules à activer les LT naïfs), les \cle{macrophages} et les \cle{lymphocytes B (LB)}.

\begin{aretenir}[Le rôle pivot des Cellules Présentatrices d'Antigène (CPA)]
\begin{enumerate}
\item \textbf{Capture :} Au site d'infection, cellules dendritiques immatures et macrophages phagocytent l'agent (phagocytose, macropinocytose, endocytose récepteur-dépendante via BCR pour les LB).
\item \textbf{Maturation / Migration :} Stimulation par PAMP (via TLR) + cytokines inflammatoires (TNF-$\alpha$, IL-1$\beta$, IFN-$\gamma$) $\rightarrow$ maturation : \textbf{up-régulation des molécules CMH de classe II} et des \textbf{molécules co-stimulatrices (B7/CD80, CD86)}. Migration vers le ganglion lymphatique drainant (via chimiotactisme CCR7 $\leftarrow$ CCL19/21).
\item \textbf{Présentation :} Dans le ganglion, la CPA présente le \cle{peptide antigénique} (13--18 aa) inséré dans la rainure du \textbf{CMH de classe II} au \textbf{TCR d'un lymphocyte T CD4$^+$ naïf} (Signal 1).
\item \textbf{Co-stimulation :} Interaction \textbf{B7 (CPA) -- CD28 (LT CD4$^+$)} (Signal 2). Sans Signal 2 $\rightarrow$ anergie ou tolérance.
\item \textbf{Activation clonale :} Le LT CD4$^+$ activé prolifère et se différencie en \cle{Th1} (immunité cellulaire, activation macrophages via IFN-$\gamma$), \cle{Th2} (immunité humorale, éosinophiles via IL-4/5/13), \cle{Th17} (neutrophiles, barrières muqueuses via IL-17), \cle{Tfh} (aide aux LB dans centres germinatifs), \cle{Treg} (régulation, tolérance).
\end{enumerate}
C'est cette \textbf{présentation antigénique} qui initie la réponse spécifique (Sections 3 et 4).
\end{aretenir}

\courssub{Exercice résolu : Analyse d'un syndrome inflammatoire}

\begin{exoresolu}[Interprétation d'un tableau clinique et biologique]
\textbf{Énoncé :} Un enfant de 7 ans, non vacciné BCG, consulte au Centre de Santé de District de Bafia pour une fièvre à 39,5°C depuis 48 h, une toux grasse, une douleur thoracique droite et un abattement. L'examen retrouve un murmure vésiculaire diminué à la base droite. La numération-formule sanguine (NFS) montre : Leucocytes = 18 500 /mm$^3$ (N : 5 000--10 000) ; Polynucléaires neutrophiles = 15 700 /mm$^3$ (N : 2 000--7 000) dont 12\% de formes jeunes (bâtonnets) ; CRP = 120 mg/L (N < 6). L'hémoculture reste stérile à J3. Le test cutané à la tuberculine (IDR) est positif (14 mm).

\begin{enumerate}
\item Identifier les signes cliniques et biologiques traduisant une réaction inflammatoire aiguë.
\item Expliquer l'origine de l'hyperleucocytose neutrophile avec formules jeunes.
\item Pourquoi l'hémoculture peut-elle rester stérile malgré une infection bactérienne sévère ?
\item Quel type d'immunité est principalement mobilisé ici ? Justifier.
\end{enumerate}

\tcblower
\textbf{Solution détaillée :}
\begin{enumerate}
\item \textbf{Signes cliniques (Celse) :} Fièvre (Calor systémique), Douleur thoracique (Dolor), Abattement. \textbf{Signes biologiques :} Hyperleucocytose (18 500 /mm$^3$ soit 18,5 G/L) à prédominance polynucléaire (15 700 /mm$^3$), \textbf{formules jeunes (12\% > 5-10\% seuil pathologique)} signe de \textbf{moelle osseuse stimulée} (libération prématurée de réserves), CRP très élevée (120 mg/L) marqueur sensible de l'inflammation aiguë (synthèse hépatique induite par IL-6).
\item \textbf{Mécanisme :} L'infection pulmonaire (pneumonie bactérienne probable : \emph{S. pneumoniae}, \emph{H. influenzae}) libère des chimiotactines (C5a, IL-8, LTB4, fMLP) et des cytokines hématopoïétiques (G-CSF, GM-CSF). Le G-CSF stimule la \textbf{granulopoïèse} dans la moelle osseuse (phase de prolifération ~5-7 jours, mais libération immédiate du stock mature et du stock de réserve « bâtonnets »). L'afflux massif de PNN dans le sang (marge vasculaire $\rightarrow$ circulation) et leur migration vers le poumon (diapedèse) expliquent la neutrophilie avec « déviation à gauche » (formules jeunes).
\item \textbf{Hémoculture stérile :} Plusieurs causes : (1) Antibiothérapie précoce non déclarée (automédication fréquente au Cameroun) ; (2) Bactérie intracellulaire ou localisée (pneumonie sans bactériémie) ; (3) Volume de sang insuffisant (chez l'enfant, 1-3 mL seulement vs 20-30 mL recommandés) ; (4) Germe fastidieux ou anaérobie strict mal cultivé sur milieux standards.
\item \textbf{Immunité non spécifique (innée) :} Réponse en \textbf{heures} (fièvre, CRP, neutrophilie), \textbf{non spécifique} (même réponse pour toute bactérie pyogène), \textbf{sans mémoire} (enfant non vacciné BCG, pas de mémoire anti-pneumocoque). L'immunité spécifique (anticorps anti-pneumocoque, LT Th1/Th17) mettra 5-7 jours pour être efficace. L'IDR positif témoigne d'une infection tuberculeuse latente (immunité cellulaire spécifique anti-\emph{M. tuberculosis}), mais le tableau aigu est dominé par l'innée.
\end{enumerate}
\end{exoresolu}

\courssub{Synthèse et ouverture}

L'immunité non spécifique constitue le \textbf{socle évolutif} de notre défense. Elle est \textbf{rapide, stéréotypée, sans mémoire}, mais \textbf{insuffisante} face aux pathogènes intracellulaires, encapsulés ou variables. Son efficacité repose sur l'intégrité des barrières (peau, muqueuses, microbiote), la puissance de la phagocytose (PNN, macrophages), l'action du complément (opsonisation, lyse, inflammation), la réponse antivirale (IFN-$\alpha/\beta$, cellules NK) et l'alarme inflammatoire (vasodilatation, exsudat, recrutement). \textbf{Son point de bascule vers l'immunité spécifique est la présentation de l'antigène par les CPA (cellules dendritiques, macrophages, LB) aux lymphocytes T CD4$^+$ naïfs dans les organes lymphoïdes secondaires.} C'est ce dialogue crucial, alliant reconnaissance innée (PRR/PAMP) et reconnaissance adaptative (TCR/pCMH), que nous étudierons dans les deux prochaines sections : la réponse humorale (anticorps) et la réponse cellulaire (LT cytotoxiques, Th1).

\coursec{La réponse immunitaire spécifique à médiation humorale}

Après avoir vu comment la réponse non spécifique constitue une première ligne de défense rapide mais non sélective, nous abordons maintenant la réponse adaptative à médiation humorale. Celle-ci repose sur la production d'anticorps solubles capables de reconnaître spécifiquement un antigène, de le neutraliser et de marquer l'agent pathogène pour son élimination par les phagocytes. C'est cette forme d'immunité qui est transférable par le sérum, comme l'ont démontré les expériences fondatrices de la sérothérapie.

\courssub{Mise en évidence historique : la sérothérapie anti-diphtérique et anti-tétanique}

\begin{experience}[Expériences de Behring et Kitasato (1890) — Transfert de l'immunité par le sérum]
\textbf{Contexte et observation :} À la fin du XIX\textsuperscript{e} siècle, la diphtérie et le tétanos sont des fléaux mortels, notamment pour les enfants. Émile Roux et Alexandre Yersin ont mis en évidence que le bacille diphtérique (\emph{Corynebacterium diphtheriae}) sécrète une toxine soluble responsable des lésions.
\textbf{Protocole :}
\begin{enumerate}
\item Des cobayes sont immunisés par injections répétées de doses croissantes de toxine diphtérique atténuée (anatoxine) ou de toxine tétanique. Ils survivent à des doses létales pour un animal naïf.
\item Du sang est prélevé sur ces animaux immunisés, on laisse coaguler et on récupère le \textbf{sérum} (sang privé de ses cellules et du fibrinogène).
\item Ce sérum immun est injecté à des animaux non immunisés (cobayes, moutons, puis enfants atteints de diphtérie).
\item Ces animaux receveurs deviennent résistants à l'injection ultérieure de la toxine correspondante.
\end{enumerate}
\textbf{Résultats :} Le sérum d'un animal immunisé protège un animal non immunisé. La protection est \emph{immédiate} mais \emph{temporaire} (quelques semaines).
\textbf{Interprétation :} La substance protectrice circule dans le sérum (donc dans l'humeur, d'où « immunité humorale »). Il ne s'agit pas de cellules vivantes (le sérum est acellulaire) mais de molécules solubles : les \textbf{anticorps} (appelés alors \emph{antitoxines}).
\textbf{Conclusion :} L'immunité spécifique peut être transférée par le sérum ; elle est médiée par des facteurs humoraux (anticorps) capables de neutraliser spécifiquement la toxine. Emil von Behring recevra le premier prix Nobel de physiologie ou médecine en 1901 pour cette découverte.
\end{experience}

\begin{savaistu}[Le saviez-vous ?]
Emil von Behring (1854--1917), avec Shibasaburo Kitasato, a mis au point la \textbf{sérothérapie anti-diphtérique}. En 1894, le traitement du premier enfant diphtérique guéri par injection de sérum antitoxique marque la naissance de l'immunothérapie. Cette découverte a sauvé des millions de vies avant l'avènement des antibiotiques et des vaccins anatoxiniques. Au Cameroun, la sérothérapie antivenineuse (sérums anti-venins de vipères, cobras, mambas) repose sur le même principe : des anticorps polyvalents produits chez le cheval neutralisent les venins de serpents responsables d'envenimations graves en zone rurale.
\end{savaistu}

\courssub{Activation des lymphocytes B et différenciation}

La réponse humorale est portée par les \textbf{lymphocytes B} (maturation dans la moelle osseuse — \emph{Bursa equivalent} chez les mammifères). Chaque lymphocyte B naïve porte à sa surface environ $10^5$ récepteurs identiques, le \textbf{BCR} (\emph{B Cell Receptor}), qui est une immunoglobuline de type IgM (ou IgD) ancrée dans la membrane plasmique.

\begin{propriete}[Activation du lymphocyte B et coopération T-B]
L'activation complète d'un lymphocyte B naïf nécessite généralement \textbf{deux signaux} :
\begin{enumerate}
\item \textbf{Signal 1 (Reconnaissance antigénique) :} Le BCR se lie spécifiquement à l'antigène natif (conformationnel) — protéine, polysaccharide, etc. L'antigène est internalisé par endocytose médiée par le BCR, traité dans les endosomes/lysosomes, et des peptides antigéniques sont présentés à la surface du lymphocyte B sur des molécules du \textbf{CMHC de classe II}.
\item \textbf{Signal 2 (Coopération avec un lymphocyte T auxiliaire CD4\textsuperscript{+}) :} Un lymphocyte T auxiliaire (LT4) spécifique du même antigène (activé préalablement par une cellule présentatrice d'antigène professionnelle comme la cellule dendritique) reconnaît le complexe peptide-CMHC II sur le lymphocyte B. Cette interaction \emph{CD40--CD40L} et la sécrétion de cytokines (IL-4, IL-5, IL-6, IL-21) par le LT4 fournissent le second signal indispensable.
\end{enumerate}
\emph{Remarque :} Certains antigènes (polysaccharides bactériens à épitopes répétitifs) peuvent activer les lymphocytes B de manière \textbf{indépendante des T} (réponse TI), mais sans changement de classe (IgM seulement) et sans mémoire immunitaire efficace.
\end{propriete}

\begin{aretenir}[Différenciation du lymphocyte B activé]
Sous l'effet des signaux d'activation et des cytokines, le lymphocyte B activé prolifère (expansion clonale) et se différencie en deux populations :
\begin{itemize}
\item \textbf{Plasmocytes :} Usines à anticorps. Cytoplasme riche en réticulum endoplasmique rugueux et en appareil de Golgi. Durée de vie courte (quelques jours à semaines pour la plupart), mais certains deviennent des \emph{plasmocytes à longue vie} nichés dans la moelle osseuse et sécrètent des anticorps pendant des années.
\item \textbf{Lymphocytes B mémoires :} Ne sécrètent pas d'anticorps. Expressent un BCR de haute affinité (grâce à l'hypermutation somatique et la sélection clonale dans les centres germinatifs des organes lymphoïdes secondaires). Durée de vie très longue (années, vie entière). Assurent la réponse secondaire.
\end{itemize}
\end{aretenir}

\courssub{Structure de l'anticorps (immunoglobuline)}

\begin{definition}[Anticorps (Immunoglobuline)]
Glycoprotéine sécrétée par les plasmocytes, capable de se fixer \textbf{spécifiquement} sur l'antigène (ou plus précisément sur son \emph{épitope} ou \emph{déterminant antigénique}) qui a induit sa production. Elle assure les fonctions effectrices de l'immunité humorale (neutralisation, opsonisation, activation du complément, etc.).
\end{definition}

L'unité de base est un hétérodimère de \textbf{quatre chaînes polypeptidiques} : deux \textbf{chaînes lourdes} (Lourdes, $\sim$ 50 kDa chacune) et deux \textbf{chaînes légères} (Légères, $\sim$ 25 kDa chacune), liées par des \textbf{ponts disulfure} (--S--S--). Il existe 5 classes (isotypes) d'anticorps chez l'homme (IgG, IgM, IgA, IgD, IgE) définies par leur chaîne lourde ($\gamma, \mu, \alpha, \delta, \varepsilon$), et deux types de chaînes légères ($\kappa$ et $\lambda$).

\begin{popfigure}
\begin{tikzpicture}[scale=0.95, every node/.style={font=\small}]
% Couleurs
\definecolor{heavy}{RGB}{0,90,180}
\definecolor{light}{RGB}{180,90,0}
\definecolor{disulf}{RGB}{100,100,100}
\definecolor{antigen}{RGB}{180,0,0}

% --- Chaînes lourdes (Heavy chains) ---
% Bras gauche (Fab gauche + Fc)
\draw[heavy, line width=3pt] (0,0) -- (0,1.5) -- (-1,2.5) -- (-1,5.5); % Montée VL/VH
\draw[heavy, line width=3pt] (0,0) -- (0,-1.5) -- (-1,-2.5) -- (-1,-5.5); % Descente (pas de Fab en bas, Fc seulement)
% Bras droit (Fab droit + Fc)
\draw[heavy, line width=3pt] (4,0) -- (4,1.5) -- (5,2.5) -- (5,5.5);
\draw[heavy, line width=3pt] (4,0) -- (4,-1.5) -- (5,-2.5) -- (5,-5.5);

% Tronc commun (Fc) - correction: les chaînes lourdes se rejoignent au niveau charnière
% On redessine proprement : deux chaînes lourdes symétriques formant un Y
% Chaîne lourde Gauche
\draw[heavy, line width=3pt] (2, -4) -- (2, -1) -- (1, 0) -- (1, 3) -- (0.5, 4.5); % Fc -> Hinge -> CH1 -> VH
% Chaîne lourde Droite
\draw[heavy, line width=3pt] (2, -4) -- (2, -1) -- (3, 0) -- (3, 3) -- (3.5, 4.5);

% --- Chaînes légères (Light chains) ---
% Gauche
\draw[light, line width=2pt] (0.5, 4.5) -- (0.5, 3) -- (1, 1.5) -- (1, 0); % VL -> CL -> (jointure CH1)
% Droite
\draw[light, line width=2pt] (3.5, 4.5) -- (3.5, 3) -- (3, 1.5) -- (3, 0);

% --- Domaines (Rectangles) ---
% Fab Gauche : VH (Heavy), VL (Light)
\draw[heavy, fill=heavy!20, thick] (0.0, 3.8) rectangle (1.0, 4.8) node[midway, font=\tiny] {VH};
\draw[light, fill=light!20, thick] (1.0, 3.8) rectangle (2.0, 4.8) node[midway, font=\tiny] {VL};
% Fab Droit : VH, VL
\draw[heavy, fill=heavy!20, thick] (3.0, 3.8) rectangle (4.0, 4.8) node[midway, font=\tiny] {VH};
\draw[light, fill=light!20, thick] (2.0, 3.8) rectangle (3.0, 4.8) node[midway, font=\tiny] {VL};

% Sites fixation antigène (Paratopes)
\node[antigen, circle, fill=antigen!20, draw, thick, minimum size=6pt] at (1, 5.1) {};
\node[antigen, circle, fill=antigen!20, draw, thick, minimum size=6pt] at (3, 5.1) {};
\draw[<-, antigen, thick] (1, 5.1) -- (0.5, 5.8) node[left, font=\tiny] {Paratope};
\draw[<-, antigen, thick] (3, 5.1) -- (3.5, 5.8) node[right, font=\tiny] {Paratope};

% Fab Constants : CH1, CL
\draw[heavy, fill=heavy!10, thick] (0.0, 2.3) rectangle (1.0, 3.8) node[midway, font=\tiny] {CH1};
\draw[light, fill=light!10, thick] (1.0, 2.3) rectangle (2.0, 3.8) node[midway, font=\tiny] {CL};
\draw[heavy, fill=heavy!10, thick] (3.0, 2.3) rectangle (4.0, 3.8) node[midway, font=\tiny] {CH1};
\draw[light, fill=light!10, thick] (2.0, 2.3) rectangle (3.0, 3.8) node[midway, font=\tiny] {CL};

% Région Charnière (Hinge)
\draw[heavy, line width=3pt] (1, 1.5) -- (3, 1.5);
\draw[disulf, dashed, thick] (1.5, 1.5) -- (1.5, 2.0) node[above, font=\tiny] {Ponts disulfure};
\draw[disulf, dashed, thick] (2.5, 1.5) -- (2.5, 2.0);

% Région Fc : CH2, CH3 (IgG a CH2, CH3 ; IgM/IgE ont CH4)
\draw[heavy, fill=heavy!10, thick] (1.5, 0.0) rectangle (2.5, 1.5) node[midway, font=\tiny] {CH2};
\draw[heavy, fill=heavy!10, thick] (1.5, -1.5) rectangle (2.5, 0.0) node[midway, font=\tiny] {CH3};
\draw[heavy, fill=heavy!10, thick] (1.5, -3.0) rectangle (2.5, -1.5) node[midway, font=\tiny] {CH3}; % IgG a 2 domaines constants Fc (CH2, CH3) -> souvent noté CH2, CH3. Pour simplifier on met CH2, CH3, CH3 ou CH2, CH3, CH4 selon isotype. Ici IgG standard.

% Glycosylation
\node[font=\small] at (2, -0.75) {$\star$ Glycosylation};

% Etiquettes globales
\node[heavy, font=\bfseries, align=center] at (-1.2, 1.5){Chaînes\\ lourdes};
\node[light, font=\bfseries, align=center] at (-1.2, -1.5){Chaînes\\ légères};
\node[font=\bfseries] at (2, 5.5) {Région Fab (Fragment antigen-binding)};
\node[font=\bfseries] at (2, -4.3) {Région Fc (Fragment crystallizable)};

% Flèches régions
\draw[<->, thick, popdark] (-1.5, 5.0) -- (-1.5, 3.8);
\draw[<->, thick, popdark] (-1.5, -3.8) -- (-1.5, -4.8);

\end{tikzpicture}
\legende{Schéma de la structure d'un anticorps de type IgG (monomère). Les domaines variables (VH, VL) forment le site de fixation de l'antigène (paratope) au sommet des deux bras Fab. La région Fc (domaines constants CH2--CH3) assure les fonctions effectrices (fixation sur récepteurs Fc des phagocytes, activation du complément, transport placentaire pour IgG). La région charnière confère la flexibilité nécessaire pour lier deux épitopes distants. La glycosylation (étoile) sur la région Fc est cruciale pour l'interaction avec les récepteurs Fc. Valence = 2 sites de fixation.}
\end{popfigure}

\begin{methode}[Réaliser une maquette de la structure d'un anticorps]
Pour visualiser la structure en Y et la symétrie :
\begin{enumerate}
\item Utiliser deux fils de fer rigides (chaînes lourdes) et deux fils plus fins (chaînes légères).
\item Les plier pour former deux bras identiques (Fab) et un tronc commun (Fc).
\item Relier les chaînes lourdes et légères par des attaches (ponts disulfure) au niveau de la charnière et dans les domaines constants.
\item Placer des billes de couleur distincte au sommet de chaque bras pour représenter les sites de fixation de l'antigène (paratopes).
\item Vérifier la valence : un IgG a \textbf{2 sites de fixation} (valence 2) ; un IgM pentamérique en a 10.
\end{enumerate}
\end{methode}

\courssub{Modes d'action des anticorps}

Les anticorps ne tuent pas directement les agents pathogènes. Ils agissent comme des \textbf{marqueurs} ou des \textbf{bloqueurs} pour faciliter leur élimination par les mécanismes effecteurs de l'organisme (phagocytose, complément, cytotoxicité cellulaire).

\begin{propriete}[Principaux modes d'action des anticorps]
\begin{enumerate}
\item \textbf{Neutralisation :} L'anticorps se fixe sur la partie fonctionnelle d'une toxine (ex: toxine diphtérique, tétanique) ou sur les protéines d'enveloppe d'un virus (ex: protéine Spike du SARS-CoV-2, gp120 du VIH), empêchant leur fixation sur les récepteurs des cellules cibles. La toxine ou le virion devient inoffensif.
\item \textbf{Agglutination (pour antigènes particulaires : bactéries, hématies) :} Grâce à leur valence $\ge$ 2, les anticorps lient entre eux plusieurs antigènes à la surface de cellules ou particules, formant des amas (agglutinats) plus facilement phagocytables et immobilisant les bactéries.
\item \textbf{Opsonisation (facilitation de la phagocytose) :} La région Fc de l'anticorps fixé sur l'antigène (bactérie, cellule tumorale) est reconnue par les \textbf{récepteurs Fc} (Fc$\gamma$R pour IgG) à la surface des phagocytes (macrophages, neutrophiles). Cela déclenche l'engouffrement (phagocytose). Le C3b (fragment du complément) agit aussi comme opsonine.
\item \textbf{Activation de la voie classique du complément :} La fixation de C1q sur la région Fc d'IgG ou d'IgM (en conformation activée) déclenche la cascade du complément $\rightarrow$ lyse cellulaire (complexe d'attaque membranaire MAC), chimiotaxie (C5a), opsonisation (C3b).
\item \textbf{Cytotoxicité cellulaire à médiation anticorps (ADCC) :} Les cellules NK (Natural Killer) expriment le récepteur Fc$\gamma$RIII (CD16). Elles lient la région Fc d'IgG fixée sur une cellule cible (infectée par un virus, tumorale) et la lysent par libération de perforines et granzymes.
\end{enumerate}
\end{propriete}

\begin{popfigure}
\begin{tikzpicture}[scale=0.9, every node/.style={font=\small}]
% Macrophage
\draw[popgreen!80!black, fill=popgreen!15, thick] (0,0) ellipse (2.2cm and 1.3cm);
\node[popgreen!80!black, font=\bfseries\small] at (0, -1.8) {Macrophage};

% Récepteurs Fc à la surface
\foreach \x in {-1.8, -0.9, 0, 0.9, 1.8} {
  \draw[popgreen!80!black, thick] (\x, 1.3) -- (\x, 1.9) node[above, font=\tiny] {Fc$\gamma$R};
}

% Bactéries (Antigènes de surface)
\foreach \i/\x/\y in {1/-1.5/1.4, 2/0/1.6, 3/1.5/1.4, 4/-0.5/0.4, 5/0.5/0.4} {
  \draw[poporange, fill=poporange!30, thick] (\x,\y) circle (0.25cm);
  \node[poporange, font=\tiny] at (\x,\y) {Ag};
}

% Anticorps IgG reliant Ag et recepteur Fc (Opsonisation)
% Exemple 1 (Bactérie gauche)
\draw[popblue, thick] (-1.5, 1.65) -- (-1.5, 2.2); % Fab
\draw[popblue, thick] (-1.5, 2.2) -- (-1.3, 2.5); % Hinge
\draw[popblue, thick] (-1.5, 2.2) -- (-1.7, 2.5);
\draw[popblue, thick] (-1.5, 2.2) -- (-1.5, 2.8); % Fc vers recepteur
\node[popblue, font=\tiny] at (-1.5, 3.0) {IgG};

% Exemple 2 (Bactérie centre)
\draw[popblue, thick] (0, 1.85) -- (0, 2.4);
\draw[popblue, thick] (0, 2.4) -- (-0.2, 2.7);
\draw[popblue, thick] (0, 2.4) -- (0.2, 2.7);
\draw[popblue, thick] (0, 2.4) -- (0, 3.0);
\node[popblue, font=\tiny] at (0, 3.2) {IgG};

% Agglutination / Complexe immun (liaison inter-bactérienne)
\draw[poppurple, thick, decorate, decoration={snake, amplitude=1.5pt, segment length=5pt}] (-1.8, 1.0) -- (-0.8, 1.6) node[above left, font=\tiny, poppurple] {Complexe immun Ag-Ac};

% Flèches phagocytose
\draw[->, thick, popdark] (-1.5, 1.4) -- (-0.7, 0.8) node[midway, above, font=\tiny, rotate=-20] {Phagocytose};
\draw[->, thick, popdark] (0, 1.6) -- (0.2, 0.7);

% Legende interne
\node[draw, fill=white, rounded corners, inner sep=3pt, font=\tiny, anchor=north east, align=center] at (3.8, 2.0){
\textbf{Légende :}\\
\textcolor{poporange}{\textbullet} Antigène (bactérie)\\
\textcolor{popblue}{\textbullet} Anticorps IgG (opsonine)\\
\textcolor{popgreen!80!black}{\textbullet} Récepteur Fc$\gamma$R\\
\textcolor{poppurple}{\textbullet} Complexe immun Ag-Ac
};

\end{tikzpicture}
\legende{Formation de complexes immuns et opsonisation. Les anticorps (IgG) font le pont entre l'antigène à la surface de la bactérie et les récepteurs Fc$\gamma$R du macrophage, déclenchant la phagocytose. L'agglutination (liaison de plusieurs bactéries par un même anticorps ou plusieurs) forme des complexes immuns de grande taille, cibles privilégiées des phagocytes.}
\end{popfigure}

\begin{attention}[Erreur fréquente : « Les anticorps détruisent les antigènes »]
\textbf{Faux.} Les anticorps sont des protéines \emph{non enzymatiques} (sauf les abzymes catalytiques rares). Ils \textbf{ne digèrent pas} et \textbf{ne lysent pas} directement l'antigène ou le pathogène.
\begin{itemize}
\item Ils \textbf{neutralisent} (blocage stérique).
\item Ils \textbf{marquent} (opsonisation) pour les phagocytes.
\item Ils \textbf{recrutent} le système du complément (voie classique) qui peut lyser certaines bactéries à Gram négatif ou cellules cibles.
\end{itemize}
Au Bac, écrire « Les anticorps neutralisent la toxine » ou « Les anticorps opsonisent la bactérie pour la phagocytose », mais \textbf{jamais} « Les anticorps détruisent la bactérie ».
\end{attention}

\courssub{Cinétique de la réponse humorale : réponse primaire et réponse secondaire}

L'étude des taux sériques d'anticorps spécifiques (séroconversion) après une primo-infection (ou primo-vaccination) puis après un second contact avec le même antigène révèle l'existence de la \textbf{mémoire immunitaire}.

\begin{popfigure}
\begin{tikzpicture}
\begin{axis}[
    width=\linewidth,
    height=9cm,
    xlabel={Temps (jours)},
    ylabel={Taux d'anticorps spécifiques (log UI/mL)},
    xmin=0, xmax=60,
    ymin=0, ymax=5.5,
    xtick={0,7,14,21,28,35,42,49,56},
    ytick={0,1,2,3,4,5},
    grid=major,
    grid style={popline},
    legend style={at={(0.02,0.98)}, anchor=north west, fill=white, draw=popdark, font=\small},
    legend cell align=left,
    clip=false,
    axis line style={popdark},
    tick style={popdark}
]
% Réponse Primaire (1ère injection jour 0) - Coordonnées pour forme réaliste
\addplot[popblue, thick, smooth, tension=0.7] coordinates {
(0,0) (3,0.1) (7,0.5) (10,1.5) (14,3.0) (21,2.5) (28,1.5) (35,0.8) (42,0.4) (49,0.2) (56,0.1)
};
\addlegendentry{Réponse primaire (1\textsuperscript{er} contact J0)}

% Réponse Secondaire (2ème injection jour 28)
\addplot[poporange, thick, smooth, tension=0.7] coordinates {
(28,1.5) (29,2.0) (30,3.0) (31,4.5) (33,4.2) (35,3.8) (42,3.0) (49,2.5) (56,2.0)
};
\addlegendentry{Réponse secondaire (2\textsuperscript{e} contact J28)}

% Injection arrows
\draw[->, thick, popdark] (0, -0.4) -- (0, -0.15) node[below, font=\tiny] {1\textsuperscript{re} inj.};
\draw[->, thick, popdark] (28, -0.4) -- (28, -0.15) node[below, font=\tiny] {2\textsuperscript{e} inj. (Rappel)};

% Latence annotations
\draw[<->, dashed, popmuted] (0, 0.3) -- (7, 0.3) node[midway, above, font=\tiny, popmuted] {Latence $\sim$ 7--10 j};
\draw[<->, dashed, popmuted] (28, 1.8) -- (30, 1.8) node[midway, above, font=\tiny, popmuted] {Latence $\sim$ 2--3 j};

% Pic annotations
\node[popblue, font=\small, anchor=south] at (axis cs:14,3.2) {Pic $\sim$ J14};
\node[poporange, font=\small, anchor=south] at (axis cs:31,4.7) {Pic $\sim$ J31};

\end{axis}
\end{tikzpicture}
\legende{Cinétique comparative des taux d'anticorps (échelle logarithmique) après primo-injection (jour 0) et injection de rappel (jour 28). La réponse secondaire est plus précoce (latence courte), plus intense (taux 10 à 100 fois supérieur), plus durable et dominée par des IgG de haute affinité. Unités : UI/mL (Unités Internationales par mL) ou titre sérique.}
\end{popfigure}

\begin{methode}[Interpréter une courbe de séroconversion (réponse primaire vs secondaire)]
Pour analyser un tel graphique au Bac, procédez en 4 étapes :
\begin{enumerate}
\item \textbf{Identifier les phases :} Repérer les deux pics (réponse primaire et secondaire) et les injections (flèches).
\item \textbf{Comparer la latence (délai avant montée) :} Primaire : 7 à 10 jours (activation naïve, expansion clonale, différenciation). Secondaire : 1 à 3 jours (lymphocytes B mémoires déjà différenciés, seuil d'activation plus bas).
\item \textbf{Comparer l'intensité (taux maximal) :} Secondaire 10 à 100 fois plus élevée (plus de plasmocytes, sécrétion plus massive).
\item \textbf{Comparer la durée et la classe d'Ig :} Primaire : IgM dominantes au début, puis IgG, taux chute vite. Secondaire : IgG dominantes dès le début (changement de classe déjà fait), haute affinité (hypermutation somatique), maintien prolongé (plasmocytes à longue vie + mémoires).
\end{enumerate}
\end{methode}

\begin{exemplebox}[Exemple chiffré : Vaccination anti-tétanique (anatoxine tétanique)]
\begin{itemize}
\item \textbf{Primaire (2 injections à 1 mois d'intervalle) :} Taux d'anticorps protecteurs ($> 0,1$ UI/mL) atteints vers J21--J28 après la 2\textsuperscript{e} dose. Pic $\sim$ 1--5 UI/mL.
\item \textbf{Rappel (1 an plus tard) :} Montée en 2--3 jours, pic $\sim$ 50--200 UI/mL (\textbf{$\times$ 50 à $\times$ 100}).
\item \textbf{Conséquence pratique :} Le calendrier vaccinal camerounais (PEV) prévoit 3 doses primaires (Pentavalent à 6, 10, 14 semaines) puis des rappels (DTCq à 15--18 mois, 4--7 ans, 12--15 ans) pour entretenir ce taux protecteur. Une femme enceinte non vaccinée reçoit 2 doses d'anatoxine tétanique (TAT) à 4 semaines d'intervalle pour protéger le nouveau-né par transfert d'IgG transplacentaire.
\end{itemize}
\end{exemplebox}

\begin{exoresolu}[Analyse d'une courbe de séroconversion post-vaccinale]
\textbf{Énoncé :} Un élève de Terminale D reçoit une primo-vaccination contre l'hépatite B (jour 0), une deuxième dose à J30 et une troisième à J180. Le graphique ci-dessous montre l'évolution du taux d'anti-HBs (anticorps anti-antigène de surface HBsAg). On considère qu'un taux $> 10$ mUI/mL est protecteur.
\begin{center}
\begin{tikzpicture}
\begin{axis}[
    width=0.9\linewidth, height=6.5cm,
    xlabel={Jours},
    ylabel={Anti-HBs (mUI/mL)},
    xmin=0, xmax=210,
    ymin=0.1, ymax=2000,
    ymode=log,
    log ticks with fixed point={1,2,3,4,5,6,7,8,9},
    ytick={1,10,100,1000},
    yticklabels={1,10,100,1000},
    grid=major,
    grid style={popline},
    axis line style={popdark},
    tick style={popdark},
    clip=false
]
% Courbe principale
\addplot[popblue, thick, smooth, tension=0.6] coordinates {
(0,0.5) (10,0.6) (20,1) (30,5) (40,50) (60,150) (90,200) (120,150) (150,60) (180,20) (185,100) (190,500) (200,900) (210,800)
};
% Points d'injection
\addplot[only marks, mark=*, mark size=2.5, poporange] coordinates {(0,0.5) (30,5) (180,20)};
\node[poporange, anchor=south, font=\tiny] at (axis cs:0,0.5) {D1};
\node[poporange, anchor=south, font=\tiny] at (axis cs:30,5) {D2};
\node[poporange, anchor=south, font=\tiny] at (axis cs:180,20) {D3};
% Seuil protecteur
\addplot[popred, dashed, thick, domain=0:210] {10};
\node[popred, anchor=west, font=\tiny] at (axis cs:210,10) {Seuil protecteur (10 mUI/mL)};
\end{axis}
\end{tikzpicture}
\end{center}
\textbf{Questions :}
1. Caractériser la réponse après la dose 1 (J0--J30).
2. Expliquer l'effet de la dose 2 (J30).
3. Justifier la nécessité de la dose 3 (J180) au vu de la courbe.
\textbf{Solution détaillée :}
\begin{enumerate}
\item \textbf{Dose 1 (Primaire) :} Latence $\sim$ 3 semaines. Taux maximal $\sim$ 200 mUI/mL (protecteur) vers J90. Chute progressive ensuite (demi-vie des IgG $\sim$ 21--28 jours, pas de plasmocytes à longue vie en nombre suffisant sans rappel). À J180, taux $\sim$ 20 mUI/mL (proche du seuil).
\item \textbf{Dose 2 (Rappel précoce à J30) :} Agit comme un \emph{boost} de la réponse primaire en cours. Montée rapide (latence $\sim$ 1 semaine), pic plus haut ($\sim$ 200 mUI/mL plus tôt). C'est encore une réponse primaire amplifiée (les lymphocytes mémoires ne sont pas encore pleinement générés/sélectionnés).
\item \textbf{Dose 3 (Rappel tardif à J180) :} À J180, le taux a chuté (perte des plasmocytes de la réponse primaire). L'injection rencontre une \textbf{population de lymphocytes B mémoires} installée. Réponse secondaire typique : latence $\sim$ 3--5 jours (J185--J188), pic très haut ($\sim$ 800--1000 mUI/mL), IgG de haute affinité. Cette dose 3 assure la \textbf{mémoire à long terme} (protection $> 20$ ans chez 95\% des vaccinés).
\end{enumerate}
\end{exoresolu}

\courssub{Mémoire immunitaire et fondement de la vaccination}

\begin{aretenir}[Base cellulaire et moléculaire de la mémoire immunitaire humorale]
La mémoire immunitaire repose sur deux piliers générés dans les \textbf{centres germinatifs} des ganglions et de la rate lors de la réponse primaire :
\begin{enumerate}
\item \textbf{Lymphocytes B mémoires :} Cellules quiescentes (G0), longue durée de vie, BCR de haute affinité (sélection par affinité), seuil d'activation bas. Se différencient rapidement en plasmocytes au second contact sans besoin de forte coopération T (ou coopération plus facile).
\item \textbf{Plasmocytes à longue vie (Long-lived plasma cells) :} Nichés dans la moelle osseuse (survie dépendante de l'IL-6, APRIL, contact avec cellules stromales). Sécrètent constitutivement des anticorps (IgG surtout) pendant des années $\rightarrow$ \textbf{immunité sérique persistante} (taux d'anticorps résiduels protecteurs sans rappel).
\end{enumerate}
La vaccination (anatoxines, vaccins sous-unitaires, conjugués, ARNm, vecteurs viraux) mime l'infection naturelle pour générer cette mémoire \textbf{sans la maladie}.
\end{aretenir}

\begin{savaistu}[Le saviez-vous ? -- Behring et le premier Nobel]
En 1901, Emil von Behring reçoit le \textbf{tout premier prix Nobel de physiologie ou médecine} « pour ses travaux sur la sérothérapie, notamment son application contre la diphtérie, par lesquels il a ouvert une nouvelle route dans le domaine de la médecine et mis entre les mains du médecin une arme victorieuse contre la maladie et la mort ». Son sérum antitoxique diphtérique, produit chez le cheval, a fait chuter la mortalité de la diphtérie de 50\% à moins de 15\% en quelques années en Europe. Au Cameroun, bien que la diphtérie soit rare grâce au PEV (vaccin DTCq/Hib/Hépatite B), la sérothérapie antivenineuse reste vitale dans les régions forestières et soudano-sahéliennes où les morsures de serpents (vipères, cobras, mambas) causent encore des décès par manque d'accès rapide aux sérums polyvalents africains (ex: SAIMR, Inoserp).
\end{savaistu}

\begin{exercice}[Entraînement Bac -- Réponse humorale et vaccination]
Un nourrisson de 3 mois reçoit sa première dose de vaccin pentavalent (DTCq-Hib-Hépatite B). Son taux d'anti-HBs est dosé à J0, J30, J60, J90, J180, J365.
\begin{enumerate}
\item Prédire l'allure de la courbe (taux vs temps) en distinguant les réponses primaires et secondaires.
\item Expliquer pourquoi le calendrier prévoit une dose à J30 (1 mois) et une à J180 (6 mois) plutôt qu'une seule dose.
\item Si ce nourrisson est exposé au VHB à J200, décrire la réponse immunitaire attendue (cellules impliquées, cinétique, classe d'anticorps).
\end{enumerate}
\end{exercice}

\begin{corrige}[Exercice n°1 -- Réponse humorale et vaccination]
\begin{enumerate}
\item \textbf{Courbe :} J0--J30 : taux basal (anticorps maternels résiduels IgG s'ils existent, sinon 0). J30--J90 : réponse primaire lente (latence 2--3 sem), pic IgM puis IgG modéré ($\sim$ 10--100 mUI/mL). J90--J180 : déclin exponentiel. J180 (Dose 3) : réponse secondaire brutale (latence 3--5 j), pic IgG haut ($> 1000$ mUI/mL). J180--J365 : plateau élevé (plasmocytes moelle osseuse).
\item \textbf{Justification doses rapprochées puis espacées :} Les 2 premières doses (J0, J30) « amorcent » la réponse primaire et génèrent les premiers lymphocytes B mémoires. La 3\textsuperscript{e} dose à J180 (6 mois) rencontre une population mémoires mature $\rightarrow$ réponse secondaire robuste, installation de plasmocytes à longue vie dans la moelle $\rightarrow$ protection décennale. Une dose unique ne génère pas de mémoire durable (tolérance ou réponse primaire avortée).
\item \textbf{Exposition à J200 :} Présence d'anticorps circulants protecteurs (IgG anti-HBs $> 10$ mUI/mL) $\rightarrow$ \textbf{neutralisation immédiate} du virion (immunité stérilisante). Si quelques virions échappent : activation rapide des lymphocytes B mémoires spécifiques $\rightarrow$ différenciation en plasmocytes en 2--3 jours $\rightarrow$ montée massive d'IgG haute affinité (réponse anamnestique). Pas de maladie clinique.
\end{enumerate}
\end{corrige}

\coursec{La réponse immunitaire spécifique à médiation cellulaire}

Dans la section précédente, nous avons vu que les anticorps sécrétés par les plasmocytes neutralisent les antigènes présents dans les liquides de l'organisme : toxines, bactéries libres, virus en circulation. Cette réponse à médiation \emph{humorale} est redoutable... mais elle a une limite fondamentale : un anticorps est une molécule qui ne peut pas pénétrer dans une cellule. Or certains agents pathogènes, comme le bacille de Koch (responsable de la tuberculose, encore très présente dans nos centres de santé) ou les virus, vivent \emph{à l'intérieur} des cellules de l'hôte. Cachés derrière la membrane plasmique, ils sont inaccessibles aux anticorps. Comment l'organisme les élimine-t-il alors ? C'est toute la question de cette section : il existe une seconde réponse spécifique, assurée non par des molécules mais par des \emph{cellules tueuses} : la réponse à médiation cellulaire.

\courssub{Mise en évidence expérimentale : l'immunité contre la tuberculose est cellulaire}

\textbf{Observation et question.} Un animal (cobaye) infecté par le bacille de Koch puis guéri est \emph{immunisé} contre la tuberculose : une seconde injection de bacilles ne le rend pas malade. Quel est le support de cette immunité : les anticorps de son sérum, ou ses cellules ?

\begin{experience}[L'expérience de Chase (1945) : transfert de sérum contre transfert de lymphocytes]
\textbf{Protocole.} On immunise un cobaye donneur contre le bacille de Koch. On prélève ensuite, séparément, son \textbf{sérum} (liquide sanguin contenant les anticorps) et ses \textbf{lymphocytes}. On transfère chacun de ces éléments à deux lots de cobayes receveurs sains, puis on inocule à tous les animaux une dose de bacilles de Koch virulents.

\begin{center}
\begin{tabularx}{\linewidth}{|l|X|X|X|}
\hline
\rowcolor{popblueL} \textbf{Lot} & \textbf{Élément transféré} & \textbf{Inoculation de BK} & \textbf{Résultat} \\
\hline
A & Sérum du cobaye immunisé (anticorps) & Oui & Les cobayes développent la tuberculose \\
\hline
B & Lymphocytes du cobaye immunisé & Oui & Les cobayes restent sains (protégés) \\
\hline
C (témoin) & Rien & Oui & Les cobayes développent la tuberculose \\
\hline
\end{tabularx}
\end{center}

\textbf{Interprétation.} Le sérum du donneur immunisé, bien qu'il contienne des anticorps, ne protège pas : les anticorps sont donc \emph{inefficaces} contre le bacille de Koch. En revanche, les lymphocytes transférés confèrent la protection : ce sont eux qui portent l'immunité antituberculeuse.

\textbf{Conclusion.} L'immunité contre la tuberculose est une immunité à \cle{médiation cellulaire} : elle est assurée par des cellules (lymphocytes) et non par des anticorps.
\end{experience}

\begin{methode}[Interpréter une expérience de transfert (sérum vs lymphocytes)]
Face à un protocole de transfert, procède toujours ainsi :
\begin{enumerate}
\item \textbf{Identifie la variable testée} : c'est la \emph{nature de l'élément transféré} (sérum = anticorps ; lymphocytes = cellules). Toutes les autres conditions doivent être identiques entre les lots.
\item \textbf{Repère le lot témoin} (aucun transfert) : il montre ce qui se passe sans protection et valide l'expérience.
\item \textbf{Compare les résultats} : si seul le transfert de sérum protège, l'immunité est \emph{humorale} ; si seul le transfert de lymphocytes protège, elle est \emph{cellulaire}.
\item \textbf{Conclus en citant les faits} : « le lot ayant reçu les lymphocytes survit, donc... », jamais l'inverse (on ne part pas de la conclusion pour expliquer les résultats).
\end{enumerate}
\end{methode}

\textbf{Pourquoi les anticorps échouent-ils ici ?} Le bacille de Koch est un agent \emph{intracellulaire} : phagocyté par les macrophages, il survit et se multiplie à l'intérieur de ces cellules. Les anticorps, grosses molécules protéiques, ne traversent pas la membrane plasmique : l'agent pathogène leur est littéralement invisible. La seule stratégie possible est de \emph{détruire la cellule infectée elle-même}, avec l'agent qu'elle abrite. C'est le rôle des lymphocytes T cytotoxiques.

\begin{definition}[Réponse immunitaire à médiation cellulaire]
La \cle{réponse à médiation cellulaire} est la réponse immunitaire spécifique assurée par les \cle{lymphocytes T cytotoxiques} (LT8), qui détruisent \emph{par contact direct} les cellules de l'organisme infectées par un agent intracellulaire (virus, bactérie intracellulaire) ou devenues anormales (cellules cancéreuses, cellules greffées).
\end{definition}

\courssub{Les lymphocytes T cytotoxiques (LT8) : reconnaissance de la cellule infectée}

\textbf{Le problème de la reconnaissance.} Pour tuer une cellule infectée, le lymphocyte doit d'abord la \emph{reconnaître}. Mais l'agent pathogène est caché à l'intérieur... Comment le LT8 « sait-il » que la cellule est infectée ? La réponse est élégante : toute cellule infectée expose à sa surface des \emph{fragments de l'agent} qu'elle abrite, comme un drapeau de détresse.

\begin{propriete}[Présentation de l'antigène par le CMH de classe I (admise)]
Toute cellule nucléée de l'organisme porte à sa surface des molécules du \cle{CMH de classe I} (complexe majeur d'histocompatibilité, appelé HLA chez l'Homme). Ces molécules présentent en permanence des fragments (peptides) des protéines fabriquées dans la cellule :
\begin{itemize}
\item chez une cellule saine, le CMH I présente des peptides du \emph{soi} : aucun LT8 ne réagit ;
\item chez une cellule infectée, le CMH I présente des peptides de l'\emph{agent pathogène} (non-soi) : le complexe \textbf{CMH I + antigène} est reconnu par le récepteur spécifique (TCR) d'un LT8, aidé de sa molécule \cle{CD8}.
\end{itemize}
\end{propriete}

Le LT8 possède un récepteur (TCR) spécifique d'un seul antigène : il ne s'active que si son TCR reconnaît le couple CMH I--antigène. C'est la \emph{double reconnaissance} : le LT8 vérifie à la fois l'identité de la cellule (CMH du soi) et la présence de l'intrus (peptide étranger). La molécule CD8, caractéristique des LT cytotoxiques (d'où leur nom de « LT8 » ou « T CD8$^+$ »), stabilise cette liaison.

\courssub{Le mécanisme de destruction : perforines et granzymes}

Une fois fixé sur sa cible, le LT8 déclenche une exécution rapide et propre, en trois étapes :

\begin{enumerate}
\item \textbf{Contact et adhésion} : le TCR du LT8 se fixe sur le complexe CMH I--antigène de la cellule infectée ; la zone de contact se resserre (synapse immunologique).
\item \textbf{Exocytose des perforines} : le LT8 libère par exocytose des granules contenant des \cle{perforines}. Ces protéines s'assemblent dans la membrane de la cellule cible et y forment des \emph{canaux} (pores), comme des poinçons.
\item \textbf{Entrée des granzymes et apoptose} : par ces pores pénètrent les \cle{granzymes}, enzymes qui activent dans la cellule cible une cascade de destruction programmée : l'\cle{apoptose}. La cellule infectée se fragmente en petits corps qui seront phagocytés par les macrophages, détruisant du même coup les agents pathogènes qu'elle contenait.
\end{enumerate}

\begin{attention}[Apoptose, pas nécrose]
La mort induite par le LT8 est une \emph{apoptose} (mort cellulaire programmée, « propre », sans libération brutale du contenu cellulaire), et non une nécrose (éclatement inflammatoire). L'apoptose limite la dissémination des virus et l'inflammation des tissus voisins. Ne rédige jamais « le LT8 fait exploser la cellule » : dis « le LT8 induit l'apoptose de la cellule cible ».
\end{attention}

\begin{popfigure}
\begin{tikzpicture}[font=\small, >=Stealth]
% Cellule cible (infectée)
\draw[fill=poppinkL, draw=poppink, thick] (4.5,0) ellipse (2.6 and 1.9);
\node[font=\bfseries\small, popdark] at (4.5,1.1) {Cellule infectée};
\node[font=\scriptsize, popdark] at (4.5,0.6) {(CMH I + antigène viral)};
% virus inside
\foreach \p in {(3.6,-0.5),(4.6,-0.9),(5.5,-0.4),(4.2,0.1)}{
  \node[popink] at \p {\scriptsize $\ast$};
}
\node[font=\scriptsize, popink] at (5.9,-1.3) {virus intracellulaires};
% CMH I complexes on top membrane
\foreach \x in {3.3,4.5,5.7}{
  \draw[popdark, thick] (\x,1.72) -- (\x,2.15);
  \draw[popdark, thick] (\x-0.12,1.72) -- (\x,2.05) -- (\x+0.12,1.72);
  \fill[poporange] (\x,2.2) circle (0.09);
}
\node[font=\scriptsize, popdark, anchor=west] at (6.6,2.1) {CMH I};
\node[font=\scriptsize, poporange, anchor=west] at (6.6,1.75) {peptide antigénique};
% LT8
\draw[fill=popblueL, draw=popblue, thick] (4.5,4.3) ellipse (2.0 and 1.15);
\node[font=\bfseries\small, popblue] at (4.5,4.75) {LT8 (lymphocyte T cytotoxique)};
% granules
\foreach \p in {(3.6,4.0),(4.2,3.85),(4.9,3.95),(5.4,4.05)}{
  \draw[fill=popgoldL, draw=popgold] \p circle (0.13);
}
\node[font=\scriptsize, popdark] at (4.5,4.25) {granules (perforines + granzymes)};
% TCR
\foreach \x in {3.3,4.5,5.7}{
  \draw[popblue, thick] (\x,3.18) -- (\x,2.62);
  \draw[popblue, thick] (\x-0.12,2.62) -- (\x+0.12,2.62);
}
\node[font=\scriptsize, popblue, anchor=east] at (2.9,2.9) {TCR + CD8};
% exocytosis arrow
\draw[->, popgreen, very thick] (1.6,3.6) -- (1.6,1.4) node[midway, left, font=\scriptsize, align=center]{exocytose\\ des granules};
% perforin pores
\foreach \x in {3.3,4.5,5.7}{
  \draw[popgreen, very thick] (\x-0.1,2.45) -- (\x-0.1,1.95);
  \draw[popgreen, very thick] (\x+0.1,2.45) -- (\x+0.1,1.95);
}
\node[font=\scriptsize, popgreen, anchor=west] at (6.6,2.5) {pores de perforines};
% apoptosis
\draw[->, popink, very thick] (7.3,-0.3) -- (9.3,-0.3);
\node[font=\scriptsize\bfseries, popink] at (8.3,0.15) {apoptose};
\foreach \p in {(9.9,0.2),(10.4,-0.3),(9.8,-0.7),(10.6,0.4)}{
  \draw[fill=poppinkL, draw=poppink] \p circle (0.22);
}
\node[font=\scriptsize, popdark] at (10.2,-1.2) {corps apoptotiques (phagocytés)};
\end{tikzpicture}
\legende{Mécanisme de destruction d'une cellule infectée par un LT8. 1 : le TCR (associé à CD8) reconnaît le complexe CMH I + peptide antigénique à la surface de la cellule infectée ; 2 : le LT8 libère par exocytose des perforines qui percent la membrane cible ; 3 : les granzymes pénètrent et déclenchent l'apoptose ; les corps apoptotiques sont éliminés par les macrophages.}
\end{popfigure}

\courssub{Le rôle coordinateur des lymphocytes T auxiliaires (LT4)}

Les LT8 ne travaillent pas seuls. Leur activation, comme celle des lymphocytes B de la réponse humorale, exige l'intervention d'un chef d'orchestre : le \cle{lymphocyte T auxiliaire} (LT4, ou T CD4$^+$).

\begin{propriete}[Les LT4, chefs d'orchestre de la réponse immunitaire spécifique (admise)]
Les LT4 reconnaissent les antigènes présentés par le \cle{CMH de classe II} des cellules présentatrices d'antigène (macrophages, cellules dendritiques). Activés, ils sécrètent des \cle{cytokines} (interleukines, notamment l'\emph{interleukine 2}, IL-2) qui :
\begin{itemize}
\item stimulent la multiplication (amplification clonale) et la différenciation des \textbf{LT8} en cellules tueuses (réponse cellulaire) ;
\item stimulent la multiplication et la différenciation des \textbf{lymphocytes B} en plasmocytes sécréteurs d'anticorps (réponse humorale) ;
\item activent les \textbf{macrophages}.
\end{itemize}
Les LT4 coordonnent donc les \emph{deux} types de réponse spécifique : leur destruction compromet toute l'immunité adaptative. C'est précisément ce que réalise le VIH, qui infecte et détruit les LT4 (voir section 6).
\end{propriete}

\begin{attention}[Ne pas confondre LT8 et LT4]
Erreur classique du Baccalauréat :
\begin{itemize}
\item \textbf{LT8} (CD8$^+$) = lymphocytes T \emph{cytotoxiques} : ils \textbf{tuent} les cellules infectées (perforines, granzymes, apoptose). Ils reconnaissent le CMH de \emph{classe I}.
\item \textbf{LT4} (CD4$^+$) = lymphocytes T \emph{auxiliaires} : ils ne tuent pas, ils \textbf{stimulent} les autres cellules par leurs cytokines. Ils reconnaissent le CMH de \emph{classe II}.
\end{itemize}
Et retiens bien : le \textbf{VIH cible les LT4}, pas les LT8. C'est pourquoi le Sida effondre l'ensemble des défenses immunitaires.
\end{attention}

\courssub{La mémoire immunitaire cellulaire}

Comme la réponse humorale, la réponse cellulaire laisse une trace durable. Lors de l'amplification clonale, une partie des LT8 produits ne devient pas effectrice mais se différencie en \cle{lymphocytes T mémoires}. Ces cellules, à vie longue, persistent dans l'organisme après l'élimination de l'agent pathogène. Lors d'un second contact avec le même antigène, elles déclenchent une réponse cellulaire \emph{plus rapide et plus intense} : c'est la mémoire immunitaire cellulaire, fondement de l'immunité durable contre les infections intracellulaires.

\begin{exemplebox}[L'IDR à la tuberculine : un test de mémoire cellulaire]
L'\emph{intradermo-réaction} (IDR) à la tuberculine, ou cuti-réaction, est utilisée dans les centres de santé camerounais pour dépister un contact ancien avec le bacille de Koch. On injecte dans le derme de la tuberculine (extraits antigéniques du bacille, incapables de provoquer la maladie). Chez une personne ayant déjà rencontré le bacille (infection ancienne ou BCG), les \textbf{LT4 mémoires (Th1)} réagissent : en \SI{48}{\hour} à \SI{72}{\hour} apparaît une \emph{induration} (durcissement) rouge au point d'injection, due à l'afflux de lymphocytes T et de macrophages (réaction d'hypersensibilité retardée de type IV). Chez une personne jamais en contact, aucune réaction. Ce test révèle donc une \emph{mémoire immunitaire cellulaire} : il ne détecte pas des anticorps, mais l'existence de LT mémoires spécifiques.
\end{exemplebox}

\courssub{Complémentarité des réponses humorale et cellulaire}

Les deux réponses spécifiques ne sont pas concurrentes mais \emph{complémentaires} : chacune agit là où l'autre est impuissante, et toutes deux sont coordonnées par les LT4.

\begin{popfigure}
\begin{tikzpicture}[font=\small]
\draw[fill=popblueL, draw=popblue, thick, rounded corners] (0,0) rectangle (6.4,6.2);
\draw[fill=popgreenL, draw=popgreen, thick, rounded corners] (7.0,0) rectangle (13.4,6.2);
\node[font=\bfseries, popblue] at (3.2,5.7) {Réponse à médiation humorale};
\node[font=\bfseries, popgreen] at (10.2,5.7) {Réponse à médiation cellulaire};
% left column
\node[font=\scriptsize, text=popdark, align=left, anchor=north west] at (0.25,5.15) {%
\textbf{Effecteurs :} lymphocytes B\\
\phantom{\textbf{Effecteurs :} }$\to$ plasmocytes $\to$ anticorps\\[2pt]
\textbf{Cibles :} antigènes \emph{extracellulaires}\\
\phantom{\textbf{Cibles :} }(toxines, bactéries et virus\\
\phantom{\textbf{Cibles :} }libres dans les humeurs)\\[2pt]
\textbf{Mécanismes :} neutralisation,\\
\phantom{\textbf{Mécanismes :} }agglutination, complexe immun,\\
\phantom{\textbf{Mécanismes :} }phagocytose facilitée\\[2pt]
\textbf{Mémoire :} lymphocytes B mémoires};
% right column
\node[font=\scriptsize, text=popdark, align=left, anchor=north west] at (7.25,5.15) {%
\textbf{Effecteurs :} lymphocytes T cytotoxiques\\
\phantom{\textbf{Effecteurs :} }(LT8, CD8$^+$)\\[2pt]
\textbf{Cibles :} cellules \emph{infectées} (agents\\
\phantom{\textbf{Cibles :} }intracellulaires), cellules cancéreuses,\\
\phantom{\textbf{Cibles :} }cellules greffées (rejet de greffe)\\[2pt]
\textbf{Mécanismes :} contact CMH I--antigène,\\
\phantom{\textbf{Mécanismes :} }perforines + granzymes,\\
\phantom{\textbf{Mécanismes :} }apoptose de la cellule cible\\[2pt]
\textbf{Mémoire :} lymphocytes T mémoires};
% central coordination
\draw[fill=popgoldL, draw=popgold, thick, rounded corners] (4.1,-1.6) rectangle (9.3,-0.5);
\node[font=\scriptsize\bfseries, popdark] at (6.7,-1.05) {Coordination par les LT4 (cytokines, IL-2)};
\draw[->, popgold, thick] (5.2,-0.5) -- (3.6,0);
\draw[->, popgold, thick] (8.2,-0.5) -- (9.8,0);
\end{tikzpicture}
\legende{Comparaison des deux réponses immunitaires spécifiques. Les anticorps n'agissent que dans les liquides (humeurs) ; les LT8 détruisent les cellules infectées ou anormales par contact. Les LT4 coordonnent les deux réponses.}
\end{popfigure}

\begin{aretenir}[L'essentiel de la section]
\begin{itemize}
\item L'expérience de Chase (1945) établit que l'immunité antituberculeuse est \emph{cellulaire} : le transfert de lymphocytes protège, celui de sérum non.
\item Les anticorps sont inefficaces contre les agents \emph{intracellulaires}, cachés dans les cellules.
\item Les \textbf{LT8} reconnaissent le complexe \textbf{CMH I + antigène} à la surface des cellules infectées et les détruisent par \textbf{perforines} (pores membranaires) et \textbf{granzymes} (déclenchement de l'apoptose).
\item Les \textbf{LT4} sécrètent des cytokines qui coordonnent les réponses humorale et cellulaire ; le VIH les détruit.
\item La réponse cellulaire produit des \textbf{LT mémoires} (révélés par l'IDR à la tuberculine) ; elle intervient aussi dans le \emph{rejet de greffe} et l'élimination des \emph{cellules cancéreuses}.
\end{itemize}
\end{aretenir}

\begin{exoresolu}[Interpréter une expérience de transfert]
\textbf{Énoncé.} Des souris A sont immunisées contre un virus qui infecte les cellules du foie. On prélève leur sérum et leurs lymphocytes. Des souris B (vierges de tout contact avec ce virus) reçoivent le sérum des souris A ; des souris C reçoivent leurs lymphocytes ; des souris D ne reçoivent rien. Toutes sont ensuite infectées par le virus. Résultats : les souris B et D développent une hépatite virale grave ; les souris C restent saines. (1) Quel est le rôle des souris D ? (2) Que peut-on conclure sur le type d'immunité anti-virale mis en jeu ? Justifier.
\tcblower
\textbf{Solution.} (1) Les souris D constituent le \emph{lot témoin} : elles montrent que, sans transfert, l'infection par le virus provoque la maladie ; elles valident la comparaison. (2) La variable testée est la nature de l'élément transféré. Le transfert de \emph{sérum} (souris B) ne protège pas : les anticorps sont inefficaces, ce qui s'explique par le caractère intracellulaire du virus, inaccessible aux anticorps une fois dans les hépatocytes. Le transfert de \emph{lymphocytes} (souris C) protège : l'immunité repose donc sur des cellules. \textbf{Conclusion :} l'immunité contre ce virus est une immunité à \emph{médiation cellulaire}, assurée par les lymphocytes T cytotoxiques qui détruisent les hépatocytes infectés par apoptose.
\end{exoresolu}

\begin{exercice}[À toi de jouer]
Un patient reçoit une greffe de rein. Quelques semaines plus tard, le rein greffé est détruit alors que le patient ne présente aucun signe d'infection. (1) Quel type de réponse immunitaire est responsable du rejet ? (2) Quelles cellules effectrices interviennent et quel signal reconnaissent-elles à la surface des cellules greffées ? (3) Pourquoi les médecins prescrivent-ils des immunosuppresseurs aux patients greffés ?
\end{exercice}

\begin{corrige}[Exercice : rejet de greffe]
(1) Il s'agit d'une réponse immunitaire spécifique \emph{à médiation cellulaire}. (2) Les effecteurs sont les lymphocytes T cytotoxiques (LT8) du receveur : ils reconnaissent les molécules du CMH de classe I du donneur, différentes des siennes, comme du non-soi ; les cellules du greffon sont alors détruites par apoptose (perforines et granzymes), sous la coordination des LT4. (3) Les immunosuppresseurs inhibent l'activation et la multiplication des lymphocytes T : ils empêchent le rejet du greffon, au prix d'une diminution des défenses immunitaires (risque accru d'infections, à surveiller).
\end{corrige}

\coursec{Les maladies auto-immunes}

Dans la section précédente, nous avons vu comment les lymphocytes T cytotoxiques (LT8) détruisent avec une redoutable précision les cellules de l'organisme infectées par un virus : la cellule cible est reconnue grâce à un fragment d'antigène \emph{étranger} présenté par son CMH de classe I. Mais que se passerait-il si cette arme de destruction massive se retournait contre des cellules \emph{saines} de l'organisme ? Si des lymphocytes ou des anticorps prenaient pour cible un constituant du soi ? C'est exactement ce qui se produit dans les \cle{maladies auto-immunes}, que nous étudions maintenant.

\courssub{La tolérance du soi et sa rupture}

\textbf{Rappel : comment le système immunitaire distingue normalement le soi du non-soi.}

Au cours de leur maturation (dans le thymus pour les lymphocytes T, dans la moelle osseuse pour les lymphocytes B), les lymphocytes qui reconnaissent fortement les antigènes du soi sont normalement \emph{éliminés} ou \emph{inactivés}. Ce processus, appelé \cle{sélection clonale négative}, établit la \cle{tolérance du soi} : l'organisme ne réagit pas contre ses propres constituants.

\begin{definition}[Tolérance du soi]
La \cle{tolérance du soi} est l'absence de réponse immunitaire de l'organisme contre ses propres constituants (auto-antigènes). Elle est acquise lors de la maturation des lymphocytes, par élimination (délétion) ou inactivation (anergie) des clones lymphocytaires dirigés contre le soi.
\end{definition}

Il arrive cependant que certains clones \og interdits \fg{} échappent à cette sélection, ou soient activés plus tard (à la suite d'une infection, d'un stress, ou par ressemblance entre un antigène microbien et un antigène du soi). La tolérance est alors \emph{rompue}.

\begin{definition}[Maladie auto-immune]
Une \cle{maladie auto-immune} est une maladie due à une réaction immunitaire de l'organisme contre ses propres constituants. Elle résulte d'une \cle{rupture de la tolérance du soi} : le système immunitaire produit des \cle{auto-anticorps} et/ou des \cle{lymphocytes T auto-réactifs} dirigés contre un organe cible, qu'ils détruisent ou dérèglent.
\end{definition}

\begin{definition}[Auto-anticorps]
Un \cle{auto-anticorps} est un anticorps dirigé contre un antigène du soi (auto-antigène). Sa présence dans le sang est un signe diagnostique majeur des maladies auto-immunes.
\end{definition}

\textbf{Mécanisme général.} Le schéma causal est toujours le même :
\begin{enumerate}
\item un \cle{auto-antigène} (constituant normal de l'organisme) est reconnu comme étranger ;
\item une réponse immunitaire \emph{spécifique} se déclenche : production d'auto-anticorps (réponse humorale) et/ou activation de LT8 auto-réactifs (réponse cellulaire) ;
\item l'organe cible est détruit ou dérèglé, ce qui produit les symptômes de la maladie.
\end{enumerate}

\begin{methode}[Décrire un cas de dysfonctionnement du système immunitaire]
Face à un document décrivant une maladie auto-immune (exercice type Bac), rédige ta réponse en quatre étapes :
\begin{enumerate}
\item \textbf{Identifier l'organe cible} : quel tissu ou quelle cellule est attaqué ? (ex. : cellules $\beta$ du pancréas).
\item \textbf{Identifier l'effecteur immunitaire fautif} : auto-anticorps, lymphocytes T cytotoxiques, ou les deux ?
\item \textbf{Décrire le symptôme} observé (ex. : hyperglycémie chronique).
\item \textbf{Établir le lien de cause à effet} : rupture de la tolérance $\Rightarrow$ attaque de l'organe $\Rightarrow$ perte ou dérèglement de sa fonction $\Rightarrow$ symptôme.
\end{enumerate}
Conclus toujours en rappelant qu'il s'agit d'une rupture de la tolérance du soi.
\end{methode}

\courssub{Le vitiligo : destruction des mélanocytes}

\textbf{Observation.} Dans nos villes et villages, on rencontre des personnes dont la peau présente des plaques blanches bien limitées, souvent au visage, aux mains ou autour de la bouche. Ces plaques ne sont ni contagieuses, ni dues à une lèpre : il s'agit le plus souvent d'un \cle{vitiligo}.

\textbf{Interprétation.} La couleur de la peau est due à la \cle{mélanine}, pigment synthétisé par des cellules de l'épiderme appelées \cle{mélanocytes}. Chez les patients atteints de vitiligo, on observe :
\begin{itemize}
\item la disparition des mélanocytes dans les plaques dépigmentées ;
\item la présence de lymphocytes T cytotoxiques et d'auto-anticorps dirigés contre des protéines des mélanocytes.
\end{itemize}

\textbf{Conclusion.} Le vitiligo est une maladie auto-immune dans laquelle des lymphocytes T auto-réactifs \emph{détruisent les mélanocytes} ; privée de mélanine, la peau se dépigmente en plaques. La fonction barrière de la peau n'est pas abolie, mais les plaques, dépourvues de protection contre les ultraviolets, sont très sensibles au soleil — un point important d'éducation à la santé sous notre climat.

\courssub{Le diabète juvénile (diabète de type 1) : destruction des cellules $\beta$}

\textbf{Observation et démarche.} Un adolescent se met à boire énormément, à uriner fréquemment et à maigrir malgré un appétit conservé. Au centre de santé, la glycémie à jeun mesurée est de \SI{3,2}{g/L} (normale : \SIrange{0,7}{1,1}{g/L}). Des analyses montrent la présence dans son sang d'\emph{auto-anticorps anti-cellules des îlots de Langerhans} et l'absence quasi totale d'insuline. À l'autopsie de pancréas de patients décédés de cette maladie, les îlots de Langerhans apparaissent infiltrés de lymphocytes (on parle d'\emph{insulite}) et les cellules $\beta$ sont détruites.

\textbf{Interprétation.} Les \cle{cellules $\beta$} des îlots de Langerhans (pancréas endocrine) sont les seules cellules de l'organisme capables de sécréter l'\cle{insuline}, hormone hypoglycémiante. Chez ces patients :
\begin{itemize}
\item des \cle{lymphocytes T cytotoxiques} auto-réactifs reconnaissent des auto-antigènes des cellules $\beta$ et les détruisent par perforines et granzymes ;
\item des \cle{auto-anticorps} dirigés contre les cellules $\beta$ (anti-insuline, anti-GAD) témoignent de la réaction auto-immune.
\end{itemize}

\textbf{Conclusion.} La destruction des cellules $\beta$ entraîne une \emph{absence totale d'insuline} : le glucose ne peut plus entrer dans les cellules et s'accumule dans le sang, d'où une \cle{hyperglycémie chronique}. C'est le \cle{diabète de type 1} (ou diabète juvénile, ou diabète insulinodépendant). Le traitement est \emph{substitutif} : injections quotidiennes d'insuline à vie.

\begin{popfigure}
\begin{center}
\begin{tikzpicture}[scale=0.95, every node/.style={font=\small}]
% Cellule beta
\fill[popblueL] (0,0) ellipse (2.6 and 1.9);
\draw[popblue, thick] (0,0) ellipse (2.6 and 1.9);
\fill[popblue!60] (0,0.3) circle (0.7);
\draw[popblue] (0,0.3) circle (0.7);
\node at (0,0.3) {noyau};
\node[popdark] at (0,-1.15) {cytoplasme (granules d'insuline)};
\foreach \a in {40,80,120,160,200,240,300,340}{
  \fill[popgold] ({1.7*cos(\a)},{1.15*sin(\a)}) circle (0.07);}
% CMH + auto-antigène
\draw[popgreen, very thick] (-1.9,1.55) -- (-1.9,2.15);
\draw[popgreen, very thick] (-2.15,2.15) -- (-1.65,2.15);
\fill[poporange] (-1.9,2.35) circle (0.12);
\node[popgreen, anchor=south west] at (-1.55,2.05) {\footnotesize CMH I + auto-antigène};
% LT8
\fill[poppinkL] (-4.6,3.2) circle (1.05);
\draw[poppink, thick] (-4.6,3.2) circle (1.05);
\fill[poppink!70] (-4.6,3.2) circle (0.45);
\node at (-4.6,3.2) {\footnotesize noyau};
\node[popdark] at (-4.6,1.85) {\textbf{LT8 auto-réactif}};
% TCR
\draw[poppink, very thick] (-3.7,2.75) -- (-3.15,2.45);
\draw[poppink, very thick] (-3.25,2.6) -- (-3.05,2.3);
\node[poppink, anchor=west] at (-3.6,3.6) {\footnotesize TCR spécifique};
% Liaison reconnaissance
\draw[popdark, dashed, thick] (-3.1,2.42) -- (-2.05,2.35);
% Granules perforine
\foreach \p in {(-3.9,3.7),(-4.2,3.9),(-3.75,3.35)}{\fill[poporange] \p circle (0.09);}
\draw[-{Stealth}, poporange, very thick] (-3.6,3.3) .. controls (-2.8,3.4) .. (-2.1,2.5);
\node[poporange, anchor=west] at (-2.9,3.75) {\footnotesize perforines + granzymes};
% Conséquence
\draw[-{Stealth}, popdark, very thick] (2.8,0) -- (4.1,0);
\node[popdark, anchor=west] at (4.2,0.55) {\footnotesize apoptose de la cellule $\beta$};
\node[popdark, anchor=west] at (4.2,-0.05) {\footnotesize $\Rightarrow$ arrêt de la sécrétion};
\node[popdark, anchor=west] at (4.2,-0.6) {\footnotesize d'insuline $\Rightarrow$ hyperglycémie};
\node[popblue, anchor=north] at (0,-2.15) {\textbf{Cellule $\beta$ d'un îlot de Langerhans}};
\end{tikzpicture}
\end{center}
\legende{Mécanisme du diabète juvénile (type 1) : un lymphocyte T cytotoxique auto-réactif reconnaît, via son TCR, un auto-antigène présenté par le CMH de classe I de la cellule $\beta$ pancréatique ; il libère des perforines et des granzymes qui déclenchent l'apoptose de la cellule. La destruction massive des cellules $\beta$ supprime la sécrétion d'insuline.}
\end{popfigure}

\begin{attention}[Ne confonds pas diabète de type 1 et diabète de type 2]
Erreur fréquente au Bac : le \cle{diabète juvénile} (type 1) est une maladie \emph{auto-immune} survenant souvent chez l'enfant ou le jeune adulte, due à la \emph{destruction des cellules $\beta$} ; il se traite par \emph{insuline}. Le \cle{diabète gras} (type 2), beaucoup plus fréquent, n'est \emph{pas} auto-immun : il est lié au mode de vie (surcharge pondérale, sédentarité) et à une \emph{résistance des cellules à l'insuline}, l'insuline étant présente mais inefficace. Écrire \og le diabète est une maladie auto-immune \fg{} sans préciser le type est une faute.
\end{attention}

\courssub{La myasthénie : des auto-anticorps qui bloquent la commande musculaire}

\textbf{Observation.} Une malade se plaint d'une fatigue musculaire anormale : ses paupières tombent en fin de journée, elle peine à mastiquer, et l'effort prolongé devient impossible, alors que le repos la soulage. L'injection d'un inhibiteur de l'acétylcholinestérase améliore spectaculairement sa force musculaire.

\textbf{Interprétation.} La contraction musculaire est déclenchée par l'\cle{acétylcholine}, neurotransmetteur libéré par le motoneurone au niveau de la \cle{plaque motrice} (jonction neuromusculaire) : l'acétylcholine se fixe sur des récepteurs spécifiques de la fibre musculaire, ce qui provoque son excitation. Chez le myasthénique, on met en évidence des \cle{auto-anticorps anti-récepteurs de l'acétylcholine}. Ces auto-anticorps :
\begin{itemize}
\item se fixent sur les récepteurs et les \emph{bloquent} (l'acétylcholine ne peut plus agir) ;
\item provoquent la destruction des récepteurs par les complexes immunologiques formés.
\end{itemize}

\textbf{Conclusion.} Le message nerveux n'est plus transmis efficacement au muscle : il en résulte une \cle{faiblesse musculaire} fluctuante, majorée par l'effort. Ici, l'organe n'est pas détruit en masse : ce sont les \emph{récepteurs} qui sont neutralisés. L'inhibiteur de l'acétylcholinestérase augmente la durée de vie de l'acétylcholine dans la fente synaptique, ce qui compense partiellement le déficit de récepteurs fonctionnels.

\courssub{Le goitre exophtalmique (maladie de Basedow) : une stimulation anormale}

\textbf{Observation.} Un patient présente un gonflement de la base du cou (\cle{goitre}), des yeux globuleux et saillants (\cle{exophtalmie}), une accélération du cœur, un amaigrissement et une nervosité permanente. Son taux d'hormones thyroïdiennes (T3, T4) est très élevé, alors que sa TSH (hormone hypophysaire qui stimule normalement la thyroïde) est effondrée.

\textbf{Interprétation.} On découvre dans son sang des auto-anticorps dirigés contre les \emph{récepteurs de la TSH} de la thyroïde. Ces auto-anticorps, appelés \cle{TRAK} (ou TSI), se fixent sur les récepteurs TSH et \emph{miment l'action de la TSH} : ils stimulent la thyroïde \emph{en permanence}, sans aucun contrôle. La thyroïde, hyperstimulée, grossit (goitre) et sécrète un excès d'hormones thyroïdiennes (\cle{hyperthyroïdie}), responsable de la tachycardie, de l'amaigrissement et de l'exophtalmie. Par rétrocontrôle négatif, cet excès d'hormones inhibe l'hypophyse, ce qui explique l'effondrement de la TSH.

\begin{popfigure}
\begin{center}
\begin{tikzpicture}[scale=0.95, every node/.style={font=\small}]
% Cellule thyroïdienne
\fill[popgreenL] (0,0) rectangle (6.4,-2.4);
\draw[popgreen, very thick] (0,0) -- (6.4,0);
\node[popgreen!60!black] at (3.2,-1.2) {\textbf{Cellule thyroïdienne (thyrocyte)}};
\node[popdark] at (3.2,-1.95) {\footnotesize sécrète T3 et T4};
% Récepteur TSH
\draw[popgreen, very thick] (1.6,0) -- (1.6,0.9);
\draw[popgreen, very thick] (2.0,0) -- (2.0,0.9);
\draw[popgreen, very thick] (1.6,0.9) arc (180:0:0.2);
\node[popgreen!60!black, anchor=west] at (2.25,0.55) {\footnotesize récepteur à la TSH};
% Auto-anticorps en Y
\draw[poppink, very thick] (1.8,1.0) -- (1.8,1.75);
\draw[poppink, very thick] (1.8,1.75) -- (1.35,2.35);
\draw[poppink, very thick] (1.8,1.75) -- (2.25,2.35);
\node[poppink, anchor=south] at (1.8,2.4) {\footnotesize auto-anticorps (TRAK) mimant la TSH};
% Flèche stimulation
\draw[-{Stealth}, poporange, very thick] (1.8,-0.15) -- (1.8,-1.05);
\node[poporange, anchor=west] at (2.0,-0.6) {\footnotesize stimulation permanente};
% Conséquences
\node[popdark, anchor=west] at (7.0,0.2) {\footnotesize \textbf{Conséquences :}};
\node[popdark, anchor=west] at (7.0,-0.5) {\footnotesize hyperthyroïdie (excès de T3/T4)};
\node[popdark, anchor=west] at (7.0,-1.1) {\footnotesize goitre (thyroïde hypertrophiée)};
\node[popdark, anchor=west] at (7.0,-1.7) {\footnotesize exophtalmie, tachycardie};
% Rétrocontrôle
\draw[-{Stealth}, poppurple, thick, dashed] (5.6,-2.4) .. controls (5.6,-3.6) and (1.0,-3.6) .. (1.0,-2.5);
\node[poppurple] at (3.3,-3.85) {\footnotesize rétrocontrôle négatif : effondrement de la TSH hypophysaire};
\end{tikzpicture}
\end{center}
\legende{Mécanisme du goitre exophtalmique (maladie de Basedow) : les auto-anticorps TRAK se fixent sur les récepteurs à la TSH des thyrocytes et les stimulent en permanence, d'où une hyperthyroïdie et un goitre. L'excès d'hormones T3/T4 inhibe l'hypophyse (rétrocontrôle), d'où une TSH sanguine très basse.}
\end{popfigure}

\begin{savaistu}[Un auto-anticorps... qui stimule au lieu de détruire]
Dans la plupart des maladies auto-immunes, les effecteurs immunitaires \emph{détruisent} leur cible. Le goitre exophtalmique est un cas remarquable : l'auto-anticorps ne détruit pas la thyroïde, il l'\emph{hyperactive} en imitant la TSH. C'est l'un des rares exemples où l'auto-immunité provoque un \emph{hyperfonctionnement} d'un organe. La maladie a été décrite au \textsc{xix}\textsuperscript{e} siècle par le médecin allemand Karl von Basedow ; elle touche plus souvent les femmes et se traite par antithyroïdiens de synthèse, parfois par chirurgie ou iode radioactif.
\end{savaistu}

\courssub{Synthèse, facteurs favorisants et traitements}

\begin{center}
\begin{tabularx}{\linewidth}{|l|X|X|X|}
\hline
\rowcolor{popblueL} \textbf{Maladie} & \textbf{Organe cible} & \textbf{Mécanisme auto-immun} & \textbf{Symptômes} \\
\hline
Vitiligo & Mélanocytes de la peau & LT8 et auto-anticorps détruisant les mélanocytes & Plaques de dépigmentation cutanée \\
\hline
Diabète juvénile (type 1) & Cellules $\beta$ des îlots de Langerhans (pancréas) & Destruction par LT8 auto-réactifs et auto-anticorps $\Rightarrow$ absence d'insuline & Hyperglycémie chronique, soif intense, amaigrissement \\
\hline
Myasthénie & Récepteurs de l'acétylcholine (plaque motrice) & Auto-anticorps bloquant et détruisant les récepteurs & Faiblesse musculaire majorée à l'effort, ptosis \\
\hline
Goitre exophtalmique (Basedow) & Récepteurs à la TSH (thyroïde) & Auto-anticorps stimulant en permanence la thyroïde & Hyperthyroïdie, goitre, exophtalmie, tachycardie \\
\hline
\end{tabularx}
\end{center}

\textbf{Facteurs favorisants.} Pourquoi la tolérance du soi se rompt-elle chez certaines personnes ? Deux grandes catégories de facteurs sont admises :
\begin{itemize}
\item une \cle{prédisposition génétique} : certains allèles du \cle{CMH} (système HLA) sont statistiquement associés à un risque accru de maladie auto-immune (par exemple certains HLA de classe II pour le diabète de type 1) ;
\item des \cle{facteurs environnementaux} : infections virales ou bactériennes (un antigène microbien ressemblant à un auto-antigène peut activer des lymphocytes qui se retournent ensuite contre le soi : c'est le \emph{mimétisme moléculaire}), stress, certains médicaments.
\end{itemize}

\textbf{Traitements et limites.} Deux stratégies principales :
\begin{itemize}
\item les \cle{immunosuppresseurs} (corticoïdes, ciclosporine\ldots) freinent la réponse immunitaire fautive ; mais ils exposent le patient aux \emph{infections} et nécessitent une surveillance permanente ;
\item la \cle{substitution} : on remplace ce que l'organe détruit ne produit plus (injections d'\emph{insuline} dans le diabète de type 1, hormones thyroïdiennes après ablation de la thyroïde). Ce traitement soulage mais ne guérit pas la cause auto-immune.
\end{itemize}

\begin{exoresolu}[Analyse d'un cas clinique]
Énoncé. Une jeune fille de 14 ans consulte pour soif intense, urines abondantes et amaigrissement. Sa glycémie à jeun est de \SI{2,8}{g/L}. Le dosage sanguin révèle une insuline quasi nulle et la présence d'auto-anticorps anti-cellules $\beta$ des îlots de Langerhans. (1) Identifie la maladie. (2) Décris le mécanisme du dysfonctionnement en précisant l'organe cible, l'effecteur fautif et le lien de cause à effet. (3) Propose un traitement.
\tcblower
Solution. (1) Il s'agit du \cle{diabète juvénile} (diabète de type 1), une maladie auto-immune. (2) \emph{Organe cible} : les cellules $\beta$ des îlots de Langerhans du pancréas, seules productrices d'insuline. \emph{Effecteurs fautifs} : des lymphocytes T cytotoxiques auto-réactifs (qui détruisent les cellules $\beta$) et des auto-anticorps anti-cellules $\beta$, signes d'une rupture de la tolérance du soi. \emph{Lien de cause à effet} : destruction des cellules $\beta$ $\Rightarrow$ absence d'insuline $\Rightarrow$ le glucose ne pénètre plus dans les cellules $\Rightarrow$ hyperglycémie chronique (\SI{2,8}{g/L} à jeun), responsable de la soif, des urines abondantes et de l'amaigrissement. (3) Traitement \emph{substitutif} : injections quotidiennes d'insuline à vie, associées à une surveillance de la glycémie ; les immunosuppresseurs ne sont pas utilisables durablement ici à cause de leurs effets secondaires (risque infectieux).
\end{exoresolu}

\begin{exercice}[Pour t'entraîner]
Chez un patient atteint de la maladie de Basedow, on mesure : T3/T4 très élevées, TSH très basse, présence d'auto-anticorps anti-récepteurs TSH. (1) Explique pourquoi la TSH est basse alors que la thyroïde est hyperactive. (2) En quoi ce cas diffère-t-il du mécanisme du diabète juvénile ? (3) Rédige une description complète de ce dysfonctionnement en suivant la méthode en quatre étapes.
\end{exercice}

\begin{corrige}[Exercice : maladie de Basedow]
(1) L'excès d'hormones thyroïdiennes T3/T4 exerce un \emph{rétrocontrôle négatif} sur l'hypophyse, qui cesse de sécréter la TSH : la TSH s'effondre alors que la thyroïde reste stimulée par les auto-anticorps, insensibles au rétrocontrôle. (2) Dans le diabète juvénile, les effecteurs auto-immuns \emph{détruisent} l'organe cible (cellules $\beta$), entraînant un \emph{hypofonctionnement} (absence d'insuline) ; dans la maladie de Basedow, l'auto-anticorps \emph{stimule} l'organe cible, entraînant un \emph{hyperfonctionnement} (hyperthyroïdie). (3) \emph{Organe cible} : la thyroïde (récepteurs à la TSH des thyrocytes). \emph{Effecteur fautif} : auto-anticorps TRAK mimant la TSH. \emph{Symptômes} : goitre, exophtalmie, tachycardie, amaigrissement. \emph{Lien de cause à effet} : rupture de la tolérance du soi $\Rightarrow$ production d'auto-anticorps stimulants $\Rightarrow$ hyperstimulation permanente de la thyroïde $\Rightarrow$ hyperthyroïdie et goitre.
\end{corrige}

\begin{aretenir}[L'essentiel de la section]
\begin{itemize}
\item Une maladie auto-immune résulte d'une \cle{rupture de la tolérance du soi} : auto-anticorps et/ou lymphocytes T auto-réactifs s'attaquent à un organe cible.
\item Quatre exemples exigibles : \cle{vitiligo} (mélanocytes détruits), \cle{diabète de type 1} (cellules $\beta$ détruites $\Rightarrow$ hyperglycémie), \cle{myasthénie} (récepteurs de l'acétylcholine bloqués $\Rightarrow$ faiblesse musculaire), \cle{goitre exophtalmique} (récepteurs TSH stimulés $\Rightarrow$ hyperthyroïdie).
\item Facteurs favorisants : prédisposition génétique (CMH) et facteurs environnementaux (infections, stress).
\item Traitements : immunosuppresseurs (avec risque infectieux) et substitution (insuline), qui traitent les conséquences sans supprimer la cause.
\end{itemize}
\end{aretenir}

\coursec{Le VIH/Sida : immunodéficience acquise et enjeux socioculturels}

Dans la section précédente, nous avons vu que le système immunitaire peut se retourner \emph{contre le soi} : c'est le cas des maladies auto-immunes, où les défenses de l'organisme détruisent ses propres cellules. Mais il existe une situation inverse et tout aussi dramatique : le système immunitaire peut être \emph{détruit de l'extérieur}, non pas parce qu'il fonctionne mal, mais parce qu'un virus s'attaque directement à ses cellules clés. C'est le cas de l'infection par le \cle{VIH} (virus de l'immunodéficience humaine), responsable du \cle{Sida}. Au Cameroun comme dans toute l'Afrique subsaharienne, cette maladie constitue un problème majeur de santé publique et un enjeu social considérable. Comprendre ses mécanismes, c'est se donner les moyens de s'en protéger et de lutter contre les préjugés.

\courssub{Le VIH : un rétrovirus qui cible les chefs d'orchestre de l'immunité}

\begin{definition}[Le VIH]
Le \cle{VIH} (virus de l'immunodéficience humaine) est un \cle{rétrovirus} : son génome est constitué d'\cle{ARN} (et non d'ADN), et il possède une enzyme caractéristique, la \cle{transcriptase inverse}, capable de fabriquer de l'ADN à partir d'ARN. Il présente un \cle{tropisme} (une affinité) marqué pour les cellules portant le récepteur \cle{CD4} : les \cle{lymphocytes T4} (LT4), mais aussi les \cle{macrophages} et les \cle{cellules dendritiques}.
\end{definition}

L'intuition est simple : pour détruire une armée, le plus efficace n'est pas de tuer les soldats un à un, mais d'éliminer le \emph{général}. Or les lymphocytes T4 sont précisément les « chefs d'orchestre » de la réponse immunitaire spécifique : ils activent les lymphocytes B (réponse humorale) et les lymphocytes T cytotoxiques (réponse cellulaire). En infectant et en détruisant les LT4, le VIH paralyse l'ensemble des défenses de l'organisme.

Le virion (particule virale) est une petite sphère d'environ \SI{100}{nm} de diamètre, entourée d'une enveloppe portant des glycoprotéines (dont la \cle{gp120}) qui reconnaissent spécifiquement le récepteur CD4. À l'intérieur se trouvent deux copies d'ARN viral et les enzymes nécessaires à la réplication, dont la transcriptase inverse.

\courssub{Le cycle de multiplication du VIH dans un lymphocyte T4}

\begin{experience}[Observation : comment le virus se multiplie-t-il ?]
Des cultures de lymphocytes T4 humains sont mises en présence de VIH. L'observation au microscope électronique et le dosage des molécules virales à différents temps montrent successivement : la fixation des virions à la surface des cellules, la disparition de l'ARN viral du cytoplasme, l'apparition d'ADN viral \emph{intégré au chromosome} de la cellule hôte, puis la libération de nouveaux virions et la mort des lymphocytes. L'ajout d'un inhibiteur de la transcriptase inverse bloque l'apparition de l'ADN viral : cela prouve le rôle central de cette enzyme.
\end{experience}

Le cycle de réplication se déroule en étapes successives :

\begin{enumerate}
\item \textbf{Fixation} : la glycoprotéine gp120 du virus se fixe sur le récepteur \cle{CD4} du lymphocyte T4 (avec l'aide d'un corécepteur membranaire, CCR5 ou CXCR4).
\item \textbf{Entrée} : l'enveloppe virale fusionne avec la membrane cellulaire ; le contenu viral (ARN + transcriptase inverse) pénètre dans le cytoplasme.
\item \textbf{Transcription inverse} : la transcriptase inverse fabrique un \cle{ADN} viral double brin à partir de l'ARN viral.
\item \textbf{Intégration} : cet ADN viral pénètre dans le noyau et s'insère dans un chromosome de la cellule hôte : il devient un \cle{provirus}. C'est une étape décisive : le virus fait désormais partie du patrimoine génétique de la cellule, ce qui explique l'impossibilité actuelle de guérir l'infection.
\item \textbf{Production de nouveaux virions} : la cellule infectée, utilisant sa propre machinerie, transcrit le provirus et fabrique les protéines virales ; de nouveaux virions s'assemblent.
\item \textbf{Bourgeonnement et libération} : les virions sortent de la cellule en emportant un morceau de sa membrane (bourgeonnement), puis infectent d'autres LT4. La cellule infectée finit par être détruite, directement par le virus ou par les lymphocytes T cytotoxiques.
\end{enumerate}

\begin{popfigure}
\begin{center}
\begin{tikzpicture}[scale=0.95, every node/.style={font=\small}]
% Cellule LT4
\draw[fill=popblueL, thick, popblue] (0,0) ellipse (4.2 and 2.6);
\draw[fill=popgreenL, thick, popgreen] (0,-0.3) ellipse (1.5 and 1.1);
\node[popdark] at (0,-0.3) {\textbf{Noyau}};
\draw[thick, popdark] (-0.9,-0.15) .. controls (-0.3,0.15) and (0.3,-0.45) .. (0.9,-0.15);
\node[popdark, font=\scriptsize] at (0,-0.75) {ADN de la cellule};
% Récepteur CD4
\draw[thick, poppurple] (-3.1,2.15) -- (-3.1,2.75);
\draw[thick, poppurple] (-2.9,2.15) -- (-2.9,2.75);
\node[font=\scriptsize, poppurple] at (-3.0,2.95) {récepteur CD4};
% Virion extérieur
\draw[fill=poporange, thick, popdark] (-3.0,3.6) circle (0.45);
\node[white, font=\scriptsize] at (-3.0,3.6) {ARN};
\draw[->, thick, poporange] (-3.0,3.1) -- (-3.0,2.8);
\node[poporange, font=\scriptsize, anchor=west] at (-2.4,3.6) {\textbf{1. Fixation}};
% Transcription inverse
\draw[->, thick, poporange] (-2.5,1.9) -- (-1.2,1.3);
\node[poporange, font=\scriptsize, anchor=west] at (-2.6,1.5) {\textbf{2. Entrée}};
\draw[->, thick, poporange] (-1.0,1.4) -- (-0.6,0.9);
\node[poporange, font=\scriptsize, anchor=east] at (-1.1,1.15) {\textbf{3. Transcription inverse} (ARN $\rightarrow$ ADN)};
% Intégration
\draw[->, thick, poporange] (0.2,0.9) -- (0.2,0.35);
\node[poporange, font=\scriptsize, anchor=west] at (0.35,0.62) {\textbf{4. Intégration (provirus)}};
% Production
\draw[->, thick, poporange] (1.4,0.3) -- (2.6,1.1);
\node[poporange, font=\scriptsize, anchor=west] at (1.7,0.55) {\textbf{5. Production}};
% Bourgeonnement
\draw[fill=poporange, thick, popdark] (3.6,2.6) circle (0.35);
\node[white, font=\scriptsize] at (3.6,2.6) {ARN};
\draw[->, thick, poporange] (3.3,1.9) -- (3.55,2.25);
\node[poporange, font=\scriptsize, anchor=west] at (3.0,1.6) {\textbf{6. Bourgeonnement}};
\node[popdark, font=\small] at (0,-3.1) {\textbf{Lymphocyte T4 infecté}};
\end{tikzpicture}
\end{center}
\legende{Cycle de réplication du VIH dans un lymphocyte T4 : 1. fixation de la gp120 sur le récepteur CD4 ; 2. entrée du contenu viral ; 3. transcription inverse de l'ARN viral en ADN par la transcriptase inverse ; 4. intégration de l'ADN viral (provirus) dans le chromosome de la cellule ; 5. production de nouveaux virions ; 6. libération par bourgeonnement et infection de nouvelles cellules.}
\end{popfigure}

\begin{attention}[Une conséquence capitale de l'intégration]
Une fois l'ADN viral intégré dans le génome de la cellule, il est \emph{impossible} de l'en déloger avec les traitements actuels. C'est pourquoi on ne guérit pas du VIH : les traitements bloquent la \emph{réplication}, mais le provirus persiste dans des cellules « réservoirs ».
\end{attention}

\courssub{L'évolution naturelle de l'infection : de la séropositivité au Sida}

\begin{definition}[Sida]
Le \cle{Sida} (syndrome d'immunodéficience acquise) est le \emph{stade ultime} de l'infection par le VIH. Il est caractérisé par un effondrement des défenses immunitaires (taux de LT4 inférieur à \SI{200}{/mm^3} de sang, contre \num{800} à \num{1200} normalement) et par l'apparition d'\cle{infections opportunistes} (tuberculose, pneumocystose, candidoses) et de \cle{cancers} (maladie de Kaposi).
\end{definition}

Le suivi médical de personnes infectées non traitées, par dosage régulier du taux de LT4 et de la \cle{charge virale} (nombre de copies d'ARN viral par millilitre de sang), a permis d'établir la courbe d'évolution de l'infection, qui comporte trois phases.

\begin{popfigure}
\begin{center}
\begin{tikzpicture}
\begin{axis}[
width=0.95\linewidth, height=7.5cm,
xlabel={Temps écoulé depuis la contamination (années)},
ylabel={Taux de LT4 (cellules/mm$^3$)},
axis y line*=left, axis x line*=bottom,
xmin=0, xmax=12, ymin=0, ymax=1300,
xtick={0,2,4,6,8,10,12},
ytick={0,200,400,600,800,1000,1200},
legend style={at={(0.35,0.97)}, anchor=north, font=\small},
]
\addplot[popcurve, color=popblue, very thick, domain=0:12, samples=200]
{1000 - 850/(1+exp(-(x-8))) - 120*exp(-((x-0.4)^2)/0.08)};
\addlegendentry{Taux de LT4}
\addplot[popcurve, color=poporange, very thick, dashed, domain=0:12, samples=200]
{1200*exp(-((x-0.5)^2)/0.5) + 60 + 1050/(1+exp(-(x-9)))};
\addlegendentry{Charge virale (échelle arbitraire)}
\addplot[popline, dotted, thick, domain=0:12]{200};
\node[popmuted, font=\scriptsize] at (axis cs:1.2,260) {seuil : 200/mm$^3$};
\node[popdark, font=\small] at (axis cs:0.9,1250) {\textbf{Phase aiguë}};
\node[popdark, font=\small] at (axis cs:5.2,1250) {\textbf{Phase de latence}};
\node[popdark, font=\small] at (axis cs:10.6,1250) {\textbf{Stade Sida}};
\end{axis}
\end{tikzpicture}
\end{center}
\legende{Évolution du taux de LT4 (courbe bleue) et de la charge virale (courbe orange, en échelle arbitraire) au cours d'une infection à VIH non traitée. Phase aiguë : multiplication virale intense et chute transitoire des LT4. Phase de latence asymptomatique (plusieurs années) : déclin lent et régulier des LT4. Stade Sida : le taux de LT4 passe sous le seuil de \SI{200}{/mm^3}, la charge virale remonte et les infections opportunistes apparaissent.}
\end{popfigure}

\begin{itemize}
\item \textbf{Phase aiguë (primo-infection)} : dans les semaines qui suivent la contamination, le virus se multiplie massivement (charge virale très élevée) et le taux de LT4 chute transitoirement. La personne peut présenter un syndrome grippal discret, souvent ignoré. Elle devient \cle{séropositive} : son organisme fabrique des anticorps anti-VIH, détectables par les tests de dépistage.
\item \textbf{Phase de latence asymptomatique} : pendant plusieurs années (souvent 8 à 10 ans, parfois plus), la personne ne présente \emph{aucun symptôme} et se sent en bonne santé. Pourtant, le virus continue de se multiplier et détruit lentement les LT4, dont le taux décline progressivement.
\item \textbf{Stade Sida} : lorsque le taux de LT4 passe sous \SI{200}{/mm^3}, les défenses immunitaires sont effondrées. Des micro-organismes normalement inoffensifs ou bien contrôlés provoquent des \cle{infections opportunistes} graves : tuberculose (très fréquente au Cameroun), pneumocystose (infection pulmonaire), candidoses (mycoses), ainsi que des cancers comme la maladie de Kaposi.
\end{itemize}

\begin{methode}[Exploiter la courbe LT4 / charge virale]
\begin{enumerate}
\item Repérer les \textbf{trois phases} sur l'axe des temps : phase aiguë (pic de charge virale précoce), phase de latence (longue période de déclin lent des LT4), stade Sida (effondrement final).
\item Comparer les \textbf{deux courbes} : quand la charge virale est élevée, le taux de LT4 diminue : la destruction des LT4 est la conséquence directe de la multiplication virale.
\item Repérer le \textbf{seuil critique} de \SI{200}{LT4/mm^3} : en dessous, l'immunité ne peut plus contrôler les micro-organismes de l'environnement, d'où l'apparition des infections opportunistes.
\item Conclure : la gravité du Sida ne vient pas du virus lui-même, mais de l'\emph{immunodépression} qu'il provoque.
\end{enumerate}
\end{methode}

\begin{attention}[Erreur fréquente : séropositif ne veut pas dire malade du Sida]
Être \cle{séropositif} signifie être \emph{porteur du VIH} (test de dépistage positif), pas être malade du Sida. Grâce à la phase de latence, qui peut durer \emph{plus de 10 ans}, une personne séropositive peut paraître en parfaite santé tout en pouvant transmettre le virus. C'est précisément ce qui rend le dépistage indispensable : on ne peut pas deviner la séropositivité de quelqu'un à son apparence.
\end{attention}

\courssub{Les modes de transmission et les moyens de prévention}

Le VIH se transmet uniquement par trois voies :

\begin{center}
\begin{tabularx}{\linewidth}{|l|X|}\hline
\rowcolor{popblueL} \textbf{Voie de transmission} & \textbf{Situations concrètes} \\ \hline
Rapports sexuels non protégés & Contact sexuel sans préservatif avec une personne infectée (voie la plus fréquente) \\ \hline
Sang contaminé & Transfusion de sang non dépisté, partage de seringues, objets tranchants souillés (lames, aiguilles, instruments de scarification ou de circoncision non stérilisés) \\ \hline
Mère-enfant & Pendant la grossesse, l'accouchement ou l'allaitement (transmission verticale) \\ \hline
\end{tabularx}
\end{center}

\begin{attention}[Ce qui ne transmet PAS le VIH]
Le VIH ne se transmet \textbf{pas} par les contacts de la vie quotidienne : poignées de main, embrassades, partage de repas ou de toilettes, piqûres de moustiques, sueur, larmes. Rejeter une personne vivant avec le VIH sur ces bases est une injustice dépourvue de tout fondement scientifique.
\end{attention}

La prévention repose sur plusieurs piliers complémentaires :

\begin{itemize}
\item l'utilisation du \cle{préservatif} lors des rapports sexuels ;
\item le \cle{dépistage volontaire et confidentiel}, seul moyen de connaître son statut sérologique ;
\item la \cle{fidélité} réciproque entre partenaires non infectés, ou l'\cle{abstinence} ;
\item la sécurisation du sang (dépistage systématique avant transfusion) et l'usage de matériel stérile ou à usage unique ;
\item la \cle{PTME} (prévention de la transmission mère-enfant) : une femme enceinte séropositive, suivie et traitée, peut mettre au monde un enfant non infecté dans la grande majorité des cas ;
\item le \cle{TARV} (traitement antirétroviral), qui bloque la réplication virale.
\end{itemize}

\begin{propriete}[Effet des antirétroviraux (admise)]
Les \cle{antirétroviraux} ne guérissent \emph{pas} du Sida : ils ne détruisent pas le provirus intégré au génome des cellules. En revanche, en bloquant des étapes clés de la réplication virale (notamment la transcriptase inverse, la protéase ou l'intégrase), ils maintiennent la charge virale à un niveau très faible, préservent les LT4 et permettent une vie quasi normale et longue, \emph{à condition d'une observance stricte} (prise quotidienne du traitement, à vie). Une charge virale devenue indétectable rend en outre la transmission sexuelle pratiquement impossible (concept I = I : Indétectable = Intransmissible).
\end{propriete}

\begin{savaistu}[Le dépistage et le traitement au Cameroun]
Au Cameroun, le dépistage du VIH est proposé \emph{gratuitement} dans de nombreux centres de santé et lors de campagnes de dépistage communautaire, et les \emph{antirétroviraux sont délivrés gratuitement depuis 2007}. Des associations de personnes vivant avec le VIH accompagnent les malades. Se faire dépister, c'est un acte de responsabilité envers soi-même et envers les autres.
\end{savaistu}

\courssub{Conséquences socioculturelles et compétences psychosociales}

Au-delà de la maladie elle-même, le VIH/Sida a des répercussions sociales profondes au Cameroun :

\begin{itemize}
\item la \cle{stigmatisation} et la \cle{discrimination} : rejet familial, perte d'emploi, exclusion sociale des personnes vivant avec le VIH (PVVIH), souvent alimentés par l'ignorance et les rumeurs ;
\item les \cle{orphelins du Sida}, enfants ayant perdu un ou deux parents, dont la prise en charge pèse sur les familles élargies et les communautés ;
\item l'\cle{impact économique} : la maladie frappe des adultes jeunes et productifs, appauvrissant les ménages (frais de soins, perte de revenus) et freinant le développement.
\end{itemize}

Face à cela, chacun peut développer des \cle{compétences psychosociales} (savoir-être) essentielles :

\begin{itemize}
\item \textbf{Vivre positivement avec le VIH} : accepter son statut, observer rigoureusement le traitement antirétroviral, adopter une alimentation équilibrée, rechercher un soutien psychologique (famille, associations, personnel de santé) et continuer à se projeter dans l'avenir.
\item \textbf{Accueillir sans discrimination} les PVVIH : leur témoigner de la solidarité, respecter la confidentialité de leur statut, refuser les rumeurs et les moqueries.
\item \textbf{Refuser les comportements à risque} : savoir dire non aux rapports non protégés, négocier l'usage du préservatif, résister à la pression du groupe.
\end{itemize}

\begin{exoresolu}[Exploiter un document médical]
\textbf{Énoncé.} Le dossier médical d'un patient non traité donne les résultats suivants : en 2015, taux de LT4 = \SI{850}{/mm^3} ; en 2019, taux de LT4 = \SI{450}{/mm^3} ; en 2023, taux de LT4 = \SI{120}{/mm^3} et apparition d'une tuberculose pulmonaire. 1) À quelle phase de l'infection se trouvait le patient en 2015 ? Justifier. 2) Expliquer l'apparition de la tuberculose en 2023. 3) Quel traitement aurait pu modifier cette évolution, et selon quel mécanisme ?
\tcblower
\textbf{Solution.} 1) En 2015, le taux de LT4 est proche de la normale (\num{800} à \SI{1200}{/mm^3}) : le patient est en \emph{phase de latence asymptomatique} ; il est séropositif mais pas malade du Sida. 2) En 2023, le taux de LT4 est tombé à \SI{120}{/mm^3}, donc sous le seuil critique de \SI{200}{/mm^3} : les défenses immunitaires sont effondrées. Le bacille de la tuberculose, normalement contrôlé par l'immunité à médiation cellulaire (dont les LT4 sont les chefs d'orchestre), peut alors se multiplier : c'est une \emph{infection opportuniste}, signe du passage au stade Sida. 3) Un traitement antirétroviral (TARV), pris dès le diagnostic de séropositivité et avec une observance stricte, aurait bloqué la réplication virale (notamment en inhibant la transcriptase inverse), maintenu la charge virale basse, préservé le taux de LT4 et empêché l'évolution vers le Sida — sans toutefois guérir l'infection.
\end{exoresolu}

\begin{exercice}[Pour t'entraîner]
Une élève affirme : « Mon voisin a l'air fatigué et maigre, il a sûrement le Sida ; il ne faut plus manger avec lui. » Relève toutes les erreurs scientifiques et les comportements inadaptés contenus dans cette phrase, et rédige la réponse argumentée que tu lui ferais.
\end{exercice}

\begin{corrige}[Exercice]
Erreurs : (1) on ne peut pas diagnostiquer le Sida à l'apparence ; seul un test de dépistage révèle la séropositivité, et la fatigue ou la maigreur ont d'innombrables causes ; (2) une personne séropositive peut être en parfaite santé apparente, et inversement ; (3) le VIH ne se transmet pas en partageant un repas ; (4) le rejet relève de la stigmatisation, injustifiée scientifiquement et blessante humainement. Réponse attendue : expliquer calmement les vraies voies de transmission (rapports non protégés, sang contaminé, mère-enfant), rappeler que les contacts quotidiens sont sans danger, encourager le dépistage volontaire et confidentiel plutôt que les soupçons, et inviter à la solidarité envers les personnes vivant avec le VIH.
\end{corrige}

\begin{aretenir}[L'essentiel de la section]
Le VIH est un rétrovirus à ARN qui, grâce à sa transcriptase inverse, intègre son génome dans celui des lymphocytes T4, des macrophages et des cellules dendritiques. La destruction progressive des LT4 conduit, après une longue phase asymptomatique, à l'immunodépression : c'est le Sida (LT4 $<$ \SI{200}{/mm^3}), marqué par les infections opportunistes. Le virus se transmet par voie sexuelle, sanguine et de la mère à l'enfant — jamais par les contacts quotidiens. Préservatif, dépistage, fidélité, PTME et antirétroviraux (gratuits au Cameroun depuis 2007) sont les armes de la lutte ; la solidarité et le refus de la stigmatisation en sont le visage humain.
\end{aretenir}

\newpage

\coursec{Activité d'intégration}

\courssub{Objectifs de l'activité}

À l'issue de cette activité d'intégration, tu seras capable de :
\begin{itemize}
\item identifier, à partir de documents cliniques et expérimentaux, le \cle{type de dysfonctionnement} du système immunitaire (maladie auto-immune ou immunodéficience) ;
\item mettre en relation un symptôme avec l'\cle{effecteur immunitaire} en cause (lymphocyte B, lymphocyte T cytotoxique, anticorps, complexe immun) ;
\item interpréter des données chiffrées (glycémie, taux de lymphocytes T4, charge virale) et des micrographies ;
\item expliquer le mécanisme d'une réaction auto-immune et celui de la destruction des LT4 par le VIH ;
\item proposer des attitudes de \cle{prévention}, de \cle{prise en charge} et de \cle{solidarité} envers les personnes vivant avec le VIH (compétences psychosociales).
\end{itemize}

\courssub{Situation}

Tu es élève en Terminale D à Bafoussam. Dans le cadre de la semaine de la santé organisée par le club sciences de ton lycée, le médecin-chef de l'hôpital régional vous soumet deux dossiers médicaux anonymisés et vous demande, par groupes de quatre, de jouer les « experts immunologues ».

\textbf{Dossier A — Patient « M. », adolescent de 15 ans.}

M., scolarisé en classe de Seconde, se plaint depuis trois mois d'une \cle{soif intense} (polydipsie), d'urines abondantes (polyurie) et d'un amaigrissement malgré un appétit conservé. Les résultats du laboratoire d'analyses médicales sont les suivants :

\begin{center}
\begin{tabularx}{\linewidth}{|l|X|X|}
\hline
\rowcolor{popblueL} \textbf{Paramètre} & \textbf{Valeur mesurée} & \textbf{Valeur de référence} \\
\hline
Glycémie à jeun & \SI{3,0}{g/L} (soit \SI{16,7}{mmol/L}) & 0,70 à \SI{1,10}{g/L} \\
\hline
Auto-anticorps anti-cellules $\beta$ des îlots de Langerhans & Présents & Absents \\
\hline
Anticorps anti-insuline & Présents & Absents \\
\hline
\end{tabularx}
\end{center}

La micrographie fournie montre un îlot de Langerhans du pancréas \cle{infiltré par de nombreux lymphocytes}, avec une destruction presque totale des cellules $\beta$ sécrétrices d'insuline.

\textbf{Dossier B — Patient « K. », adulte de 35 ans.}

K., commerçant, consulte pour des infections répétées : une \cle{tuberculose} pulmonaire traitée il y a deux ans, une \cle{candidose} buccale récidivante et des diarrhées chroniques. Le bilan immunologique et virologique donne :

\begin{center}
\begin{tabularx}{\linewidth}{|l|X|X|}
\hline
\rowcolor{popblueL} \textbf{Paramètre} & \textbf{Valeur mesurée} & \textbf{Valeur de référence} \\
\hline
Taux de lymphocytes T4 (LT4) & 180 cellules/mm\textsuperscript{3} & 800 à 1200 cellules/mm\textsuperscript{3} \\
\hline
Charge virale VIH & 85\,000 copies/mL & Indétectable chez un sujet sain \\
\hline
Test de dépistage VIH (ELISA, confirmé par Western blot) & Positif & Négatif \\
\hline
\end{tabularx}
\end{center}

Le suivi du patient sur 8 ans a permis de tracer l'évolution de son taux de LT4 :

\begin{popfigure}
\begin{center}
\begin{tikzpicture}
\begin{axis}[
width=0.85\linewidth, height=6cm,
xlabel={Temps (années)}, ylabel={LT4 (cellules/mm\textsuperscript{3})},
xmin=0, xmax=8, ymin=0, ymax=1200,
xtick={0,1,2,3,4,5,6,7,8},
grid=both, grid style={popline!40},
legend style={at={(0.02,0.02)}, anchor=south west, font=\small}
]
\addplot[popcurve, domain=0:8, samples=100] {1100 - 115*x};
\addlegendentry{Taux de LT4 du patient K.}
\addplot[dashed, poporange, thick, domain=0:8] {200};
\addlegendentry{Seuil critique : 200 cellules/mm\textsuperscript{3}}
\end{axis}
\end{tikzpicture}
\end{center}
\legende{Évolution du taux de lymphocytes T4 du patient K. sur 8 ans. En dessous de 200 cellules/mm\textsuperscript{3}, les infections opportunistes se multiplient : c'est le stade Sida.}
\end{popfigure}

\courssub{Consignes}

\begin{enumerate}
\item Pour chaque dossier, identifie le \textbf{type de dysfonctionnement immunitaire} en cause (réaction auto-immune ou immunodéficience) en justifiant ta réponse par \textbf{deux arguments} tirés des documents.
\item \textbf{Dossier A :} nomme l'effecteur immunitaire responsable de la destruction des cellules $\beta$ et explique, en trois étapes, le mécanisme qui conduit à l'hyperglycémie. Quelle maladie le médecin diagnostique-t-il ?
\item \textbf{Dossier B :} à l'aide de la courbe, détermine l'année approximative où le patient est passé sous le seuil critique de 200 cellules/mm\textsuperscript{3}. Explique pourquoi la chute des LT4 entraîne des infections \textbf{opportunistes} (tuberculose, candidose) alors qu'un sujet sain y résiste.
\item \textbf{Dossier B :} propose un plan de prise en charge réaliste (traitement, observance, suivi biologique) et rédige un court message (5 à 8 lignes) de sensibilisation contre la stigmatisation, destiné à être lu à la communauté scolaire lors de la semaine de la santé.
\item \textbf{Restitution orale :} à l'aide de la maquette d'anticorps ou de complexe immun construite en classe, illustre devant le groupe comment les auto-anticorps du dossier A se fixent sur les cellules $\beta$ et pourquoi cette fixation entraîne leur destruction.
\end{enumerate}

\courssub{Ressources à mobiliser}

\begin{aretenir}[Ce qu'il faut avoir en tête]
\begin{itemize}
\item \textbf{Réponse immunitaire spécifique :} les lymphocytes B (LB) produisent des \cle{anticorps} (réponse à médiation humorale) ; les lymphocytes T cytotoxiques (LT8 ou CD8\textsuperscript{+}) détruisent les cellules reconnues comme étrangères (réponse à médiation cellulaire) ; les lymphocytes T4 (LT4, auxiliaires ou CD4\textsuperscript{+}) \textbf{coordonnent} les deux réponses en sécrétant des interleukines (IL-2, IFN-$\gamma$, etc.).
\item \textbf{Tolérance et auto-immunité :} normalement, le système immunitaire tolère le « soi » (tolérance centrale dans le thymus et la moelle osseuse, tolérance périphérique). Une \cle{maladie auto-immune} apparaît quand des lymphocytes ou des anticorps échappent à la tolérance et attaquent les propres constituants de l'organisme (exemples au programme : diabète juvénile, myasthénie, goitre exophtalmique, vitiligo).
\item \textbf{VIH/Sida :} le VIH (VIH-1 au Cameroun) infecte les cellules portant le récepteur \textbf{CD4} (principalement les LT4) après fixation sur un co-récepteur (CCR5 ou CXCR4) ; il s'y multiplie et entraîne leur destruction progressive. Quand le taux de LT4 passe sous \textbf{200 cellules/mm\textsuperscript{3}}, l'immunodépression permet le développement d'\cle{infections opportunistes} (tuberculose, candidose, pneumocystose, toxoplasmose...) : c'est le stade \cle{Sida}. Les \cle{TARV} (traitements antirétroviraux, trithérapies) bloquent la multiplication virale (inhibiteurs de la transcriptase inverse, de la protéase, de l'intégrase, etc.) et permettent de vivre longtemps avec une charge virale indétectable.
\item \textbf{Structures de reconnaissance :} l'anticorps (immunoglobuline) est une molécule en Y (deux chaînes lourdes, deux chaînes légères) dont les deux bras portent les \cle{sites de fixation} (domaines variables) spécifiques d'un \cle{antigène} (épitope) ; le pied (domaine constant, Fc) permet la fixation des cellules effectrices (macrophages, NK, complément). Le complexe antigène-anticorps formé est un \cle{complexe immun}, éliminé par phagocytose (opsonisation) ou activation du complément.
\end{itemize}
\end{aretenir}

\courssub{Démarche et corrigé commenté}

\begin{exoresolu}[Deux dossiers médicaux à l'hôpital régional]
Énoncé : à partir des deux dossiers ci-dessus, identifier le dysfonctionnement immunitaire de chaque patient, expliquer les mécanismes, proposer une prise en charge pour le dossier B et un message de sensibilisation, puis illustrer le dossier A avec une maquette.

\tcblower

\textbf{Question 1 — Identification des dysfonctionnements.}

\emph{Dossier A :} il s'agit d'une \cle{maladie auto-immune}. 
\begin{itemize}
\item \textbf{Argument 1 :} la présence d'\textbf{auto-anticorps} dirigés contre les cellules $\beta$ des îlots de Langerhans et contre l'insuline elle-même (marqueurs spécifiques de l'auto-immunité anti-îlots).
\item \textbf{Argument 2 :} la micrographie montre un \textbf{infiltrat lymphocytaire dense} (insulite) au sein de l'îlot, signe que les propres lymphocytes T du patient reconnaissent les antigènes des cellules $\beta$ comme non-soi et les agressent. Le système immunitaire attaque ici le « soi ».
\end{itemize}

\emph{Dossier B :} il s'agit d'une \cle{immunodéficience acquise} (Sida). 
\begin{itemize}
\item \textbf{Argument 1 :} le test de dépistage VIH est \textbf{positif} (ELISA confirmé par Western blot) avec une charge virale élevée (85\,000 copies/mL), prouvant une réplication virale active.
\item \textbf{Argument 2 :} le taux de LT4 est effondré (180 cellules/mm\textsuperscript{3}, soit environ 6 fois moins que la normale), ce qui explique cliniquement les infections répétées. Le système immunitaire est ici \textbf{déficient} (incapable de défendre l'organisme), et non retourné contre le soi.
\end{itemize}

\textbf{Question 2 — Mécanisme du dossier A : le diabète juvénile (diabète de type 1).}

L'effecteur en cause associe les deux bras de la réponse spécifique : des \textbf{lymphocytes T cytotoxiques auto-réactifs (LT8/CD8\textsuperscript{+})} qui détruisent directement les cellules $\beta$ (médiation cellulaire) et des \textbf{auto-anticorps} produits par des lymphocytes B auto-réactifs (médiation humorale, marqueurs diagnostiques). Le mécanisme se déroule en trois étapes :
\begin{enumerate}
\item \textbf{Rupture de la tolérance au soi :} des lymphocytes T (LT4 et LT8) et B reconnaissant les antigènes propres des cellules $\beta$ (insuline, GAD65, IA-2, ZnT8), qui auraient dû être éliminés ou anergisés lors de la sélection centrale (thymus) ou périphérique, s'activent anormalement (facteurs génétiques HLA-DR3/DR4 + facteurs environnementaux : virus, alimentation...).
\item \textbf{Destruction des cellules $\beta$ :} les LT8 cytotoxiques lyse nt les cellules $\beta$ par relargage de perforines et granzymes (apoptose) ; parallèlement, les auto-anticorps forment des \cle{complexes immuns} à la surface des cellules $\beta$, amplifiant leur élimination par les macrophages (phagocytose médiée par les récepteurs Fc) et l'activation du complément (d'où l'infiltrat lymphocytaire et macrophagique visible sur la micrographie).
\item \textbf{Conséquence métabolique :} la masse des cellules $\beta$ s'effondre ; la sécrétion d'\textbf{insuline} (hormone hypoglycémiante unique permettant l'entrée du glucose dans les cellules cibles : foie, muscle, tissu adipeux via GLUT4) devient insuffisante. Le glucose s'accumule dans le sang : la glycémie atteint \SI{3,0}{g/L} (hyperglycémie majeure), dépassant le seuil de réabsorption rénale $\rightarrow$ glycosurie $\rightarrow$ polyurie osmotique $\rightarrow$ polydipsie compensatrice ; l'absence d'insuline empêche l'anabolisme $\rightarrow$ catabolisme protéique et lipidique $\rightarrow$ amaigrissement malgré l'hyperphagie.
\end{enumerate}
Le diagnostic est donc le \cle{diabète juvénile} (diabète insulinodépendant de type 1), l'une des maladies auto-immunes citées au programme, au même titre que la myasthénie (anticorps anti-récepteur à l'acétylcholine), le goitre exophtalmique de Basedow (anticorps stimulants du récepteur à la TSH) ou le vitiligo (destruction auto-immune des mélanocytes).

\textbf{Question 3 — Mécanisme du dossier B : de l'infection à VIH au Sida.}

Sur la courbe, le taux de LT4 décroît régulièrement selon une relation affine $y = 1100 - 115x$ (avec $y$ en cellules/mm\textsuperscript{3} et $x$ en années), passant d'environ 1100 à 180 cellules/mm\textsuperscript{3}. La droite franchit le seuil critique de 200 cellules/mm\textsuperscript{3} pour $1100 - 115x = 200 \Rightarrow x = 900/115 \approx 7,8$ ans, soit vers la \textbf{8\textsuperscript{e} année} de suivi. Cela correspond à l'apparition des infections graves documentées : le patient entre au stade \cle{Sida}.

\textbf{Mécanisme moléculaire et cellulaire :}
\begin{enumerate}
\item Le VIH se fixe via sa glycoprotéine gp120 sur le récepteur \textbf{CD4} des LT4, puis sur un \textbf{co-récepteur} (CCR5 pour les souches macrophagotropes initiales, CXCR4 pour les souches lymphotropes tardives).
\item La fusion virale permet l'entrée du capside ; la \textbf{transcriptase inverse} synthétise l'ADN viral à partir de l'ARN génomique ; l'\textbf{intégrase} insère cet ADN (provirus) dans le génome de la cellule hôte.
\item La cellule devient une « usine à virus » : transcription, traduction, assemblage, bourgeonnement. Chaque cycle de réplication détruit la cellule (lyse, apoptose, syncytia) et infecte de nouveaux LT4.
\end{enumerate}
Or les LT4 sont les \textbf{chefs d'orchestre} de l'immunité adaptative : par leurs interleukines (IL-2, IL-4, IFN-$\gamma$, etc.), ils activent la prolifération et la différenciation des LB en plasmocytes (anticorps) et des LT8 en effecteurs cytotoxiques. Leur disparition progressive \textbf{paralyse les deux bras de la réponse spécifique}. Des germes normalement inoffensifs ou bien contrôlés par l'immunité cellulaire (\emph{Mycobacterium tuberculosis}, \emph{Candida albicans}, \emph{Pneumocystis jirovecii}, \emph{Toxoplasma gondii}) ne sont plus éliminés : ce sont les \cle{infections opportunistes}, caractéristiques de l'immunodépression profonde.

\textbf{Question 4 — Prise en charge et message de sensibilisation.}

Plan de prise en charge du patient K. (conforme aux directives OMS et Programme National de Lutte contre le Sida du Cameroun) :
\begin{itemize}
\item \textbf{Mise sous TARV immédiate} (trithérapie antirétrovirale : 2 inhibiteurs de la transcriptase inverse + 1 inhibiteur de l'intégrase, ex. Ténofovir + Lamivudine + Dolutégravir), gratuite dans les centres de traitement agréés (CTA) au Cameroun.
\item \textbf{Observance stricte et à vie :} prise quotidienne à heure fixe, sans interruption, pour maintenir la charge virale indétectable et éviter la sélection de mutants résistants.
\item \textbf{Suivi biologique régulier :} charge virale plasmatique à 6 mois puis tous les 6 à 12 mois (objectif : \textbf{charge virale indétectable} $<$ 50 copies/mL = non-transmissibilité sexuelle, « I = I » : Indétectable = Intransmissible) ; numération des LT4 tous les 3 à 6 mois pour vérifier la reconstitution immunitaire.
\item \textbf{Prise en charge des co-infections et comorbidités :} traitement antituberculeux (si TB active), antifongiques (candidose), réhydratation et nutrition adaptée (diarrhées chroniques), vaccination (hépatite B, pneumocoque, grippe), dépistage IST, accompagnement psychosocial et groupe de parole (GAS - Groupes d'Auto-Soutien).
\end{itemize}

Message de sensibilisation (exemple, $\approx$ 7 lignes) : 
\emph{« Chers camarades, le VIH ne se transmet pas en serrant la main, en partageant un repas, en jouant ensemble ou par les moustiques. Une personne vivant avec le VIH, correctement traitée et dont la charge virale est indétectable, vit, étudie, travaille et aime normalement : elle ne transmet pas le virus. Rejeter ou moquer un camarade séropositif, c'est le pousser à cacher sa maladie et à abandonner son traitement ; l'entourer de respect et d'amitié, c'est l'aider à observer son traitement et à rester en bonne santé. Dépistons-nous volontairement, protégeons-nous (préservatifs, PrEP), et faisons preuve de solidarité : la discrimination tue, la fraternité soigne. »}

\textbf{Question 5 — Restitution avec la maquette.}

Sur la maquette, l'anticorps est figuré par une molécule en \textbf{Y} : les deux bras portent les \cle{sites de fixation} (domaines variables $V_H/V_L$) spécifiques de l'antigène (ici l'antigène de surface des cellules $\beta$ ou l'insuline), le pied (domaine constant $F_c$) permet la fixation des récepteurs des macrophages (ou du complément C1q). En apposant les « anticorps » de la maquette sur les « cellules $\beta$ » (pastilles portant l'antigène $\beta$), on montre la formation du \cle{complexe immun} antigène-anticorps à la surface cellulaire, puis sa reconnaissance et sa phagocytose par un « macrophage » (opsonisation). On illustre ainsi que, dans le diabète juvénile, la spécificité et l'efficacité de la réaction antigène-anticorps — normalement protectrice contre le non-soi (bactéries, virus) — se retournent tragiquement contre un constituant du soi.
\end{exoresolu}

\begin{attention}[Pièges à éviter dans ce type d'activité]
\begin{itemize}
\item Ne confonds pas \textbf{maladie auto-immune} (le système immunitaire est hyperactif mais mal dirigé : il attaque le soi) et \textbf{immunodéficience} (le système immunitaire est affaibli, hypoactif : il ne défend plus l'organisme contre le non-soi).
\item Une glycémie normale à jeun ne dépasse pas \SI{1,10}{g/L} (critère OMS 1999 utilisé dans le programme) : \SI{3,0}{g/L} est une hyperglycémie majeure, \textbf{pas} une hypoglycémie. L'hypoglycémie est $<$ 0,70 g/L.
\item Le VIH n'est pas le Sida : le Sida est le \textbf{stade ultime} de l'infection, défini cliniquement (infections opportunistes, cancers) et/ou biologiquement (LT4 $<$ 200/mm\textsuperscript{3}). On peut être séropositif (VIH+) sans être au stade Sida.
\item Les TARV ne \textbf{guérissent} pas l'infection à VIH (le provirus intégré dans les réservoirs latents persiste) : ils bloquent la multiplication virale et restaurent l'immunité \textbf{tant que l'observance est respectée}. L'arrêt du traitement entraîne une reprise virale rapide.
\item Attention au vocabulaire : on dit « \textbf{personne vivant avec le VIH} » (PVVIH) et non « sidéen » (terme stigmatisant, réservé au stade Sida) ni « contaminé » (terme péjoratif).
\end{itemize}
\end{attention}

\courssub{Bilan de l'activité}

Cette activité t'a permis de mobiliser l'ensemble des savoir-faire de la leçon dans une démarche proche de celle du médecin et du biologiste :
\begin{itemize}
\item tu as \textbf{interprété des documents} (résultats de laboratoire, micrographie, courbe physiologique) pour distinguer deux grandes catégories de perturbations du système immunitaire : l'\cle{auto-immunité} (dossier A : diabète juvénile) et l'\cle{immunodéficience acquise} (dossier B : VIH/Sida) ;
\item tu as relié chaque dysfonctionnement à un \textbf{effecteur précis} : lymphocytes T cytotoxiques (LT8) et auto-anticorps (LB) auto-réactifs dans le premier cas, destruction des LT4 (CD4\textsuperscript{+}) par le VIH dans le second ;
\item tu as exploité une \textbf{courbe} pour dégager un seuil clinique (200 cellules/mm\textsuperscript{3}) et comprendre la notion d'infection opportuniste comme conséquence de la perte de coordination immunitaire ;
\item tu as utilisé une \textbf{maquette} pour matérialiser la structure de l'anticorps (domaines variables et constants) et la formation du complexe immun, compétence expérimentale exigée par le programme (réaliser et utiliser des maquettes) ;
\item enfin, tu as développé des \cle{compétences psychosociales} : proposer une prise en charge réaliste (TARV, observance, suivi), lutter contre la stigmatisation et adopter des attitudes de solidarité envers les personnes vivant avec le VIH — un enjeu de santé publique majeur au Cameroun, où la prévalence reste significative et l'accès au traitement gratuit un droit fondamental.
\end{itemize}

Tu es désormais prêt à aborder l'évaluation sommative du module, où ce type de dossier documentaire, mêlant données chiffrées, micrographies et raisonnement scientifique, est très fréquent au Baccalauréat D.

\newpage

\coursec{Méthodes et exercices résolus}

\coursec{Lutte contre les perturbations du système immunitaire}

\courssub{Origine des cellules immunitaires et structures de reconnaissance du non-soi}

\begin{methode}[Localiser les organes lymphoïdes et les structures de reconnaissance]
Pour répondre à une question portant sur l'origine et la maturation des cellules immunitaires :
\begin{enumerate}
\item \cle{Distinguer} organes lymphoïdes \textbf{primaires} (moelle osseuse : formation de tous les leucocytes par hématopoïèse et maturation des lymphocytes B ; thymus : maturation des lymphocytes T) et organes lymphoïdes \textbf{secondaires} (rate, ganglions lymphatiques, amygdales, plaques de Peyer : lieux de rencontre avec l'antigène, d'activation et de prolifération clonale).
\item \cle{Raisonner} à partir d'expériences classiques : ablation du thymus chez la souris nouveau-née (Miller, 1961) $\Rightarrow$ disparition de l'immunité à médiation cellulaire ; ablation de la bourse de Fabricius chez l'oiseau ou irradiation de la moelle osseuse chez le mammifère $\Rightarrow$ disparition de la production d'anticorps (immunité humorale).
\item \cle{Nommer} les structures de reconnaissance : les \textbf{récepteurs T} (TCR) à la surface des lymphocytes T, les \textbf{récepteurs B} (BCR, anticorps membranaires) à la surface des lymphocytes B, et le \textbf{CMH} (complexe majeur d'histocompatibilité, HLA chez l'Homme) qui marque le « soi » et présente les peptides antigéniques.
\item \cle{Conclure} toujours en reliant la structure à sa fonction : un lymphocyte ne reconnaît qu'\textbf{un seul} antigène (spécificité clonale) grâce à la diversité générée par le réarrangement somatique des gènes des récepteurs.
\end{enumerate}
\end{methode}

\begin{popfigure}
\begin{tikzpicture}[scale=0.85, every node/.style={font=\small}]
% Bone Marrow
\node[draw, rounded corners, fill=popblueL, minimum width=3.5cm, minimum height=2.5cm] (moelle) at (0,0) {};
\node[align=center, text width=3cm] at (0,0.5) {\textbf{Moelle osseuse}\\(Organe lymphoïde primaire)\\\textit{Hématopoïèse : formation de tous les leucocytes}\\\textit{Maturation des lymphocytes B}};
\node[align=center, text width=3cm] at (0,-1.2) {Cellules souches hématopoïétiques $\rightarrow$ Pro-lymphocytes B $\rightarrow$ Lymphocytes B matures};

% Thymus
\node[draw, rounded corners, fill=poppurpleL, minimum width=3.5cm, minimum height=2.5cm] (thymus) at (8,0) {};
\node[align=center, text width=3cm] at (8,0.5) {\textbf{Thymus}\\(Organe lymphoïde primaire)\\\textit{Maturation des lymphocytes T}};
\node[align=center, text width=3cm] at (8,-1.2) {Pro-thymocytes $\rightarrow$ Sélection positive/négative $\rightarrow$ Lymphocytes T4/T8 matures};

% Arrow from bone marrow to thymus
\draw[->, thick, popdark] (moelle.east) -- node[above] {Pro-thymocytes (via sang)} (thymus.west);

% Secondary organs
\node[draw, rounded corners, fill=popgreenL, minimum width=3cm, minimum height=1.8cm] (rate) at (0,-4) {};
\node[align=center, text width=2.5cm] at (0,-4) {\textbf{Rate}\\(Filtration sang, réponse antigènes sanguins)};

\node[draw, rounded corners, fill=popgreenL, minimum width=3cm, minimum height=1.8cm] (ganglion) at (4,-4) {};
\node[align=center, text width=2.5cm] at (4,-4) {\textbf{Ganglions lymphatiques}\\(Filtration lymphe, activation lymphocytes)};

\node[draw, rounded corners, fill=popgreenL, minimum width=3cm, minimum height=1.8cm] (muqueuses) at (8,-4) {};
\node[align=center, text width=2.5cm] at (8,-4) {\textbf{Muqueuses (MALT)}\\(Amygdales, plaques de Peyer,\\réponse antigènes muqueux)};

% Arrows from primary to secondary
\draw[->, thick, popdark] (moelle.south) -- ++(0,-0.5) -| (rate.north) node[pos=0.25, left] {Lymphocytes B, T matures};
\draw[->, thick, popdark] (thymus.south) -- ++(0,-0.5) -| (ganglion.north);
\draw[->, thick, popdark] (thymus.south) -- ++(0,-0.5) -| (muqueuses.north);

% Legend
\node[align=left, anchor=north west] at (11,1) {\textbf{Légende :}\\\textcolor{popblue}{Bleu} : Formation + Maturation B\\\textcolor{poppurple}{Violet} : Maturation T\\\textcolor{popgreen}{Vert} : Activation, prolifération, effet};
\end{tikzpicture}
\legende{Organisation des organes lymphoïdes : primaires (moelle osseuse, thymus) pour la génèse et la maturation, secondaires (rate, ganglions, MALT) pour la réponse immune adaptative.}
\end{popfigure}

\begin{exoresolu}[Thymectomie et greffes de peau chez la souris (Expérience de Miller, 1961)]
On réalise des greffes de peau entre souris de lignées différentes (antigéniquement distinctes, allogreffes) et on mesure le délai de rejet de la greffe :
\begin{center}
\begin{tabularx}{\linewidth}{|l|X|X|}\hline
\rowcolor{popblueL} \textbf{Lot} & \textbf{Traitement} & \textbf{Délai de rejet de la greffe} \\ \hline
A & Souris nouveau-nées témoins (non opérées) & 11 jours \\ \hline
B & Souris thymectomisées à la naissance (jour 1) & Pas de rejet (greffe tolérée indéfiniment) \\ \hline
C & Souris thymectomisées à la naissance puis greffées d'un thymus syngénique & 12 jours \\ \hline
\end{tabularx}
\end{center}
Interpréter ces résultats et préciser le rôle du thymus.
\tcblower
\textbf{Analyse des résultats.}
\begin{itemize}
\item \textbf{Comparaison Lot A / Lot B :} La thymectomie néonatale supprime totalement le rejet de la greffe (passage de 11 jours à « pas de rejet »). Le thymus est donc \textbf{indispensable} à la mise en œuvre de la réaction de rejet.
\item \textbf{Comparaison Lot B / Lot C :} La greffe d'un thymus chez la souris thymectomisée \textbf{restaure} la capacité de rejet (12 jours, valeur proche du témoin A). Cela confirme que le défaut observé au lot B est spécifiquement dû à l'absence de thymus et non à un traumatisme chirurgical accessoire.
\end{itemize}
\textbf{Interprétation mécanistique.} Le rejet d'allogreffe est une réponse immunitaire à médiation \textbf{cellulaire}, assurée principalement par les lymphocytes T cytotoxiques (LTc ou CD8$^{+}$) qui reconnaissent les molécules CMH de classe I étrangères sur les cellules du greffon et les détruisent par apoptose (perforines/granzymes). Or, les lymphocytes T acquièrent leur maturité fonctionnelle (expression du TCR, apprentissage de la restriction CMH et de la tolérance au soi par sélection positive et négative) dans le \textbf{thymus}, organe lymphoïde primaire. Sans thymus, pas de lymphocytes T matures fonctionnels $\Rightarrow$ pas de rejet.
\textbf{Conclusion.} Le thymus est le lieu exclusif de \cle{maturation} des lymphocytes T ; la moelle osseuse est le lieu de leur \cle{formation} (hématopoïèse, comme toutes les cellules sanguines) et de la maturation des lymphocytes B. La reconnaissance du non-soi par les lymphocytes T repose sur le TCR qui reconnaît un peptide antigénique présenté par une molécule du CMH (restriction CMH).
\end{exoresolu}

\courssub{Rappels : la réponse immunitaire non spécifique (innée)}

\begin{propriete}[Principaux effecteurs et mécanismes de l'immunité innée]
L'immunité innée est la première ligne de défense, rapide (heures), non spécifique (reconnaissance de motifs conservés : PAMPs par les PRR), sans mémoire. Elle prépare et oriente la réponse adaptative.
\begin{center}
\begin{tabularx}{\linewidth}{|l|X|X|}\hline
\rowcolor{popblueL} \textbf{Composante} & \textbf{Effecteurs / Molécules} & \textbf{Mécanismes d'action principaux} \\ \hline
\textbf{Barrières physico-chimiques} & Peau (kératine, pH acide, flore), muqueuses (cilia, mucus, lysozyme, défensines) & Blocage de l'entrée des pathogènes \\ \hline
\textbf{Cellules phagocytes} & \textbf{Neutrophiles} (polynucléaires, 1ers arrivés), \textbf{Macrophages} (tissulaires, dérivés des monocytes), \textbf{Dendritiques} (pont vers adaptative) & \textbf{Phagocytose} : reconnaissance (opsonines, PRR) $\rightarrow$ ingestion $\rightarrow$ phagosome $\rightarrow$ fusion lysosome $\rightarrow$ digestion enzymatique (NO, ROS, lysozyme) \\ \hline
\textbf{Cellules tueuses naturelles (NK)} & Lymphocytes NK (CD56$^{+}$, CD16$^{+}$, sans TCR/BCR) & \textbf{Cytotoxicité naturelle} : reconnaissance de l'absence de CMH-I (cellules tumorales, infectées) $\rightarrow$ lysis par perforines/granzymes ; ADCC (anticorps-dépendante) via CD16 \\ \hline
\textbf{Inflammation aiguë} & Mediateurs : Histamine, Prostaglandines, Leucotriènes, Cytokines (IL-1, TNF-$\alpha$, IL-6), Chémokines (IL-8) & Vasodilatation, perméabilité vasculaire $\uparrow$, chimiotaxie des neutrophiles/monocytes, fièvre, protéines de phase aiguë (CRP) \\ \hline
\textbf{Système du complément} & Protéines plasmatiques (C1 à C9, Facteurs B, D, Properdine) & 3 voies d'activation (Classique, Lectine, Alternative) $\rightarrow$ \textbf{C3b} (opsonisation), \textbf{C5a} (chémotaxine, anaphylatoxine), \textbf{C5b-9 (MAC)} (lyse bactérienne) \\ \hline
\textbf{Interférons de type I (IFN-$\alpha$/$\beta$)} & Cellules infectées, plasmacytoïdes & Inhibition réplication virale dans cellules voisines, activation NK, augmentation expression CMH-I \\ \hline
\end{tabularx}
\end{center}
\textbf{Articulation avec l'immunité adaptative :} Les cellules dendritiques (phagocytes « professionnels ») capturent l'antigène en périphérie, migrent vers le ganglion (organe secondaire), mûrissent (expression co-stimulation B7, CMH-II), et présentent l'antigène aux lymphocytes T naïfs $\rightarrow$ déclenchement réponse spécifique.
\end{propriete}

\courssub{La réponse immunitaire spécifique à médiation humorale}

\begin{methode}[Interpréter une expérience d'immunisation (type Behring/Kitasato, Roux)]
Face à un tableau ou un graphe de résultats expérimentaux (survie, mort, dosage d'anticorps) :
\begin{enumerate}
\item \cle{Décrire} d'abord les faits bruts, lot par lot, avec les valeurs chiffrées et les unités (survie en jours, dose injectée, titre d'anticorps).
\item \cle{Comparer} les lots deux à deux en ne faisant varier \textbf{qu'un seul facteur} (lot témoin / lot traité) pour dégager le facteur responsable de la protection.
\item \cle{Identifier} le type d'immunité mis en jeu par le \textbf{test de transfert} : 
\begin{itemize}
\item Immunité transférable par le \textbf{sérum} (fraction acellulaire contenant les anticorps) $\Rightarrow$ immunité à médiation \textbf{humorale} (cas de la diphtérie, du tétanos : antitoxines).
\item Immunité transférable \textbf{uniquement par les cellules lymphoïdes} (lymphocytes) $\Rightarrow$ immunité à médiation \textbf{cellulaire} (cas de la tuberculose, greffes, infections virales intracellulaires).
\end{itemize}
\item \cle{Conclure} en formulant le mécanisme général : Antigène $\rightarrow$ Présentation (CMH-II) $\rightarrow$ Activation Lymphocyte T4 auxiliaire (Th2) $\rightarrow$ Activation Lymphocyte B (cognate) $\rightarrow$ Différenciation en \textbf{plasmocytes} (usines à anticorps) + \textbf{lymphocytes B mémoires} $\rightarrow$ Élimination de l'antigène par les anticorps.
\end{enumerate}
\end{methode}

\begin{exoresolu}[L'expérience de Behring et Kitasato sur la diphtérie (1890) — Fondement de la sérothérapie]
On injecte à des cobayes une dose mortelle de toxine diphtérique purifiée :
\begin{center}
\begin{tabularx}{\linewidth}{|l|X|X|}\hline
\rowcolor{popblueL} \textbf{Lot} & \textbf{Traitement préalable} & \textbf{Résultat après injection de toxine mortelle} \\ \hline
1 & Aucun (témoins sains) & Mort en 3 jours \\ \hline
2 & Immunisés par injections répétées de toxine atténuée (anatoxine) & Survie (immunisation active) \\ \hline
3 & Cobayes sains recevant du \textbf{sérum} de cobayes immunisés (lot 2) & Survie (immunisation passive) \\ \hline
4 & Cobayes sains recevant des \textbf{lymphocytes} de cobayes immunisés & Mort en 3 jours \\ \hline
\end{tabularx}
\end{center}
Interpréter ces résultats et préciser le type d'immunité mis en jeu contre la diphtérie.
\tcblower
\textbf{Description des faits.} Les cobayes témoins (lot 1) meurent rapidement ; l'immunisation active préalable (lot 2) confère une protection solide. Le transfert de \textbf{sérum} immun (lot 3) transfère la protection, alors que le transfert de \textbf{lymphocytes} (lot 4) ne la transfère pas.
\textbf{Comparaisons critiques.}
\begin{itemize}
\item \textbf{Lots 1 vs 2 :} L'exposition à l'anatoxine induit la synthèse d'effecteurs protecteurs par l'organisme.
\item \textbf{Lots 3 vs 4 :} La protection est transférable par le sérum (liquide acellulaire riche en immunoglobulines) mais \textbf{pas} par les cellules lymphoïdes. Les effecteurs protecteurs sont donc des molécules solubles : les \textbf{anticorps} (ici des \cle{antitoxines} qui neutralisent spécifiquement la toxine diphtérique).
\end{itemize}
\textbf{Conclusion.} L'immunité contre la toxine diphtérique (et le tétanos) est une immunité à médiation \textbf{humorale}. Des lymphocytes B spécifiques, activés avec l'aide de lymphocytes T4 (coopération), se différencient en plasmocytes sécrétant des anticorps neutralisants. Le lot 3 illustre la \textbf{sérothérapie} (immunisation passive : anticorps reçus tout faits, protection immédiate mais transitoire, $\sim$ 3 semaines, sans mémoire) ; le lot 2 illustre le principe de la \textbf{vaccination} (immunisation active : stimulation de la réponse propre, lente à mettre en place mais durable grâce à la mémoire immunitaire).
\end{exoresolu}

\begin{methode}[Décrire et exploiter la structure d'un anticorps (Immunoglobuline G) et d'un complexe immun]
Pour construire ou exploiter une maquette (schéma, modèle moléculaire) :
\begin{enumerate}
\item \cle{Représenter} la molécule en \textbf{Y} : 4 chaînes polypeptidiques — 2 chaînes \textbf{lourdes} ($\gamma$, $\approx$ 50 kDa) et 2 chaînes \textbf{légères} ($\kappa$ ou $\lambda$, $\approx$ 25 kDa) — reliées par des \textbf{ponts disulfure} (liaisons covalentes S-S).
\item \cle{Distinguer} les régions structurales et fonctionnelles :
\begin{itemize}
\item Régions \textbf{variables} (VL et VH) : extrémités des bras du Y = fragments \textbf{Fab} (Fragment antigen-binding). Elles portent les 2 \textbf{sites de liaison à l'antigène} (paratopes), uniques pour chaque clone (spécificité). Variabilité générée par réarrangement V(D)J.
\item Régions \textbf{constantes} (CL, CH1, CH2, CH3) : base du Y = fragment \textbf{Fc} (Fragment crystallizable). Identiques pour tous les anticorps d'une même classe (isotype : IgG, IgM, IgA, IgE, IgD). Déterminent la demi-vie, le passage placentaire (IgG), la fixation sur cellules effectrices (récepteurs Fc sur macrophages, NK, mastocytes) et l'activation du complément (voie classique via C1q).
\end{itemize}
\item \cle{Exploiter} la maquette du complexe immun : un anticorps IgG est \textbf{bivalent} (2 sites Fab) ; un antigène est souvent \textbf{multivalent} (ex: protéine de surface bactérienne, hématie). La fixation croisée forme un \textbf{réseau tridimensionnel} antigène--anticorps $\rightarrow$ \textbf{agglutination} (antigène particulaire : bactéries, hématies) ou \textbf{précipitation} (antigène soluble).
\item \cle{Relier} structure et fonction : Variabilité Fab $\rightarrow$ Spécificité ; Région Fc $\rightarrow$ Fonctions effectrices (opsonisation, ADCC, activation complément, transport placentaire).
\end{enumerate}
\end{methode}

\begin{exoresolu}[Agglutination des hématies et détermination des groupes sanguins ABO]
On mélange du sang de groupe A (hématies portant l'agglutinogène A à leur surface) avec du sérum anti-A (contenant des agglutinines anti-A, anticorps IgM bivalents — en réalité décavalents pour IgM pentamérique). On observe la formation d'amas d'hématies visibles à l'œil nu. À l'aide de la structure de l'anticorps, expliquer cette observation et prévoir le résultat du mélange sang B + sérum anti-A.
\tcblower
\textbf{Explication de l'agglutination (réaction Ag/Ac visible).}
\begin{enumerate}
\item L'agglutinogène A (antigène) est un glucide membranaire porté par des millions de copies sur chaque hématie (antigène \textbf{multivalent} particulaire).
\item Les agglutinines anti-A (anticorps, ici IgM pentamérique = 10 sites de liaison) se fixent spécifiquement sur l'agglutinogène A via leurs régions variables (Fab).
\item Grâce à leur \textbf{multivalence} (10 sites pour IgM, 2 pour IgG), un seul anticorps peut lier \textbf{deux hématies différentes} simultanément (pontage).
\item Il se forme un \textbf{réseau tridimensionnel} étendu hématies--anticorps, visible macroscopiquement sous forme d'amas : c'est l'\textbf{agglutination}, cas particulier de formation de complexe immun insoluble.
\end{enumerate}
\textbf{Prévision pour Sang B + Sérum anti-A.} Le sang B porte l'agglutinogène B, structurellement distinct de l'agglutinogène A (différence d'un sucre terminal : GalNAc pour A, Galactose pour B). Les sites de liaison des agglutinines anti-A sont spécifiques de l'épitope A (relation \textbf{clé--serrure}). Il n'y a \textbf{pas de fixation}, pas de réseau, \textbf{pas d'agglutination} : le mélange reste homogène.
\textbf{Application clinique.} Ce test de \textbf{phénotypage ABO} (et Rhésus) est obligatoire avant toute transfusion sanguine (ex: au Centre National de Transfusion Sanguine de Yaoundé). Une transfusion ABO-incompatible provoquerait une agglutination massive intravasculaire des hématies reçues, hémolyse aiguë, insuffisance rénale et choc, accident souvent mortel.
\end{exoresolu}

\begin{methode}[Rédiger le devenir d'un antigène : Neutralisation puis Élimination (Phase effectrice humorale)]
Pour expliquer comment un antigène (toxine, virus, bactérie) est neutralisé puis éliminé grâce aux anticorps :
\begin{enumerate}
\item \cle{Reconnaissance spécifique} : Les sites de liaison des régions variables (Fab) de l'anticorps épousent les déterminants antigéniques (épitopes) de l'antigène (relation clé--serrure, complémentarité de forme et charges). Formation du \textbf{complexe immun} (Ag--Ac).
\item \cle{Neutralisation (blocage fonctionnel)} : Le complexe immun masque les sites fonctionnels de l'antigène : 
\begin{itemize}
\item Toxine : empêche sa fixation sur son récepteur cellulaire (ex: toxine tétanique/diphtérique).
\item Virus : empêche son attache aux récepteurs cellulaires (neutralisation virale) ou sa fusion.
\item Bactérie : empêche son adhésion aux muqueuses (immunité muqueuse IgA).
\end{itemize}
\item \cle{Élimination (effet effecteur via Fc)} : La région constante (Fc) des anticorps fixés sert de « drapeau » reconnu par les effecteurs de l'immunité innée :
\begin{itemize}
\item \textbf{Opsonisation} : Fixation des Fc sur récepteurs Fc$\gamma$R des \textbf{macrophages} et neutrophiles $\rightarrow$ phagocytose facilitée $\rightarrow$ digestion dans phagolysosomes.
\item \textbf{Activation du complément (voie classique)} : Fixation de C1q sur Fc agrégés $\rightarrow$ cascade $\rightarrow$ C3b (opsonine puissante) et C5b-9 (MAC, lyse bactérienne, surtout Gram-négatif).
\item \textbf{ADCC (Cytotoxicité à médiation cellulaire dépendante des anticorps)} : Fixation Fc sur CD16 (Fc$\gamma$RIII) des cellules NK $\rightarrow$ lysis de la cellule cible (ex: cellule infectée par virus).
\end{itemize}
\item \cle{Conclure} : L'antigène est détruit/digéré. Les peptides issus de la digestion peuvent être présentés par les macrophages/dendritiques (CMH-II) aux lymphocytes T4 auxiliaires $\rightarrow$ amplification de la réponse (coopération).
\end{enumerate}
\textbf{Réflexe rédaction Bac :} Toujours nommer les \textbf{deux temps distincts} (1. Neutralisation par Fab, 2. Élimination via Fc) et la molécule effectrice (l'anticorps). Ne jamais écrire « l'anticorps tue la bactérie ».
\end{methode}

\begin{exoresolu}[Devenir de la toxine tétanique après sérothérapie d'urgence (Contexte camerounais)]
Un enfant de 8 ans, non vacciné, se blesse profondément au pied avec un clou rouillé dans une cour à Yaoundé. Au centre de santé de district, on lui injecte immédiatement des \textbf{immunoglobulines antitétaniques humaines} (sérum antitétanique, IgG purifiées) par voie intramusculaire. Expliquer le devenir de la toxine tétanique dans son organisme.
\tcblower
\textbf{Étape 1 — Reconnaissance spécifique.} La toxine tétanique (TeNT, $\approx$ 150 kDa), neurotoxine puissante sécrétée par \emph{Clostridium tetani} (bacille anaérobie sporulé présent dans la terre), diffuse dans l'espace interstitiel. Elle est reconnue spécifiquement par les antitoxines (IgG anti-TeNT) du sérum injecté : les paratopes des régions variables (Fab) épousent les épitopes de la toxine.
\textbf{Étape 2 — Neutralisation.} Il se forme des \textbf{complexes immuns} toxine--antitoxine. La toxine, masquée par les anticorps, ne peut plus se fixer sur son récepteur spécifique (ganglioside GT1b) au niveau des terminaisons nerveuses périphériques, ni être internalisée par endocytose et transportée rétrogradement vers les interneurones inhibiteurs de la moelle épinière. Son action bloquante sur la libération de GABA/glycine (cause des contractures tétaniques généralisées, trismus, opisthotonos) est \textbf{empêchée}.
\textbf{Étape 3 — Élimination.} Les régions constantes (Fc) des anticorps fixés sur la toxine sont reconnues par les récepteurs Fc$\gamma$R des \textbf{macrophages} tissulaires et des neutrophiles recrutés par l'inflammation. Les complexes immuns sont phagocytés (\textbf{opsonisation}), incorporés dans des phagosomes qui fusionnent avec des lysosomes ; la toxine est alors \textbf{digérée} par les enzymes hydrolytiques (protéases, nucleases) à pH acide.
\textbf{Conclusion.} Grâce aux antitoxines apportées par la sérothérapie, la toxine tétanique est neutralisée sous forme de complexes immuns solubles puis éliminée par phagocytose. Il s'agit d'une \cle{immunisation passive} (anticorps reçus tout faits, protection immédiate mais brève, demi-vie IgG $\approx$ 21-28 jours, sans mémoire immunitaire) : elle \textbf{doit être complétée impérativement} par une \textbf{vaccination} antitétanique (anatoxine tétanique, immunisation active durable avec mémoire) selon le calendrier EPI du Cameroun.
\end{exoresolu}

\courssub{La réponse immunitaire spécifique à médiation cellulaire}

\begin{methode}[Décrire les mécanismes de la réponse à médiation cellulaire]
La réponse cellulaire élimine les agents \textbf{intracellulaires} (virus, bactéries intracellulaires comme \emph{Mycobacterium tuberculosis}, \emph{Listeria}, tumeurs). Elle repose sur les lymphocytes T (T4/CD4$^{+}$ auxiliaires et T8/CD8$^{+}$ cytotoxiques).
\begin{enumerate}
\item \cle{Présentation de l'antigène} : 
\begin{itemize}
\item Voie \textbf{endogène} (CMH-I) : Protéines synthétisées dans le cytosol (virus, tumeurs) $\rightarrow$ Protéasome $\rightarrow$ Peptides $\rightarrow$ RER $\rightarrow$ Assemblage CMH-I/peptide $\rightarrow$ Surface cellulaire. Reconnue par \textbf{Lymphocytes T8 (CD8$^{+}$)}.
\item Voie \textbf{exogène} (CMH-II) : Antigènes phagocytés/endocytés $\rightarrow$ Endosome/lysosome $\rightarrow$ Peptides $\rightarrow$ Assemblage CMH-II/peptide (via chaîne invariante/DM) $\rightarrow$ Surface des CPA (Dendritiques, Macrophages, Lymphocytes B). Reconnue par \textbf{Lymphocytes T4 (CD4$^{+}$)}.
\end{itemize}
\item \cle{Activation (Double signal)} : 
\begin{itemize}
\item Signal 1 : TCR + CD4/CD8 (co-récepteur) reconnaissent Complexe CMH/Peptide.
\item Signal 2 (Co-stimulation) : B7 (CD80/86) sur CPA $\leftrightarrow$ CD28 sur Lymphocyte T. Sans signal 2 $\rightarrow$ anergie/tolérance.
\item Cytokines (IL-12, IFN-$\gamma$, IL-2) $\rightarrow$ Prolifération clonale (expansion) et Différenciation.
\end{itemize}
\item \cle{Différenciation des effecteurs T4 (sous-populations Th)} :
\begin{itemize}
\item \textbf{Th1} (induits par IL-12/IFN-$\gamma$) : Sécrètent IFN-$\gamma$, TNF-$\beta$ $\rightarrow$ \textbf{Activation macrophagique} (microbicide intracellulaire, tuberculose), activation LTc, commutation IgG opsonisantes. Immunité cellulaire « classique ».
\item \textbf{Th2} (induits par IL-4) : Sécrètent IL-4, IL-5, IL-13 $\rightarrow$ Aide lymphocytes B (IgE, IgG1/IgG4), éosinophiles. Immunité humorale / anti-parasites / allergie.
\item \textbf{Th17} (IL-6, TGF-$\beta$, IL-23) : IL-17 $\rightarrow$ recrutement neutrophiles, barrières muqueuses (champignons, bactéries extracellulaires).
\item \textbf{Tfh} (Folliculaires) : Aide spécialisée aux lymphocytes B dans les centres germinatifs (affinité, commutation, mémoire).
\item \textbf{Treg} (Régulateurs, FoxP3$^{+}$) : Tolérance, frein réponse immune (IL-10, TGF-$\beta$).
\end{itemize}
\item \cle{Effecteurs T8 (LTc)} : Reconnaissance CMH-I/peptide sur cellule cible $\rightarrow$ Lysis par \textbf{perforines} (pores) + \textbf{granzymes} (apoptose caspase-dépendante) + voie Fas/FasL. Élimination directe cellule infectée/tumorale.
\item \cle{Mémoire} : Formation de lymphocytes T4 et T8 \textbf{mémoires} (central $T_{CM}$ et effecteur $T_{EM}$) : réponse secondaire plus rapide, plus forte, plus efficace.
\end{enumerate}
\end{methode}

\begin{exoresolu}[Immunité cellulaire contre la tuberculose : Expérience de transfert adoptif]
On cherche à transférer l'immunité contre \emph{Mycobacterium tuberculosis} (bacille de Koch, intracellulaire macrophage) de souris immunisées à des souris naïves.
\begin{center}
\begin{tabularx}{\linewidth}{|l|X|X|}\hline
\rowcolor{popblueL} \textbf{Lot} & \textbf{Cellules transférées (de souris immunisées)} & \textbf{Résistance à l'infection par M. tuberculosis (charge bactérienne poumon)} \\ \hline
1 & Aucune (témoins naïfs) & Faible (infection sévère) \\ \hline
2 & \textbf{Sérum} (anticorps) & Faible (infection sévère) \\ \hline
3 & \textbf{Lymphocytes T} (purifiés, CD4$^{+}$ et CD8$^{+}$) & \textbf{Forte} (charge bactérienne faible) \\ \hline
4 & \textbf{Macrophages} activés (par IFN-$\gamma$ in vitro) & Forte (transitoire) \\ \hline
\end{tabularx}
\end{center}
Interpréter. Quel type d'immunité protège contre la tuberculose ?
\tcblower
\textbf{Analyse.}
\begin{itemize}
\item Lot 1 vs Lot 3 : Le transfert de lymphocytes T confère une protection significative (charge bactérienne faible), prouvant que les cellules immunisées portent l'information de la mémoire/protection.
\item Lot 2 : Le sérum (anticorps) \textbf{ne protège pas}. L'immunité n'est pas humorale.
\item Lot 4 : Des macrophages activés \textit{in vitro} par l'IFN-$\gamma$ (cytokine Th1) protègent transitoirement, confirmant que l'effecteur final est le macrophage activé, mais l'initiateur/spécifique est le lymphocyte T.
\end{itemize}
\textbf{Conclusion.} L'immunité protectrice contre la tuberculose est une immunité à médiation \textbf{cellulaire} (Th1). Les lymphocytes T4 Th1 spécifiques sécrètent de l'IFN-$\gamma$ qui \textbf{active les macrophages} infectés, augmentant leur pouvoir microbicide (oxyde nitrique NO, espèces réactives de l'oxygène, acidification phagolysosomale) pour détruire le bacille intracellulaire. Les lymphocytes T8 cytotoxiques lysent aussi les macrophages trop lourdement infectés. C'est la base du test cutané à la tuberculine (IDR, réaction de type IV, hypersensibilité retardée, pic à 48-72h, médiation par lymphocytes T4 Th1/macrophages).
\end{exoresolu}

\courssub{Les maladies auto-immunes}

\begin{definition}[Maladie auto-immune]
Pathologie résultant d'une rupture de la \textbf{tolérance immunologique au soi} (défaillance de la sélection négative au thymus/moelle ou de la tolérance périphérique : Treg, anergie, ignorance). Le système immunitaire reconnaît anormalement des constituants propres (auto-antigènes) comme non-soi et déclenche une réponse effectrice (auto-anticorps et/ou lymphocytes T auto-réactifs) causant des lésions tissulaires.
\end{definition}

\begin{propriete}[Quatre exemples types au programme (Mécanisme $\rightarrow$ Cible $\rightarrow$ Symptômes)]
\begin{center}
\begin{tabularx}{\linewidth}{|l|X|X|X|}\hline
\rowcolor{popblueL} \textbf{Maladie} & \textbf{Cible auto-antigénique} & \textbf{Mécanisme effecteur principal} & \textbf{Conséquences cliniques principales} \\ \hline
\textbf{Diabète juvénile (Type 1)} & Cellules $\beta$ des îlots de Langerhans (pancréas) : Insuline, GAD65, IA-2, ZnT8 & \textbf{Lymphocytes T8 cytotoxiques} (destruction cellulaire) + Auto-anticorps (marqueurs diagnostiques) & Déficit absolu en insuline $\rightarrow$ \textbf{Hyperglycémie} sévère, cétose, polydipsie, polyurie, amaigrissement. Traitement : Insulinothérapie vitale. \\ \hline
\textbf{Myasthénie (Myasthénie grave)} & Récepteur nicotinique de l'acétylcholine (AChR) plaque motrice (jonction neuromusculaire) & \textbf{Auto-anticorps anti-AChR (IgG)} : \textbf{Blocage} fixation ACh + \textbf{Destruction} récepteurs (complément, phagocytose) + modulation interne & \textbf{Fatigabilité musculaire} anormale (aggravée par l'effort, améliorée au repos) : ptosis, diplopie, dysphagie, faiblesse membres. Traitement : Anticholinestérasiques (pyridostigmine), immunosuppresseurs, thymectomie. \\ \hline
\textbf{Goitre exophtalmique (Maladie de Basedow / Graves)} & Récepteur à la TSH (TSH-R) membrane cellules thyroïdiennes & \textbf{Auto-anticorps \textbf{stimulants} (TSI, IgG)} : Agonistes mimant la TSH $\rightarrow$ activation constitutive récepteur (cAMP $\uparrow$) & \textbf{Hyperthyroïdie} : Goitre diffus, tachycardie, amaigrissement, thermophobie, \textbf{exophtalmie} (infiltration orbitaire), myxœdème pré-tibial. Traitement : Antithyroïdiens de synthèse, iode 131, chirurgie. \\ \hline
\textbf{Vitiligo} & Mélanocytes (cellules pigmentaires de la peau) & \textbf{Lymphocytes T8 cytotoxiques} auto-réactifs (destruction directe) + Auto-anticorps (secondaires) & \textbf{Dépigmentation} cutanée en plaques bien limitées, symétriques souvent (visage, mains, orifices). Pas de symptôme général. Impact psychosocial majeur. \\ \hline
\end{tabularx}
\end{center}
\textbf{Point clé Bac :} Différencier \textbf{auto-anticorps bloquants} (Myasthénie : perte de fonction) et \textbf{auto-anticorps stimulants} (Basedow : gain de fonction / hyperfonction). Ne pas confondre avec l'hypothyroïdie de Hashimoto (auto-anticorps anti-TPO/Tg + lymphocytes T $\rightarrow$ destruction thyroïde $\rightarrow$ hypothyroïdie).
\end{propriete}

\begin{exoresolu}[La myasthénie : une synapse neuromusculaire sabotée par des auto-anticorps bloquants]
Madame N., 35 ans, consulte à l'Hôpital Central de Yaoundé pour une fatigabilité musculaire anormale : ses paupières tombent en fin de journée (ptosis) et elle peine à mastiquer (dysphagie). L'électromyogramme (test de stimulation répétitive) montre une diminution pathologique de la réponse musculaire (décrément $>$ 10\%). On met en évidence dans son sérum un taux élevé d'anticorps dirigés contre les récepteurs membranaires de l'acétylcholine (anti-AChR). Décrire ce dysfonctionnement du système immunitaire et expliquer les symptômes.
\tcblower
\textbf{Diagnostic mécanistique.} Il s'agit d'une \textbf{maladie auto-immune} (myasthénie auto-immune) : le système immunitaire de Madame N. a perdu la tolérance au soi et produit des \textbf{auto-anticorps} (IgG1/IgG3) dirigés contre un constituant propre, les récepteurs de l'acétylcholine (AChR) de la membrane postsynaptique des plaques motrices.
\textbf{De l'auto-anticorps au symptôme (Physiopathologie).}
\begin{enumerate}
\item \textbf{Physiologie normale :} L'influx nerveux arrive au bouton axonal $\rightarrow$ entrée Ca$^{2+}$ $\rightarrow$ libération vésiculaire d'ACh dans la fente synaptique $\rightarrow$ fixation ACh sur ses récepteurs (AChR) $\rightarrow$ ouverture canaux Na$^{+}$/K$^{+}$ $\rightarrow$ dépolarisation (potentiel de plaque) $\rightarrow$ potentiel d'action musculaire $\rightarrow$ contraction (couplage excitation-contraction).
\item \textbf{Pathologie (Madame N.) :} Les auto-anticorps anti-AChR se fixent sur les récepteurs postsynaptiques :
\begin{itemize}
\item \textbf{Blocage compétitif :} Ils empêchent physiquement la fixation de l'ACh sur son site de liaison.
\item \textbf{Destruction accélérée :} Les complexes immuns (AChR--Auto-Ac) activent le complément (voie classique $\rightarrow$ MAC) et favorisent la phagocytose des récepteurs $\rightarrow$ \textbf{diminution drastique de la densité des AChR fonctionnels} (de $\sim$ 10 000/$\mu$m$^2$ à $<$ 1 000/$\mu$m$^2$).
\item \textbf{Modulation antigénique :} Internalisation et dégradation des récepteurs liés aux anticorps.
\end{itemize}
\item \textbf{Conséquence fonctionnelle :} La marge de sécurité de la transmission neuromusculaire s'effondre. Lors de stimulations répétées (effort), la quantité d'ACh libérée diminue (déplétion vésiculaire) ; avec trop peu de récepteurs, le potentiel de plaque n'atteint plus le seuil de déclenchement du potentiel d'action musculaire $\rightarrow$ \textbf{bloc de la transmission} $\rightarrow$ faiblesse musculaire fluctuante (fatigabilité), ptosis, diplopie, dysphagie, risque de crise myasthénique (détresse respiratoire).
\end{enumerate}
\textbf{Traitement ciblé.}
\begin{itemize}
\item \textbf{Symptomatique :} Anticholinestérasiques (pyridostigmine) : inhibent l'acétylcholinestérase $\rightarrow$ augmentent la concentration et la durée d'action de l'ACh dans la fente $\rightarrow$ compensent la perte de récepteurs.
\item \textbf{Étioopathogénique :} Immunosuppresseurs (corticoïdes, azathioprine, mycophénolate) pour réduire la production d'auto-anticorps ; Plasmaphérèse / Immunoglobulines IV (crises) ; \textbf{Thymectomie} (thymome fréquent ou hyperplasie thymique) améliore la rémission.
\end{itemize}
\end{exoresolu}

\courssub{Le VIH/Sida : immunodéficience acquise et enjeux socioculturels}

\begin{methode}[Raisonner sur le VIH/Sida : Mécanisme, Prévention, Attitudes (Savoir-être)]
Pour traiter une situation liée au VIH/Sida (cas clinique, document, question de société) :
\begin{enumerate}
\item \cle{Mécanisme physiopathologique} : Le VIH (Virus de l'Immunodéficience Humaine, \textbf{rétrovirus} à ARN, genre \emph{Lentivirus}) infecte préférentiellement les lymphocytes \textbf{T4 (CD4$^{+}$)} via fixation de sa glycoprotéine d'enveloppe \textbf{gp120} sur le récepteur \textbf{CD4} et un co-récepteur (CCR5 ou CXCR4). Son enzyme clé, la \textbf{transcriptase inverse} (RT), transcrit l'ARN viral en ADN proviral ; l'\textbf{intégrase} l'insère dans le génome hôte. La réplication virale entraîne la destruction des T4 (lyse, apoptose, syncytia, mort par activation). Chute progressive des T4 $\rightarrow$ effondrement de l'orchestration immune $\rightarrow$ Stade \textbf{Sida} = \textbf{Immunodéficience acquise} sévère avec \textbf{infections opportunistes} (Tuberculose, Pneumocystose à \emph{P. jirovecii}, Candidose œsophagienne, Toxoplasmose cérébrale, Cryptococcose, CMV) et cancers (Sarcome de Kaposi, Lymphomes).
\item \cle{Modes de transmission (UNIQUEMENT 3)} : 
\begin{itemize}
\item \textbf{Sanguin} : Sang contaminé (transfusion non dépistée, partage seringues usagers de drogues, accidents d'exposition au sang - AES - personnel soignant, scarifications/tatouages non stériles).
\item \textbf{Sexuel} : Rapports non protégés (vaginaux, anaux, oraux) avec personne séropositive non traitée (charge virale detectable). Risque : Anus $>$ Vagin $>$ Pénis $>$ Oral.
\item \textbf{Mère $\rightarrow$ Enfant (Verticale)} : Grossesse (transplacentaire rare), Accouchement (contact sang/sécrétions, principal), Allaitement (lait maternel).
\end{itemize}
\textbf{JAMAIS} par : contacts quotidiens (poignée de main, embrassade, bisous, repas, toilettes, piscine, moustiques, sueur, larmes, salive sauf saignement buccal important).
\item \cle{Prévention combinée (Stratégie 95-95-95 ONUSIDA)} :
\begin{itemize}
\item \textbf{Primaire :} Abstinence, Fidélité mutuelle testée, \textbf{Préservatif} (masculin/féminin) correct et systématique ; \textbf{PrEP} (Prophylaxie Pré-Exposition, ARV chez séronégatifs à risque) ; \textbf{PEP} (Post-Exposition, urgence $<$ 48h/72h) ; Circoncision médicale volontaire (réduit risque H$\leftarrow$F $\sim$ 60\%).
\item \textbf{Secondaire :} \textbf{Dépistage Volontaire (DV)} ou Dépistage Proposé par le Prestataire (DPP) ; Accès universel au \textbf{TAR} (Traitement Antirétroviral) immédiat (« Treat All ») $\rightarrow$ Charge Virale (CV) \textbf{indétectable} ($<$ 50 copies/mL) = \textbf{Intransmissible} (I=I).
\item \textbf{Tertiaire / PTME :} Prise en charge globale femme enceinte séropositive (TAR vie entière, accouchement sécurisé, prophylaxie ARV nouveau-né, conseil alimentation : lait maternel exclusif 6 mois si TAR efficace + CV indétectable, sinon lait de substitution sûr).
\end{itemize}
\item \cle{Savoir-être (Compétences psychosociales - Vie scolaire/professionnelle)} :
\begin{itemize}
\item \textbf{Lutter contre la stigmatisation et la discrimination :} Le VIH n'est pas une maladie honteuse. Une personne vivant avec le VIH (PVVIH) sous TAR efficace mène une vie normale (études, travail, famille, sexualité protégée).
\item \textbf{Confidentialité absolue} du statut sérologique (secret médical).
\item \textbf{Soutien, empathie, non-jugement} : Favorisent l'observance (prise quotidienne des ARV à heure fixe), le dépistage précoce, la rétention dans les soins.
\item \textbf{Éducation sexuelle complète} : Déconstruire les rumeurs (moustiques, sorcellerie, guérison par plantes seules), promouvoir le consentement, l'égalité de genre.
\end{itemize}
\end{enumerate}
\end{methode}

\begin{savaistu}[Le VIH/Sida au Cameroun : Contexte épidémiologique et riposte]
\begin{itemize}
\item \textbf{Prévalence :} $\approx$ 2,7\% (adultes 15-49 ans, EDSC 2018), $\sim$ 500 000 PVVIH. Féminisation de l'épidémie (femmes 2$\times$ plus touchées que hommes jeunes). Hétérogénéité régionale (Sud, Est, Centre plus touchés).
\item \textbf{Riposte nationale (CNLS, Ministère Santé Publique) :} Gratuité du dépistage et des ARV (depuis 2007) ; Programme PTME (Prévention Transmission Mère-Enfant) performant (taux transmission $<$ 5\% visé) ; Dépistage index (partenaires/enfants) ; Lutte contre les violences basées sur le genre (VBG) et stigmatisation ; Intégration VIH/Tuberculose (co-infection fréquente).
\item \textbf{Défis :} Observance TAR long terme (résistances), populations clés (HSH, TS, UD, prisonniers) difficiles d'accès, financement pérenne, élimination transmission mère-enfant (validation OMS).
\end{itemize}
\end{savaistu}

\begin{exoresolu}[Évolution de la numération des lymphocytes T4 chez un patient infecté par le VIH non traité]
Chez un patient séropositif non traité, on suit la numération des lymphocytes T4 (CD4$^{+}$) dans le sang périphérique :
\begin{center}
\begin{tabularx}{\linewidth}{|l|X|X|X|X|X|}\hline
\rowcolor{popblueL} \textbf{Années après infection estimée} & 0 (Séroconversion) & 1 & 3 & 6 & 9 \\ \hline
\textbf{T4 (cellules/$\mu$L de sang)} & 1000 & 900 & 650 & 350 & 120 \\ \hline
\end{tabularx}
\end{center}
Sachant que le stade Sida correspond à une numération T4 $<$ 200 cellules/$\mu$L (ou survenue d'une infection opportuniste définissante), analyser ces données, expliquer l'apparition des infections opportunistes et indiquer les attitudes à adopter envers ce patient.
\tcblower
\textbf{Analyse quantitative des données.} La numération des T4 passe de \num{1000} cellules/$\mu$L (valeur normale adulte : 500--1500) à \num{120} cellules/$\mu$L en 9 ans. La perte moyenne est d'environ $\frac{1000-120}{9} \approx 98$ cellules/$\mu$L par an. La cinétique n'est pas linéaire (perte plus rapide aux stades symptomatiques). Le seuil critique de \textbf{200 cellules/$\mu$L} est franchi entre la 6\textsuperscript{e} année (350) et la 9\textsuperscript{e} année (120) : le patient entre au stade \textbf{Sida} (immunodéficience sévère) vers la 7\textsuperscript{e}--8\textsuperscript{e} année.
\textbf{Explication physiopathologique des infections opportunistes.}
\begin{enumerate}
\item Le VIH infecte massivement les lymphocytes T4 (réservoir principal) via gp120/CD4/CCR5. La réplication virale (RT, Intégrase, Protéase) détruit les cellules (lyse, mort par activation, apoptose, élimination par LTc).
\item Les lymphocytes T4 (Th) sont les \textbf{chefs d'orchestre} de l'immunité adaptative : ils sécrètent l'IL-2 (autocrine/paracrine pour prolifération T/B), l'IFN-$\gamma$ (activation macrophages, Th1), l'IL-4/IL-5/IL-13 (aide B, Th2), l'IL-17 (neutrophiles, Th17).
\item Leur chute ($<$ 200/$\mu$L) entraîne une \textbf{défaillance globale de l'immunité cellulaire et humorale} : 
\begin{itemize}
\item Pas d'activation macrophagique $\rightarrow$ prolifération intracellulaire de \emph{M. tuberculosis}, \emph{Leishmania}, champignons.
\item Pas d'aide aux lymphocytes B $\rightarrow$ hypogammaglobulinémie relative, défaut réponse vaccinale.
\item Pas de surveillance tumorale (LTc, NK) $\rightarrow$ Sarcome de Kaposi (HHV-8), Lymphomes (EBV).
\end{itemize}
\item Des germes habituellement contrôlés (commensaux ou latents) deviennent pathogènes : \textbf{infections opportunistes} (IO). Au Cameroun, la \textbf{tuberculose} est l'IO la plus fréquente et première cause de décès au stade Sida.
\end{enumerate}
\textbf{Attitudes à adopter (Savoir-être / Compétences psychosociales).}
\begin{itemize}
\item \textbf{Urgence médicale :} Orientation immédiate vers un Centre de Prise en Charge (PEC) agréé pour bilan (Charge Virale, T4, recherche IO, bilan hépatique/rénal) et \textbf{initiation rapide du TAR} (Trithérapie / TAR : 2 INI + 1 INNTI ou IP, ex: TDF/3TC/DTG). Le TAR bloque la réplication, permet la remontée des T4 (reconstitution immune), rend la CV indétectable (I=I) et redonne une espérance de vie quasi normale.
\item \textbf{Non-stigmatisation / Inclusion :} Rappeler que le VIH ne se transmet pas par les contacts quotidiens (classe, travail, sport, repas, latrines, moustiques). Le patient a droit à l'éducation, l'emploi, les soins, le respect de sa vie privée. Le rejet social pousse au cachot, au non-dépistage, à l'abandon du traitement.
\item \textbf{Soutien psychosocial :} Écoute active, conseil observance (pilulier, alarme, groupe de pairs), dépistage partenaire(s) et enfants (index testing), planning familial, nutrition, santé mentale (dépression fréquente).
\item \textbf{Confidentialité :} Le statut sérologique ne doit être divulgué qu'avec le consentement éclairé du patient.
\end{itemize}
\textbf{Conclusion.} La destruction progressive des lymphocytes T4 par le VIH explique l'effondrement des défenses immunitaires et la survenue d'infections opportunistes définissant le stade Sida. L'accès précoce au TAR, gratuit au Cameroun, et la solidarité communautaire (lutte contre la stigmatisation) permettent aujourd'hui de vivre longtemps et en bonne santé avec le VIH.
\end{exoresolu}

\begin{attention}[Les 5 erreurs qui coûtent des points au Bac (SVT Terminale D)]
\begin{enumerate}
\item \textbf{Confondre formation et maturation des lymphocytes.} Tous les leucocytes naissent dans la \textbf{moelle osseuse} (hématopoïèse) ; la maturation des B a lieu dans la moelle (bourse de Fabricius chez l'oiseau), celle des T dans le \textbf{thymus}. Écrire « les lymphocytes T sont formés dans le thymus » est \textbf{faux} (ce sont des pro-thymocytes qui y migrent pour mûrir).
\item \textbf{Affirmer que l'anticorps « tue » ou « digère » l'antigène.} L'anticorps \textbf{neutralise} (formation complexe immun, blocage fonctionnel) et \textbf{marque} l'antigène (opsonisation) ; l'élimination physique est assurée par la \textbf{phagocytose} (macrophages, neutrophiles) ou la lyse par le \textbf{complément} (MAC). Respecter scrupuleusement les deux temps : \textbf{neutralisation \emph{puis} élimination}.
\item \textbf{Inverser immunité humorale et cellulaire dans l'interprétation expérimentale.} Le critère décisif est le \textbf{transfert adoptif} : immunité transférable par le \textbf{sérum} $\Rightarrow$ anticorps $\Rightarrow$ médiation \textbf{humorale} (diphtérie, tétanos, venins) ; transférable \textbf{uniquement par les cellules lymphoïdes} $\Rightarrow$ médiation \textbf{cellulaire} (tuberculose, leishmaniose, greffes, virus intracellulaires). Toujours comparer un lot témoin et un lot expérimental avant de conclure.
\item \textbf{Décrire les auto-anticorps comme toujours destructeurs/bloquants.} Dans le \textbf{goitre exophtalmique (Basedow)}, les auto-anticorps (TSI) sont \textbf{stimulants} (agonistes du récepteur TSH $\rightarrow$ hyperthyroïdie) ; dans la \textbf{myasthénie}, ils sont \textbf{bloquants} (antagonistes récepteur ACh $\rightarrow$ perte de fonction). Adapter le mécanisme à la maladie citée.
\item \textbf{Sur le VIH/Sida : confondre séropositivité et Sida, et citer de faux modes de transmission.} Être \textbf{séropositif} signifie être porteur du VIH (anticorps anti-VIH positifs), stade souvent asymptomatique pendant des années ; le \textbf{Sida} est le stade d'immunodéficience avancée (T4 $<$ 200/$\mu$L ou IO). Ne \textbf{jamais} écrire que le VIH se transmet par les moustiques, la sueur, les larmes, les embrassades, la nourriture, les toilettes : cela entretient la stigmatisation, contredit les données scientifiques et est \textbf{sanctionné} dans les questions de savoir-être.
\end{enumerate}
\end{attention}

\newpage

\coursec{Exercices}

\courssub{Je vérifie mes connaissances}

\begin{exercice}[QCM : les bases de l'immunité]
Pour chaque question, une seule réponse est exacte. Recopie la lettre correspondante.

\textbf{1.} Les lymphocytes T acquièrent leur maturité immunologique dans :
\begin{itemize}
\item[a)] la moelle osseuse rouge ;
\item[b)] le thymus ;
\item[c)] la rate ;
\item[d)] les ganglions lymphatiques.
\end{itemize}

\textbf{2.} Un anticorps est une molécule :
\begin{itemize}
\item[a)] lipidique sécrétée par les macrophages ;
\item[b)] protéique sécrétée par les plasmocytes ;
\item[c)] glucidique présente à la surface des virus ;
\item[d)] protéique fabriquée par les hématies.
\end{itemize}

\textbf{3.} La réaction du phagocyte qui englobe puis digère un micro-organisme s'appelle :
\begin{itemize}
\item[a)] l'agglutination ;
\item[b)] la lyse ;
\item[c)] la phagocytose ;
\item[d)] la précipitation.
\end{itemize}

\textbf{4.} Le VIH infecte préférentiellement :
\begin{itemize}
\item[a)] les lymphocytes T4 (LT4) ;
\item[b)] les globules rouges ;
\item[c)] les neurones ;
\item[d)] les plaquettes.
\end{itemize}
\end{exercice}

\begin{exercice}[Vrai ou faux, en justifiant]
Indique si chaque affirmation est vraie ou fausse, puis justifie ta réponse en une ou deux phrases.
\begin{enumerate}
\item Toutes les cellules immunitaires dérivent de cellules souches de la moelle osseuse.
\item La réponse immunitaire non spécifique met en jeu des anticorps.
\item La réponse secondaire à un antigène est plus rapide et plus intense que la réponse primaire grâce aux cellules mémoires.
\item Le diabète juvénile (diabète de type 1) est une maladie auto-immune.
\item On peut contracter le VIH en serrant la main d'une personne séropositive.
\end{enumerate}
\end{exercice}

\begin{exercice}[Définitions éclair]
Définis en deux lignes maximum chacun des termes suivants : \cle{antigène}, \cle{anticorps}, \cle{complexe immun}, \cle{mémoire immunitaire}, \cle{maladie auto-immune}, \cle{immunodéficience}.
\end{exercice}

\begin{exercice}[Lecture de données chiffrées]
Chez un patient suivi au Centre Hospitalier d'Essos (Yaoundé), la numération des lymphocytes T4 donne les résultats suivants :
\begin{center}
\begin{tabular}{|l|c|c|c|}
\hline
\rowcolor{popblueL} Date & Janvier 2022 & Janvier 2023 & Janvier 2024 \\
\hline
Taux de LT4 (cellules/$\mu$L de sang) & 850 & 420 & 150 \\
\hline
\end{tabular}
\end{center}
On rappelle qu'un taux normal de LT4 est compris entre 500 et \num{1500} cellules/$\mu$L, et que le stade Sida est diagnostiqué en dessous de 200 cellules/$\mu$L.
\begin{enumerate}
\item Calcule le pourcentage de diminution du taux de LT4 entre janvier 2022 et janvier 2024.
\item À partir de quelle date le patient entre-t-il dans le stade d'immunodéficience sévère ? Justifie.
\item Quelle conséquence clinique directe cette chute des LT4 entraîne-t-elle ?
\end{enumerate}
\end{exercice}

\courssub{Je m'entraîne}

\begin{exercice}[Sérum ou lymphocytes : l'immunité antituberculeuse]
On immunise un lot A de souris contre le bacille de Koch (agent de la tuberculose). Quelques semaines plus tard, on prélève chez ces souris, d'une part du sérum, d'autre part des lymphocytes, que l'on transfère à deux lots de souris saines. Toutes les souris sont ensuite infectées par le bacille de Koch. Les taux de survie après 30 jours sont consignés ci-dessous :
\begin{center}
\begin{tabular}{|l|c|}
\hline
\rowcolor{popblueL} Lot de souris & Survie à 30 jours (\%) \\
\hline
Souris témoins (aucun transfert) & 10 \\
\hline
Souris ayant reçu le sérum & 15 \\
\hline
Souris ayant reçu les lymphocytes & 85 \\
\hline
\end{tabular}
\end{center}
\begin{enumerate}
\item Analyse et compare les taux de survie des trois lots.
\item Déduis-en le type de réponse immunitaire spécifique efficace contre le bacille de Koch (à médiation humorale ou à médiation cellulaire). Justifie ta réponse.
\item Propose une explication au niveau cellulaire : quel type de lymphocyte est impliqué et quel est son mode d'action sur les cellules infectées ?
\end{enumerate}
\end{exercice}

\begin{exercice}[Vaccination antitétanique : réponses primaire et secondaire]
Un enfant reçoit à Ngaoundéré une primo-vaccination contre le tétanos (injection $I_1$ au jour 0), puis un rappel (injection $I_2$ au jour 30). Le taux d'anticorps antitétaniques (en unités arbitraires, UA) est mesuré régulièrement :
\begin{center}
\begin{tabular}{|l|c|c|c|c|c|c|c|}
\hline
\rowcolor{popblueL} Jours après $I_1$ & 0 & 10 & 20 & 30 & 35 & 40 & 60 \\
\hline
Taux d'anticorps (UA) & 0 & 5 & 25 & 15 & 90 & 200 & 180 \\
\hline
\end{tabular}
\end{center}
\begin{enumerate}
\item Trace sur papier millimétré la courbe du taux d'anticorps en fonction du temps (échelle : 1 cm pour 5 jours ; 1 cm pour 20 UA).
\item Compare la réponse primaire (après $I_1$) et la réponse secondaire (après $I_2$) selon trois critères : délai d'apparition, intensité maximale, durée.
\item Explique, en faisant intervenir les cellules mémoires, le principe des rappels vaccinaux.
\item Le sérum antitétanique injecté en urgence après une blessure souillée confère-t-il le même type d'immunité que le vaccin ? Précise la différence entre sérothérapie et vaccination.
\end{enumerate}
\end{exercice}

\begin{exercice}[Une faiblesse musculaire inexpliquée]
Une jeune femme de 28 ans consulte à l'Hôpital Général de Douala pour une fatigabilité musculaire : ses paupières tombent en fin de journée et elle peine à mastiquer. Le dosage sanguin révèle la présence d'auto-anticorps dirigés contre les récepteurs de l'acétylcholine de la plaque motrice (jonction nerf-muscle).
\begin{enumerate}
\item Nomme la maladie dont souffre cette patiente.
\item Définis le type de dysfonctionnement immunitaire dont relève cette affection.
\item À l'aide de tes connaissances sur le fonctionnement de la synapse neuromusculaire, explique le mécanisme par lequel ces auto-anticorps provoquent la faiblesse musculaire.
\item Cite deux autres maladies relevant du même dysfonctionnement, en précisant pour chacune la cible des auto-anticorps ou des lymphocytes autoréactifs.
\end{enumerate}
\end{exercice}

\begin{exercice}[Structure d'un anticorps et complexe immun]
Le document ci-dessous représente de façon simplifiée la structure d'un anticorps (immunoglobuline G) : il est constitué de deux chaînes lourdes et de deux chaînes légères reliées par des ponts disulfures ; chaque chaîne possède une partie constante et une partie variable.
\begin{enumerate}
\item Réalise le schéma légendé d'un anticorps en faisant apparaître : les chaînes lourdes, les chaînes légères, les parties constantes, les parties variables et les deux sites de fixation de l'antigène.
\item Explique pourquoi les parties variables confèrent à l'anticorps sa spécificité.
\item Représente, par un schéma simple, un \cle{complexe immun} formé entre des anticorps et des antigènes solubles.
\item Décris le devenir de ce complexe immun dans l'organisme (rôle des phagocytes).
\end{enumerate}
\end{exercice}

\courssub{Je me prépare au Bac}

\begin{exercice}[Type Bac — De la diphtérie au Sida : spécificité des réponses immunitaires]
\textbf{Partie A.} Au début du XX\textsuperscript{e} siècle, von Behring réalise les expériences suivantes sur des cobayes, avec la toxine diphtérique :
\begin{center}
\begin{tabular}{|l|l|c|}
\hline
\rowcolor{popblueL} Lot & Traitement préalable & Effet d'une dose mortelle de toxine \\
\hline
1 (témoin) & Aucun & Mort \\
\hline
2 & Injection de toxine atténuée 15 jours avant & Survie \\
\hline
3 & Injection de sérum d'un cobaye du lot 2 & Survie \\
\hline
4 & Injection de sérum d'un cobaye immunisé contre le tétanos & Mort \\
\hline
\end{tabular}
\end{center}
\begin{enumerate}
\item Analyse les résultats des lots 1 et 2. Que montre la survie du lot 2 ?
\item Que démontre la comparaison des lots 3 et 4 quant à la nature de l'immunité antidiphtérique ? Nomme ce type d'immunité et précise les effecteurs en cause.
\item Schématise les étapes de la réponse immunitaire ainsi mise en évidence, depuis la reconnaissance de l'antigène jusqu'à la sécrétion des effecteurs.
\end{enumerate}

\textbf{Partie B.} Au Cameroun, la tuberculose reste endémique. Des travaux montrent que le transfert de sérum de souris immunisées contre le bacille de Koch ne protège pas des souris saines, alors que le transfert de leurs lymphocytes les protège.
\begin{enumerate}
\setcounter{enumi}{3}
\item Compare cette conclusion à celle de la partie A. Que révèle cette différence sur les mécanismes de défense selon la nature de l'antigène ?
\item Décris le mécanisme d'élimination d'une cellule infectée par le bacille de Koch.
\end{enumerate}

\textbf{Partie C.} Un patient séropositif non traité développe une tuberculose pulmonaire.
\begin{enumerate}
\setcounter{enumi}{5}
\item Explique pourquoi la destruction des LT4 par le VIH favorise l'apparition d'infections opportunistes comme la tuberculose.
\end{enumerate}
\end{exercice}

\begin{exercice}[Type Bac — Maladies auto-immunes : deux cas cliniques camerounais]
\textbf{Document 1.} Un nourrisson de 18 mois, suivi au Centre Mère-Enfant de Yaoundé, présente une soif intense, une émission fréquente d'urines abondantes et une glycémie à jeun de \SI{2,4}{g/L} (normale : 0,7 à \SI{1,1}{g/L}). La biopsie et les dosages révèlent une destruction des cellules $\beta$ des îlots de Langerhans du pancréas, infiltrées par des lymphocytes, et la présence d'auto-anticorps anti-cellules $\beta$.

\textbf{Document 2.} Une patiente de 35 ans présente un gonflement du cou (goitre), des yeux exorbités et une nervosité permanente. Son sang contient des auto-anticorps qui se fixent sur les récepteurs de la TSH (hormone hypophysaire stimulant la thyroïde) et les activent de façon permanente ; son taux d'hormones thyroïdiennes est très élevé alors que sa TSH est effondrée.
\begin{enumerate}
\item Nomme les deux maladies décrites et indique le dysfonctionnement immunitaire commun qu'elles illustrent. Définis ce dysfonctionnement.
\item Pour le document 1 : explique le lien entre la destruction des cellules $\beta$ et l'hyperglycémie. Quel traitement substitutif est indispensable ?
\item Pour le document 2 : à l'aide d'un schéma fonctionnel de la régulation de la thyroïde (axe hypothalamo-hypophyso-thyroïdien avec rétrocontrôle), explique pourquoi la TSH est effondrée alors que la thyroïde est hyperactive.
\item Compare les mécanismes des deux maladies : dans quel cas les auto-anticorps détruisent-ils leur cible, et dans quel cas la stimulent-ils ?
\item Le vitiligo (dépigmentation de la peau par destruction des mélanocytes) est fréquemment observé au Cameroun. Explique en quoi il peut relever du même mécanisme et cite une conséquence socioculturelle de cette maladie dans notre contexte.
\end{enumerate}
\end{exercice}

\begin{exercice}[Type Bac — Le VIH/Sida : de la biologie aux comportements citoyens]
\textbf{Document 1.} Le VIH est un rétrovirus à ARN dont l'enveloppe porte des glycoprotéines gp120 qui se fixent sur les récepteurs CD4 des lymphocytes T4. Après pénétration, son ARN est copié en ADN par la transcriptase inverse ; cet ADN s'intègre dans le génome de la cellule hôte et dirige la production de nouveaux virions qui bourgeonnent à la surface cellulaire.

\textbf{Document 2.} Évolution simplifiée du taux de LT4 et de la charge virale (nombre de copies d'ARN viral par mL de plasma) chez une personne infectée non traitée :
\begin{center}
\begin{tabular}{|l|c|c|c|c|}
\hline
\rowcolor{popblueL} Années après l'infection & 0,1 & 1 & 5 & 9 \\
\hline
LT4 (cellules/$\mu$L) & 800 & 700 & 400 & 120 \\
\hline
Charge virale (copies/mL) & \num{100000} & \num{5000} & \num{30000} & \num{300000} \\
\hline
\end{tabular}
\end{center}
\begin{enumerate}
\item À partir du document 1, décris en quatre étapes numérotées le cycle de réplication du VIH dans un LT4. Pourquoi qualifie-t-on le VIH de rétrovirus ?
\item Analyse l'évolution croisée du taux de LT4 et de la charge virale (document 2). Quelle relation peut-on établir entre les deux courbes ?
\item Explique, en t'appuyant sur le rôle central des LT4 dans les réponses immunitaires spécifiques, pourquoi le stade Sida s'accompagne d'infections opportunistes et de cancers.
\item Cite trois modes de transmission du VIH et trois idées reçues sur la transmission qu'il faut combattre dans nos communautés.
\item Les antirétroviraux (ARV), distribués gratuitement au Cameroun, inhibent notamment la transcriptase inverse. Explique leur mode d'action et précise pourquoi ils ne guérissent pas l'infection mais transforment le Sida en maladie chronique contrôlée.
\item Rédige un court texte (8 à 10 lignes) destiné à sensibiliser les élèves de ton lycée à la lutte contre la stigmatisation des personnes vivant avec le VIH.
\end{enumerate}
\end{exercice}

\newpage

\coursec{Corrigés des exercices}

\coursec{Corrigés des exercices d'application et type Bac — Lutte contre les perturbations du système immunitaire}

\begin{corrige}[Exercice 1 — QCM : les bases de l'immunité]

\textbf{1. Réponse b) : le thymus.} Les lymphocytes T naissent dans la moelle osseuse rouge, mais c'est dans le \cle{thymus} qu'ils acquièrent leur maturité immunologique (apprentissage de la reconnaissance du soi et du non-soi, expression des récepteurs T). La rate et les ganglions sont des organes lymphoïdes \emph{périphériques}, lieux d'activation et non de maturation.

\textbf{2. Réponse b) : protéique sécrétée par les plasmocytes.} Un anticorps est une \cle{immunoglobuline}, glycoprotéine sécrétée par les \cle{plasmocytes}, issus de la différenciation des lymphocytes B activés.

\textbf{3. Réponse c) : la phagocytose.} La \cle{phagocytose} comporte l'adhérence, l'endocytose (formation du phagosome), la digestion (phagolysosome) et le rejet des résidus. L'agglutination, la précipitation et la lyse sont des modes d'action des anticorps.

\textbf{4. Réponse a) : les lymphocytes T4 (LT4).} La glycoprotéine $gp_{120}$ du VIH se fixe spécifiquement sur le récepteur \cle{CD4} des $LT_4$ (et des macrophages).
\end{corrige}

\begin{attention}[Points de vigilance]
Ne pas confondre \emph{lieu de production} (moelle osseuse pour tous les leucocytes) et \emph{lieu de maturation} (thymus pour les $LT$, moelle pour les $LB$). Ne pas écrire « les anticorps phagocytent » : ce sont les phagocytes qui phagocytent ; les anticorps marquent, neutralisent, agglutinent.
\end{attention}

\begin{corrige}[Exercice 2 — Vrai ou faux, en justifiant]

\textbf{1. VRAI.} Toutes les cellules immunitaires (lymphocytes B et T, polynucléaires, monocytes/macrophages) dérivent de \cle{cellules souches hématopoïétiques} de la moelle osseuse rouge (hématopoïèse).

\textbf{2. FAUX.} La réponse non spécifique met en jeu l'inflammation et la phagocytose (polynucléaires neutrophiles, macrophages) ; elle est identique quel que soit l'agresseur. Les anticorps sont les effecteurs de la réponse \emph{spécifique} à médiation humorale.

\textbf{3. VRAI.} Lors du premier contact, des lymphocytes B et T \cle{mémoires} sont produits. Au second contact, ils sont réactivés immédiatement : la réponse secondaire est plus rapide (délai de latence réduit), plus intense (pic d'anticorps plus élevé) et plus durable : c'est la \cle{mémoire immunitaire}, principe de la vaccination.

\textbf{4. VRAI.} Le diabète de type 1 résulte de la destruction des cellules $\beta$ des îlots de Langerhans par des lymphocytes autoréactifs et des auto-anticorps : c'est une \cle{maladie auto-immune} organospécifique.

\textbf{5. FAUX.} Le VIH ne se transmet que par le sang, les rapports sexuels non protégés et de la mère à l'enfant. Aucun contact quotidien (poignée de main, embrassade, repas partagé, piqûre de moustique) ne transmet le virus : c'est une idée reçue source de stigmatisation.
\end{corrige}

\begin{corrige}[Exercice 3 — Définitions éclair]

\begin{itemize}
\item \cle{Antigène} : substance étrangère à l'organisme (ou modifiée), reconnue comme non-soi, capable de déclencher une réponse immunitaire spécifique.
\item \cle{Anticorps} : glycoprotéine (immunoglobuline) sécrétée par les plasmocytes, capable de se fixer spécifiquement sur l'antigène qui a induit sa production.
\item \cle{Complexe immun} : ensemble formé par la fixation d'anticorps sur des antigènes ; il neutralise l'antigène et facilite son élimination par les phagocytes.
\item \cle{Mémoire immunitaire} : propriété du système immunitaire, liée aux lymphocytes mémoires, de répondre plus vite et plus fort lors d'un second contact avec un antigène déjà rencontré.
\item \cle{Maladie auto-immune} : maladie résultant d'une réaction immunitaire dirigée contre des constituants du soi (auto-anticorps et/ou lymphocytes autoréactifs), par rupture de la tolérance.
\item \cle{Immunodéficience} : défaillance quantitative ou qualitative du système immunitaire, congénitale ou acquise (ex. infection à VIH), exposant l'organisme aux infections opportunistes.
\end{itemize}
\end{corrige}

\begin{corrige}[Exercice 4 — Lecture de données chiffrées]

\textbf{1. Pourcentage de diminution entre janvier 2022 et janvier 2024.}
\[
\text{Diminution} = \frac{850 - 150}{850} \times 100 = \frac{700}{850} \times 100 \approx 82{,}4\ \%
\]
Le taux de $LT_4$ a diminué d'environ \textbf{82,4 \%} en deux ans (perte de 700 cellules/$\mu$L).

\textbf{2. Entrée dans le stade d'immunodéficience sévère.} Le seuil du stade Sida est fixé à 200 cellules/$\mu$L. En janvier 2023, le taux (420 cellules/$\mu$L) est déjà anormalement bas (norme : 500 à \num{1500} cellules/$\mu$L) mais supérieur à 200. En janvier 2024, il tombe à 150 cellules/$\mu$L $< 200$ cellules/$\mu$L : le patient entre donc dans le stade d'immunodéficience sévère (stade Sida) \textbf{entre janvier 2023 et janvier 2024}, constaté en janvier 2024.

\textbf{3. Conséquence clinique directe.} Les $LT_4$ étant les chefs d'orchestre des réponses immunitaires spécifiques (activation des $LB$ et des $LT_8$), leur effondrement rend l'organisme incapable de se défendre : apparaissent des \cle{infections opportunistes} (tuberculose, pneumocystose, candidoses, méningite à cryptocoques) et certains cancers (maladie de Kaposi).
\end{corrige}

\begin{attention}[Points de vigilance]
Le pourcentage de diminution se calcule toujours par rapport à la valeur \emph{initiale} (850), non à la valeur finale. Vérification : $850 \times (1 - 0{,}824) \approx 150$ \checkmark.
\end{attention}

\begin{corrige}[Exercice 5 — Sérum ou lymphocytes : l'immunité antituberculeuse]

\textbf{1. Analyse et comparaison.} Les souris témoins ont une survie de 10 \% : l'infection par le bacille de Koch est presque toujours mortelle. Le transfert de sérum (15 \%) n'améliore pas significativement la survie (écart de 5 points, non significatif) : les anticorps seuls ne protègent pas. En revanche, le transfert de lymphocytes porte la survie à 85 \% (soit 8,5 fois celle des témoins) : les lymphocytes des souris immunisées confèrent une protection efficace.

\textbf{2. Type de réponse efficace.} L'immunité antituberculeuse est transférable par les \emph{lymphocytes} mais pas par le \emph{sérum} (qui contient les anticorps) : il s'agit donc d'une réponse immunitaire spécifique \textbf{à médiation cellulaire}, et non humorale.

\textbf{3. Explication cellulaire.} Le bacille de Koch est une bactérie \emph{intracellulaire} : il survit dans les macrophages, à l'abri des anticorps. Les effecteurs sont les \cle{lymphocytes T8 cytotoxiques} ($LT_8$), activés avec l'aide des $LT_4$. Un $LT_8$ reconnaît spécifiquement, via son récepteur T, le complexe [peptide du bacille + $CMH_I$] présenté à la surface d'une cellule infectée ; il s'y fixe et libère des enzymes (perforines, granzymes) qui provoquent la \textbf{lyse de la cellule infectée}, détruisant ainsi le bacille qu'elle héberge.
\end{corrige}

\begin{attention}[Points de vigilance]
Raisonner toujours par comparaison avec le lot témoin. La conclusion doit être justifiée par la \emph{dissociation} des résultats : sérum inefficace / lymphocytes efficaces. Citer le caractère intracellulaire du bacille est un argument décisif au Bac.
\end{attention}

\begin{corrige}[Exercice 6 — Vaccination antitétanique : réponses primaire et secondaire]

\textbf{1. Courbe.} On reporte les points (0 ; 0), (10 ; 5), (20 ; 25), (30 ; 15), (35 ; 90), (40 ; 200), (60 ; 180), en respectant les échelles (1 cm = 5 jours ; 1 cm = 20 UA). Allure attendue :

\begin{popfigure}
\begin{center}
\begin{tikzpicture}
\begin{axis}[
  width=0.85\linewidth, height=6cm,
  xlabel={Temps (jours)}, ylabel={Taux d'anticorps (UA)},
  xmin=0, xmax=65, ymin=0, ymax=220,
  xtick={0,10,20,30,40,50,60}, ytick={0,50,100,150,200},
  grid=major, grid style={popline!40},
  axis lines=left]
\addplot[popcurve, mark=*, mark size=1.5pt] coordinates {(0,0) (10,5) (20,25) (30,15) (35,90) (40,200) (60,180)};
\draw[dashed, poporange] (axis cs:30,0) -- (axis cs:30,200) node[pos=0.95, anchor=south east, font=\small]{$I_2$};
\node[popblue, font=\small] at (axis cs:18,45) {Réponse primaire};
\node[poppurple, font=\small] at (axis cs:47,120) {Réponse secondaire};
\end{axis}
\end{tikzpicture}
\end{center}
\legende{Taux d'anticorps antitétaniques après primo-injection $I_1$ (jour 0) et rappel $I_2$ (jour 30). La réponse secondaire est plus rapide, plus intense et plus durable.}
\end{popfigure}

\textbf{2. Comparaison des deux réponses.}

\begin{center}
\begin{tabularx}{\linewidth}{|l|X|X|}
\hline
\rowcolor{popblueL} Critère & Réponse primaire (après $I_1$) & Réponse secondaire (après $I_2$) \\
\hline
Délai d'apparition & Long : anticorps détectables vers le jour 10 & Court : forte élévation dès le jour 35 (5 jours après $I_2$) \\
\hline
Intensité maximale & Faible : pic de 25 UA vers le jour 20 & Élevée : pic de 200 UA vers le jour 40 (8 fois plus) \\
\hline
Durée & Courte : déclin dès le jour 20 (15 UA au jour 30) & Longue : taux encore élevé (180 UA) au jour 60 \\
\hline
\end{tabularx}
\end{center}

\textbf{3. Principe des rappels.} Lors de la primo-injection, l'antigène tétanique (anatoxine) active des lymphocytes B spécifiques qui se différencient en plasmocytes (sécréteurs d'anticorps) \emph{et} en \cle{lymphocytes B mémoires}. Lors du rappel, ces cellules mémoires, déjà présentes et en grand nombre, sont réactivées immédiatement : elles prolifèrent et produisent massivement des plasmocytes, d'où une réponse secondaire rapide, intense et durable. Les rappels vaccinaux exploitent donc la \cle{mémoire immunitaire} pour maintenir une protection à long terme.

\textbf{4. Sérothérapie versus vaccination.} Non. Le sérum antitétanique apporte des anticorps \emph{tout faits} : c'est une \textbf{immunisation passive} — protection immédiate mais brève (quelques semaines, le temps de l'élimination des anticorps injectés), sans mémoire. Le vaccin stimule la production propre d'anticorps et de cellules mémoires : c'est une \textbf{immunisation active} — protection différée mais durable. En pratique, après une blessure souillée, on associe souvent les deux (sérum pour l'urgence, vaccin pour la durée).
\end{corrige}

\begin{attention}[Points de vigilance]
Ne pas confondre immunisation passive (transfert d'anticorps, pas de mémoire) et active (stimulation du système immunitaire, mémoire). Au Bac, préciser que le sérum contient des immunoglobulines polyclonales et que sa demi-vie est courte (environ 3 semaines pour les IgG).
\end{attention}

\begin{corrige}[Exercice 7 — Une faiblesse musculaire inexpliquée]

\textbf{1. Diagnostic.} Il s'agit de la \cle{myasthénie} (myasthenia gravis), caractérisée par une fatigabilité musculaire fluctuante (ptosis, troubles de la mastication) et des auto-anticorps anti-récepteurs de l'acétylcholine.

\textbf{2. Type de dysfonctionnement.} C'est une \cle{maladie auto-immune} : rupture de la tolérance au soi, avec production d'auto-anticorps dirigés contre un constituant propre de l'organisme (ici les récepteurs membranaires de la plaque motrice).

\textbf{3. Mécanisme.} Normalement, l'influx nerveux libère l'acétylcholine dans la fente synaptique ; sa fixation sur les récepteurs de la membrane musculaire déclenche un potentiel d'action postsynaptique, donc la contraction. Les auto-anticorps se fixent sur ces récepteurs, les bloquent et provoquent leur destruction (internalisation, action du complément). Le nombre de récepteurs fonctionnels diminue : l'acétylcholine libérée ne peut plus dépolariser efficacement la fibre musculaire, d'où une transmission neuromusculaire défaillante et une faiblesse musculaire qui s'aggrave à l'effort (épuisement du stock de médiateur).

\textbf{4. Autres exemples.}
\begin{itemize}
\item \cle{Diabète juvénile} (type 1) : destruction des cellules $\beta$ des îlots de Langerhans par des lymphocytes autoréactifs et des auto-anticorps anti-cellules $\beta$.
\item \cle{Goitre exophtalmique} (maladie de Basedow) : auto-anticorps qui se fixent sur les récepteurs de la $TSH$ et stimulent en permanence la thyroïde (hyperthyroïdie).
\item (Autre possibilité : le \cle{vitiligo}, destruction auto-immune des mélanocytes.)
\end{itemize}
\end{corrige}

\begin{corrige}[Exercice 8 — Structure d'un anticorps et complexe immun]

\textbf{1. Schéma légendé de l'anticorps.}

\begin{popfigure}
\begin{center}
\begin{tikzpicture}[scale=1]
\draw[very thick, popblue] (-0.5,0) -- (-0.5,1.8) -- (-1.6,3.2);
\draw[very thick, popblue] (0.5,0) -- (0.5,1.8) -- (1.6,3.2);
\draw[very thick, poporange] (-1.05,2.2) -- (-1.85,3.2);
\draw[very thick, poporange] (1.05,2.2) -- (1.85,3.2);
\draw[popdark, dashed] (-0.5,1.6) -- (0.5,1.6);
\draw[popdark, dashed] (-0.95,2.35) -- (-0.62,2.35);
\draw[popdark, dashed] (0.95,2.35) -- (0.62,2.35);
\fill[popgreen] (-1.72,3.35) circle (0.12);
\fill[popgreen] (1.72,3.35) circle (0.12);
\node[popblue, font=\small, anchor=east] at (-2.6,1.0) {Chaîne lourde};
\draw[->, popblue] (-2.55,1.0) -- (-0.55,1.0);
\node[poporange, font=\small, anchor=west] at (2.6,2.6) {Chaîne légère};
\draw[->, poporange] (2.55,2.6) -- (1.5,2.7);
\node[popgreen, font=\small, anchor=west] at (2.6,3.5) {Site de fixation de l'antigène};
\draw[->, popgreen] (2.55,3.5) -- (1.85,3.4);
\node[popdark, font=\small, anchor=east] at (-2.6,1.6) {Ponts disulfures};
\draw[->, popdark] (-2.55,1.6) -- (-0.6,1.6);
\node[font=\small, anchor=west] at (2.6,0.4) {Parties constantes};
\draw[->, popmuted] (2.55,0.4) -- (0.55,0.4);
\node[font=\small, anchor=east] at (-2.6,3.0) {Parties variables};
\draw[->, popmuted] (-2.55,3.0) -- (-1.35,2.9);
\end{tikzpicture}
\end{center}
\legende{Schéma d'une immunoglobuline G : deux chaînes lourdes (bleu) et deux chaînes légères (orange) unies par des ponts disulfures ; les extrémités variables forment deux sites de fixation spécifiques de l'antigène (vert) ; la molécule a une forme en Y.}
\end{popfigure}

\textbf{2. Rôle des parties variables.} La séquence en acides aminés des parties variables diffère d'un anticorps à l'autre : elle détermine la configuration spatiale (structure tertiaire) du site de fixation. Ce site est complémentaire d'un déterminant antigénique précis, comme une clé pour une serrure : c'est ce qui confère à l'anticorps sa \cle{spécificité}. La partie constante, identique pour une classe d'immunoglobulines, assure les fonctions effectrices (fixation du complément, reconnaissance par les phagocytes).

\textbf{3. Complexe immun.}

\begin{popfigure}
\begin{center}
\begin{tikzpicture}[scale=0.9]
\fill[popgreen] (0,0) circle (0.18);
\fill[popgreen] (1.6,0.6) circle (0.18);
\fill[popgreen] (0.8,-0.9) circle (0.18);
\fill[popgreen] (2.2,-0.5) circle (0.18);
\draw[very thick, popblue] (0.1,0.15) -- (0.7,0.4) -- (1.45,0.55);
\draw[very thick, popblue] (0.7,0.4) -- (0.75,-0.7);
\draw[very thick, popblue] (1.65,0.45) -- (1.9,-0.1) -- (2.1,-0.4);
\draw[very thick, popblue] (0.15,-0.1) -- (0.45,-0.5) -- (0.7,-0.8);
\draw[very thick, popblue] (1.9,-0.1) -- (1.35,-0.5) -- (0.95,-0.85);
\node[popgreen, font=\small] at (3.2,0.6) {Antigènes};
\node[popblue, font=\small] at (3.2,-0.6) {Anticorps};
\end{tikzpicture}
\end{center}
\legende{Complexe immun : chaque anticorps, divalent, relie deux antigènes ; un réseau (précipité) se forme lorsque les antigènes sont solubles et en proportions convenables.}
\end{popfigure}

\textbf{4. Devenir du complexe immun.} Le complexe immun neutralise l'antigène (il ne peut plus agir) et le rend plus facilement capturable : les macrophages et polynucléaires possèdent des récepteurs pour la partie constante des anticorps (phagocytose facilitée = \cle{opsonisation}). Le complexe est alors \cle{phagocyté} et digéré dans les phagolysosomes ; les résidus sont rejetés. Le complément peut aussi être activé et contribuer à la lyse ou à l'élimination.
\end{corrige}

\begin{attention}[Points de vigilance]
Un anticorps est divalent (deux sites de fixation) ; un $IgM$ pentamère est décavalent. La formation de réseau (précipitation) nécessite des antigènes multivalents (protéines, polysaccharides) et un rapport antigène/anticorps optimal (zone d'équivalence).
\end{attention}

\begin{corrige}[Exercice 9 — Type Bac : de la diphtérie au Sida]

\textbf{Barème indicatif (sur 20) : Partie A : 7 pts ; Partie B : 6 pts ; Partie C : 4 pts ; présentation et rigueur du vocabulaire : 3 pts.}

\textbf{Partie A.}

\textbf{1.} Le lot 1 (témoin) meurt : la dose de toxine est bien mortelle. Le lot 2, ayant reçu 15 jours plus tôt de la toxine atténuée, survit : le contact préalable avec l'antigène atténué a déclenché un état de \cle{résistance acquise} (immunisation active) ; l'organisme a fabriqué ses propres moyens de défense.

\textbf{2.} Le sérum d'un cobaye immunisé contre la diphtérie protège (lot 3), alors que le sérum d'un cobaye immunisé contre le tétanos ne protège pas (lot 4). Conclusions :
\begin{itemize}
\item l'immunité antidiphtérique est transférable par le \emph{sérum} : elle repose sur des molécules solubles, les \cle{anticorps} (antitoxines) — c'est une immunité \textbf{à médiation humorale} ;
\item elle est \textbf{spécifique} : les anticorps antitétaniques sont inefficaces contre la toxine diphtérique.
\end{itemize}
Effecteurs : les plasmocytes (issus des $LB$ activés avec l'aide des $LT_4$) sécrètent les anticorps qui neutralisent la toxine.

\textbf{3.} Étapes de la réponse à médiation humorale :
\begin{enumerate}
\item phagocytose de l'antigène par un macrophage et présentation des peptides antigéniques (complexe $CMH_{II}$) ;
\item reconnaissance par un $LT_4$ spécifique, qui s'active et sécrète des interleukines ;
\item activation du lymphocyte B spécifique (double signal : antigène + $LT_4$) ;
\item multiplication clonale du $LB$ et différenciation en \cle{plasmocytes} et en lymphocytes B mémoires ;
\item sécrétion massive d'anticorps spécifiques qui neutralisent l'antigène (complexes immuns éliminés par les phagocytes).
\end{enumerate}

\textbf{Partie B.}

\textbf{4.} Contre la diphtérie (toxine soluble, extracellulaire), le sérum suffit : immunité humorale. Contre le bacille de Koch, seuls les lymphocytes protègent : immunité \textbf{à médiation cellulaire}. Cette différence révèle que le type de réponse efficace dépend de la \emph{nature de l'antigène} : les anticorps n'atteignent pas un micro-organisme \emph{intracellulaire} ; seule la destruction de la cellule hôte infectée permet de l'éliminer.

\textbf{5.} Élimination d'une cellule infectée : le macrophage infecté présente à sa surface des peptides du bacille associés au $CMH_I$ ; un \cle{$LT_8$ cytotoxique} spécifique (activé au préalable avec l'aide des $LT_4$) reconnaît ce complexe via son récepteur T, s'accole à la cellule cible et libère des enzymes lytiques (perforines qui percent la membrane, granzymes qui déclenchent l'apoptose) : la cellule infectée est lysée et le bacille détruit.

\textbf{Partie C.}

\textbf{6.} Les $LT_4$ sont indispensables à l'activation des $LB$ (réponse humorale) et des $LT_8$ (réponse cellulaire) : ils coordonnent toute la réponse spécifique. En détruisant les $LT_4$, le VIH paralyse les deux branches de l'immunité spécifique ; or la défense antituberculeuse repose précisément sur l'immunité à médiation cellulaire (partie B). Le bacille de Koch, présent à l'état latent ou rencontré dans l'environnement, n'est plus contenu : la tuberculose se déclare. C'est une \cle{infection opportuniste} révélatrice du stade Sida.
\end{corrige}

\begin{attention}[Points de vigilance]
\begin{itemize}
\item Toujours exploiter \emph{tous} les lots (le lot 4 est le contrôle de spécificité : ne pas l'oublier).
\item Justifier chaque conclusion par la comparaison chiffrée ou qualifiée des lots (survie/mort).
\item Vocabulaire exigé : « médiation humorale », « médiation cellulaire », « spécificité », « $LT_8$ cytotoxique », « infection opportuniste ».
\end{itemize}
\end{attention}

\begin{corrige}[Exercice 10 — Type Bac : maladies auto-immunes, deux cas cliniques]

\textbf{Barème indicatif (sur 20) : Q1 : 3 pts ; Q2 : 3 pts ; Q3 : 5 pts (schéma compris) ; Q4 : 3 pts ; Q5 : 4 pts ; présentation : 2 pts.}

\textbf{1.} Document 1 : \cle{diabète juvénile} (diabète de type 1). Document 2 : \cle{goitre exophtalmique} (maladie de Basedow, hyperthyroïdie auto-immune). Dysfonctionnement commun : une \cle{maladie auto-immune}, c'est-à-dire une réaction immunitaire anormale dirigée contre des constituants du soi (auto-anticorps et/ou lymphocytes autoréactifs), par défaut de la tolérance immunitaire.

\textbf{2.} Les cellules $\beta$ des îlots de Langerhans sécrètent l'\cle{insuline}, seule hormone hypoglycémiante (elle favorise l'entrée du glucose dans les cellules et son stockage). Leur destruction entraîne un déficit total en insuline : le glucose s'accumule dans le sang (glycémie à \SI{2,4}{g/L}, très supérieure à la norme), d'où hyperglycémie, glycosurie, polyurie et polydipsie. Le traitement substitutif indispensable est l'injection quotidienne d'\textbf{insuline} (insulinothérapie à vie), associée à la surveillance glycémique et à des règles hygiéno-diététiques.

\textbf{3.} Schéma fonctionnel de l'axe hypothalamo-hypophyso-thyroïdien :

\begin{popfigure}
\begin{center}
\begin{tikzpicture}[node distance=1.1cm, every node/.style={font=\small}]
\node[draw, rounded corners, fill=popblueL] (hypo) {Hypothalamus ($TRH$)};
\node[draw, rounded corners, fill=poppurpleL, below=of hypo] (hyp) {Hypophyse ($TSH$)};
\node[draw, rounded corners, fill=popgreenL, below=of hyp] (thy) {Thyroïde ($T_3$--$T_4$)};
\draw[->, thick, popblue] (hypo) -- node[right]{+} (hyp);
\draw[->, thick, poppurple] (hyp) -- node[right]{+} (thy);
\draw[->, thick, poporange] (thy.east) .. controls (3.2,-2.2) .. node[right]{rétrocontrôle $-$} (hyp.east);
\draw[->, thick, poporange] (thy.west) .. controls (-3.2,-2.2) .. node[left]{rétrocontrôle $-$} (hypo.west);
\node[draw, rounded corners, fill=poppinkL, right=2.6cm of hyp] (ac) {Auto-anticorps anti-récepteurs $TSH$};
\draw[->, thick, poppink] (ac.south) .. controls (2.8,-3.4) .. node[right]{activation permanente +} (thy.east);
\end{tikzpicture}
\end{center}
\legende{Régulation de la thyroïde : la $TRH$ stimule l'hypophyse, la $TSH$ stimule la thyroïde ; les hormones thyroïdiennes $T_3$--$T_4$ freinent (rétrocontrôle négatif) l'hypothalamus et l'hypophyse. Les auto-anticorps activent la thyroïde en permanence, hors de tout contrôle.}
\end{popfigure}

Les auto-anticorps miment l'action de la $TSH$ et stimulent la thyroïde de façon continue : sécrétion excessive de $T_3$--$T_4$ (hyperthyroïdie, d'où nervosité, goitre, exophtalmie). L'excès d'hormones thyroïdiennes exerce un rétrocontrôle négatif puissant sur l'hypophyse : la sécrétion de $TSH$ est bloquée, d'où une \textbf{$TSH$ effondrée malgré une thyroïde hyperactive}. La stimulation par les auto-anticorps échappe à toute régulation.

\textbf{4.} Comparaison :
\begin{itemize}
\item Diabète de type 1 : les auto-anticorps et surtout les lymphocytes autoréactifs \textbf{détruisent} leur cible (cellules $\beta$) $\Rightarrow$ déficit hormonal (hypo-insulinémie).
\item Goitre exophtalmique : les auto-anticorps \textbf{stimulent} leur cible (récepteurs de la $TSH$) $\Rightarrow$ excès hormonal (hyperthyroïdie).
\end{itemize}
Une maladie auto-immune peut donc être \emph{destructrice} ou \emph{stimulatrice} selon la nature de la cible.

\textbf{5.} Dans le \cle{vitiligo}, des lymphocytes autoréactifs et des auto-anticorps détruisent les \cle{mélanocytes} (cellules productrices de mélanine) : la peau se dépigmente par plaques. C'est donc le même mécanisme auto-immun destructeur que le diabète de type 1. Conséquence socioculturelle au Cameroun : les plaques très visibles exposent à la stigmatisation, aux idées reçues (maladie surnaturelle, contagiosité supposée), au rejet social et au recours tardif aux soins ; l'éducation à la santé doit rappeler que le vitiligo n'est \emph{ni contagieux ni une malédiction}.
\end{corrige}

\begin{attention}[Points de vigilance]
Ne pas écrire que « la $TSH$ stimule trop la thyroïde » dans le document 2 : la $TSH$ est \emph{effondrée} ; ce sont les auto-anticorps qui stimulent. Le schéma de régulation doit obligatoirement faire figurer les flèches de rétrocontrôle négatif.
\end{attention}

\begin{corrige}[Exercice 11 — Type Bac : le VIH/Sida, de la biologie aux comportements citoyens]

\textbf{Barème indicatif (sur 20) : Q1 : 4 pts ; Q2 : 3 pts ; Q3 : 4 pts ; Q4 : 3 pts ; Q5 : 3 pts ; Q6 : 3 pts.}

\textbf{1. Cycle de réplication du VIH.}
\begin{enumerate}
\item \textbf{Fixation et pénétration} : la $gp_{120}$ du virus se fixe sur le récepteur $CD_4$ du $LT_4$ ; l'enveloppe virale fusionne avec la membrane et le contenu viral pénètre dans la cellule.
\item \textbf{Transcription inverse} : la transcriptase inverse copie l'$ARN$ viral en $ADN$.
\item \textbf{Intégration} : l'$ADN$ viral s'insère dans le génome de la cellule hôte (provirus).
\item \textbf{Expression et libération} : l'$ADN$ viral dirige la synthèse des protéines et de l'$ARN$ viraux ; de nouveaux virions s'assemblent et bourgeonnent à la surface du $LT_4$, qu'ils finissent par détruire.
\end{enumerate}
Le VIH est un \cle{rétrovirus} car il réalise une transcription \emph{inverse} : de l'$ARN$ vers l'$ADN$, contrairement au sens habituel de l'information génétique ($ADN \rightarrow ARN$).

\textbf{2. Évolution croisée.} Juste après l'infection (0,1 an), la charge virale est très élevée (\num{100000} copies/mL) puis chute (\num{5000} copies/mL à 1 an) grâce aux défenses immunitaires encore efficaces ; le taux de $LT_4$ reste proche de la normale (800 puis 700 cellules/$\mu$L). Ensuite, la charge virale remonte progressivement (\num{30000} puis \num{300000} copies/mL) tandis que les $LT_4$ s'effondrent (400 puis 120 cellules/$\mu$L). Relation : les deux grandeurs évoluent en sens inverse — plus le virus se multiplie, plus il détruit de $LT_4$ ; et moins il reste de $LT_4$, moins l'immunité contrôle la réplication virale (cercle vicieux).

\textbf{3.} Les $LT_4$ coordonnent les réponses spécifiques : ils activent les $LB$ (production d'anticorps) et les $LT_8$ cytotoxiques (destruction des cellules infectées ou cancéreuses), et stimulent les macrophages. Leur destruction par le VIH abolit ces deux bras de l'immunité : les agents infectieux habituellement inoffensifs ou contenus (champignons, bacille de Koch, virus) provoquent des \cle{infections opportunistes}, et la surveillance immunitaire des cellules cancéreuses étant levée, des cancers apparaissent (maladie de Kaposi, lymphomes). C'est le stade Sida ($LT_4 < 200$ cellules/$\mu$L).

\textbf{4. Modes de transmission} : (1) rapports sexuels non protégés avec une personne infectée ; (2) sang contaminé (transfusion non contrôlée, partage d'aiguilles ou d'objets tranchants non stérilisés) ; (3) transmission de la mère à l'enfant (grossesse, accouchement, allaitement).

\textbf{Idées reçues à combattre} : le VIH ne se transmet \emph{pas} par la poignée de main, les embrassades ou les caresses ; ni par le partage des repas, de la vaisselle ou des toilettes ; ni par les piqûres de moustiques.

\textbf{5. Mode d'action des ARV.} En inhibant la transcriptase inverse, les antirétroviraux bloquent l'étape 2 du cycle : l'$ARN$ viral ne peut plus être copié en $ADN$, donc aucun nouveau provirus ne se forme et la production de virions s'effondre (charge virale indétectable). Mais les ARV n'éliminent pas l'$ADN$ viral déjà intégré dans le génome de certains $LT_4$ (réservoirs) : l'infection persiste. Le traitement, pris à vie, maintient néanmoins des $LT_4$ suffisants : le Sida devient une \textbf{maladie chronique contrôlée}, et une personne sous traitement efficace ne transmet plus le virus (charge virale indétectable = intransmissible).

\textbf{6. Texte de sensibilisation (exemple).} « Camarades, une personne vivant avec le VIH est d'abord une personne. Le virus ne se transmet ni en lui serrant la main, ni en partageant son repas, ni en l'embrassant : la rejeter, c'est ajouter une souffrance injuste à sa maladie. Grâce aux antirétroviraux gratuits, elle peut vivre longtemps, étudier, travailler et fonder une famille. Le vrai danger, c'est l'ignorance : faisons le dépistage, protégeons-nous, et tendons la main à celles et ceux qui vivent avec le VIH. Ensemble, contre le virus — jamais contre les personnes. »
\end{corrige}

\begin{attention}[Points de vigilance]
\begin{itemize}
\item Distinguer \emph{infection à VIH} (portage du virus) et \emph{Sida} (stade de la maladie, $LT_4 < 200$ cellules/$\mu$L ou infections opportunistes).
\item Ne jamais écrire que les ARV « tuent le virus » : ils bloquent sa réplication sans éradiquer les réservoirs.
\item Pour l'analyse du document 2, citer les valeurs chiffrées avec leurs unités et expliciter la relation inverse entre les deux courbes.
\end{itemize}
\end{attention}

\newpage

\coursec{Fiche bilan}

\courssub{Carte mentale de la leçon}

\begin{popfigure}
\begin{tikzpicture}[mindmap, grow cyclic, every node/.style={concept, text=white, font=\small\bfseries},
root concept/.append style={concept color=popdark, font=\large\bfseries},
level 1/.append style={level distance=3.6cm, sibling angle=60, font=\small},
level 2/.append style={level distance=2.2cm, sibling angle=45, font=\scriptsize}]
\node[root concept]{Perturbations du système immunitaire}
  child[concept color=popblue]{ node {Origine des cellules immunitaires}
    child { node {Moelle osseuse : LB, LT, phagocytes} }
    child { node {Thymus : maturation des LT} }
  }
  child[concept color=popgreen]{ node {Reconnaissance du non-soi}
    child { node {CMH (soi)} }
    child { node {Antigènes, anticorps, récepteurs T} }
  }
  child[concept color=poporange]{ node {Réponse non spécifique}
    child { node {Inflammation, phagocytose} }
  }
  child[concept color=poppurple]{ node {Réponse humorale}
    child { node {LB $\to$ plasmocytes $\to$ anticorps} }
    child { node {Complexes immuns} }
  }
  child[concept color=poppink]{ node {Réponse cellulaire}
    child { node {LT8 cytotoxiques} }
    child { node {LT4 auxiliaires} }
  }
  child[concept color=popgold]{ node {Maladies auto-immunes}
    child { node {Vitiligo, diabète juvénile, myasthénie, goitre exophtalmique} }
  }
  child[concept color=popmuted]{ node {VIH/Sida}
    child { node {Destruction des LT4} }
    child { node {Prévention, dépistage, solidarité} }
  };
\end{tikzpicture}
\legende{Carte mentale de la leçon : de l'origine des cellules immunitaires aux dysfonctionnements (auto-immunité et Sida).}
\end{popfigure}

\begin{bilan}
\textbf{Définitions clés}
\begin{itemize}
  \item \cle{Antigène} : substance étrangère (non-soi) reconnue par le système immunitaire et capable de déclencher une réponse immunitaire.
  \item \cle{Anticorps} (immunoglobuline) : glycoprotéine en Y produite par les plasmocytes, formée de 4 chaînes polypeptidiques (2 lourdes, 2 légères) ; ses sites de fixation sont spécifiques d'un antigène.
  \item \cle{CMH} (complexe majeur d'histocompatibilité) : ensemble de glycoprotéines membranaires, marqueurs du \cle{soi}, propres à chaque individu (sauf vrais jumeaux).
  \item \cle{Maladie auto-immune} : maladie due à une réaction du système immunitaire contre les constituants du propre organisme (rupture de la tolérance du soi).
  \item \cle{Sida} : syndrome d'immunodéficience acquise, stade ultime de l'infection par le \cle{VIH}, rétrovirus qui détruit les lymphocytes T4.
\end{itemize}

\textbf{Origine et maturation des cellules immunitaires}
\begin{itemize}
  \item Toutes les cellules immunitaires dérivent des \cle{cellules souches hématopoïétiques} de la moelle osseuse rouge.
  \item \cle{LB} : formation et maturation dans la moelle osseuse. \cle{LT} : naissance dans la moelle, maturation dans le \cle{thymus}.
  \item Organes lymphoïdes secondaires (rate, ganglions, amygdales) : lieux de rencontre avec les antigènes et d'activation.
\end{itemize}

\textbf{Les deux réponses immunitaires spécifiques (à connaître par cœur)}
\begin{center}
\begin{tabularx}{\linewidth}{|l|X|X|}
\hline
\rowcolor{popblueL} & \textbf{Réponse humorale} & \textbf{Réponse cellulaire} \\
\hline
Effecteur & LB activés $\to$ plasmocytes & LT8 cytotoxiques (activés par les LT4) \\
\hline
Cible & Antigènes libres (toxines, microbes extracellulaires) & Cellules infectées ou cancéreuses (antigène + CMH) \\
\hline
Arme & Anticorps : agglutination, neutralisation, formation de complexes immuns & Contact direct : libération de perforines $\to$ lyse de la cellule cible \\
\hline
Élimination & Phagocytose des complexes immuns par les macrophages & Phagocytose des débris cellulaires \\
\hline
\end{tabularx}
\end{center}
\begin{itemize}
  \item \cle{Mémoire immunitaire} : lors du premier contact, des lymphocytes mémoires sont conservés ; au second contact, la réponse est plus \textbf{rapide}, plus \textbf{intense} et plus \textbf{durable} (principe de la vaccination).
  \item Réponse non spécifique (rappel) : barrières naturelles, réaction inflammatoire, phagocytose — immédiate mais sans mémoire ni spécificité.
\end{itemize}

\textbf{Maladies auto-immunes : mécanismes}
\begin{itemize}
  \item \cle{Vitiligo} : destruction des mélanocytes $\to$ dépigmentation cutanée.
  \item \cle{Diabète juvénile} (type 1) : destruction des cellules $\beta$ des îlots de Langerhans $\to$ absence d'insuline, hyperglycémie chronique.
  \item \cle{Myasthénie} : auto-anticorps contre les récepteurs de l'acétylcholine $\to$ faiblesse musculaire.
  \item \cle{Goitre exophtalmique} (Basedow) : auto-anticorps qui stimulent les récepteurs de la TSH $\to$ hyperthyroïdie.
\end{itemize}

\textbf{VIH/Sida}
\begin{itemize}
  \item Le VIH infecte et détruit les \cle{LT4} : effondrement des défenses $\to$ infections opportunistes (tuberculose, méningites, candidoses) et cancers.
  \item Transmission : sang, rapports sexuels non protégés, mère $\to$ enfant. \textbf{Pas de transmission} par les contacts quotidiens (poignée de main, repas, moustiques).
  \item Lutte : abstinence, fidélité, préservatif, dépistage volontaire, antirétroviraux (ARV), rejet de la stigmatisation et solidarité avec les personnes vivant avec le VIH.
\end{itemize}

\textbf{Méthodes express}
\begin{itemize}
  \item Interpréter une expérience d'immunisation (tétanos/diphtérie) : comparer les lots témoin/expérimental, identifier sérum (anticorps) ou cellules (lymphocytes) transférées, conclure sur le type d'immunité (humorale ou cellulaire).
  \item Sérologie : anticorps anti-VIH présents = sujet \cle{séropositif} (infecté), pas forcément malade du Sida.
  \item Complexe immun = antigène + anticorps ; sa formation facilite la phagocytose.
\end{itemize}
\end{bilan}

\begin{aretenir}[Checklist avant le Bac]
\begin{itemize}
  \item $\square$ Je sais situer la moelle osseuse (formation des LB et LT) et le thymus (maturation des LT).
  \item $\square$ Je définis antigène, anticorps, CMH et j'explique la distinction soi / non-soi.
  \item $\square$ Je décris les étapes de la phagocytose et de la réaction inflammatoire.
  \item $\square$ Je décris la structure d'un anticorps (4 chaînes, 2 sites de fixation spécifiques) et d'un complexe immun.
  \item $\square$ Je distingue réponse humorale (anticorps, antigènes libres) et réponse cellulaire (LT8, cellules infectées).
  \item $\square$ J'explique la mémoire immunitaire et le principe de la vaccination.
  \item $\square$ Je sais interpréter les expériences historiques sur l'immunité (sérum antitétanique, transfert de lymphocytes).
  \item $\square$ Je cite les 4 maladies auto-immunes du programme avec leur mécanisme.
  \item $\square$ J'explique pourquoi la destruction des LT4 par le VIH conduit aux infections opportunistes.
  \item $\square$ Je connais les modes de transmission du VIH et les moyens de prévention, et je combats la stigmatisation.
  \item $\square$ Je sais lire une courbe de concentration d'anticorps après primo-infection puis second contact.
  \item $\square$ Je rédige mes réponses avec le vocabulaire exact : plasmocyte, lymphocyte mémoire, séropositif, tolérance du soi.
\end{itemize}
\end{aretenir}

\findecours

\end{document}
